% Espaces vectoriels et applications linéaires en dimension 3 — Mathématiques Tle E — Cours signé OnBuch+
% Programme officiel MINESEC. Compiler avec : tectonic M16-espaces-vectoriels-dimension-3.tex
\newcommand{\DOCMATIERE}{Mathématiques}
\newcommand{\DOCNIVEAU}{Tle E}
\newcommand{\DOCMODULE}{Géométrie dans l'espace}
\newcommand{\DOCLECON}{16}
\newcommand{\DOCTITRE}{Espaces vectoriels et applications linéaires en dimension 3}
% OnBuch+ — préambule « cours signés » ENRICHI (compilé avec tectonic / XeLaTeX)
% Système visuel « Pop » d'OnBuch+ : fond crème (jamais blanc), bordures encre
% épaisses, ombres « dures » décalées, titres Archivo Black en capitales.
% Version MATHÉMATIQUES : graphes (pgfplots/tikz), boîtes pédagogiques
% (définition, propriété, méthode, attention, le savais-tu, exercice résolu,
% exercice, TP, bilan), page de garde et signature OnBuch+ ancrée en bas.
%
% Macros attendues dans main.tex AVANT \input{preamble} :
%   \DOCMATIERE  \DOCNIVEAU  \DOCMODULE  \DOCLECON (numéro)  \DOCTITRE
\documentclass[11pt]{article}

\usepackage{fontspec}
\usepackage[french]{babel}
\usepackage[a4paper,margin=2cm,top=3.4cm,bottom=2.8cm,headheight=1.4cm,footskip=1.3cm]{geometry}
\usepackage{amsmath,amssymb,amsfonts}
\usepackage{mathtools} % \xmapsto, \coloneqq... souvent utilisés par les modèles
\usepackage{cancel} % simplifications barrées
% \abs{z} (module, valeur absolue) et \norm{u} (norme), utilisés par les modèles
\providecommand{\abs}{}\let\abs\relax\DeclarePairedDelimiter{\abs}{\lvert}{\rvert}
\providecommand{\norm}{}\let\norm\relax\DeclarePairedDelimiter{\norm}{\lVert}{\rVert}
\usepackage[table]{xcolor} % [table] : fournit aussi \rowcolors (alternance de lignes)
\usepackage{pagecolor}
\usepackage{tikz}
\usetikzlibrary{arrows.meta,positioning,calc,shapes.geometric,shapes.misc,shapes.arrows,decorations.pathmorphing,decorations.markings,patterns,shadows,mindmap,trees,fit,backgrounds,matrix,intersections}
\usepackage{pgfplots}
\usepackage{pgf-pie} % diagrammes circulaires \pie (statistiques)
\pgfplotsset{compat=1.18}
\usepgfplotslibrary{fillbetween,groupplots} % aires entre courbes (calcul intégral)
\usepackage{enumitem}
\usepackage{fancyhdr}
% [most] charge toutes les bibliothèques tcolorbox (listings, minted, raster,
% documentation...) et gonfle inutilement la mémoire TeX sur un document long
% (~300 boîtes) : on ne charge que ce qui est réellement utilisé.
\usepackage[skins,breakable]{tcolorbox}
\usepackage{graphicx}
\usepackage{colortbl}
\usepackage{booktabs}
\usepackage{tabularx}
\usepackage{diagbox} % case d'en-tête diagonale des tableaux à double entrée
% Colonnes extensibles centrée / gauche / droite (C, L, R), souvent utilisées par les modèles
\newcolumntype{C}{>{\centering\arraybackslash}X}
\newcolumntype{L}{>{\raggedright\arraybackslash}X}
\newcolumntype{R}{>{\raggedleft\arraybackslash}X}
\usepackage{array}
\usepackage{multicol}
\usepackage{siunitx}
% parse-numbers=false : accepte aussi \SI{2,5 \times 10^{-3}}{...}
\sisetup{locale=FR,output-decimal-marker={,},group-minimum-digits=4,parse-numbers=false}
\DeclareSIUnit\ohm{\ensuremath{\symup{\Omega}}} % U+2126 absent de la police
\usepackage{qrcode}
% \degree : commande fréquemment utilisée par les modèles pour un angle ou une
% température en dehors de \SI{}{} (non fournie par les paquets chargés).
\providecommand{\degree}{^{\circ}}
% \qed (fin de démonstration), utilisé par les modèles sans amsthm
\providecommand{\qed}{\hfill$\square$}
% \barre : double barre des valeurs interdites dans un tableau de variations (tkz-tab)
\providecommand{\barre}{\ensuremath{\|}}
% environnement proof (amsthm non chargé) utilisé par les modèles
\ifdefined\proof\else\newenvironment{proof}[1][Démonstration]{\par\noindent\textit{#1.}\ }{\qed\par}\fi
\usepackage{lastpage}
\usepackage{hyperref}
\hypersetup{hidelinks}

% ============================================================
% Palette « Pop » OnBuch+ — mêmes hex que la classe P (plus_ui.dart)
% ============================================================
\definecolor{poporange}{HTML}{F59321}
\definecolor{popdark}{HTML}{E07A0C}
\definecolor{popcream}{HTML}{FFF6E8}
\definecolor{popink}{HTML}{1C1714}
\definecolor{poppurple}{HTML}{7A5AE0}
\definecolor{popgreen}{HTML}{2E9E6A}
\definecolor{poppink}{HTML}{E0457B}
\definecolor{popblue}{HTML}{2F7DE1}
\definecolor{popgold}{HTML}{C98A12}
\definecolor{popmuted}{HTML}{8A7F76}
\definecolor{popline}{HTML}{EADFD2}
\definecolor{popgray}{HTML}{8A7F76}
\colorlet{popgrey}{popgray}
% Alias tolérants (le modèle nomme parfois les couleurs différemment)
\colorlet{popwhite}{white}
\colorlet{popblack}{popink}
\colorlet{popred}{poppink}
\colorlet{popyellow}{popgold}
% Teintes claires pour fonds de tableaux / zones
\colorlet{popblueL}{popblue!10!white}
\colorlet{poporangeL}{poporange!14!white}
\colorlet{popgreenL}{popgreen!12!white}
\colorlet{poppurpleL}{poppurple!12!white}
\colorlet{poppinkL}{poppink!12!white}
\colorlet{popgoldL}{popgold!14!white}

\pagecolor{popcream}

% ============================================================
% Typographie
% ============================================================
\defaultfontfeatures{Ligatures=TeX}
\setmainfont{PlusJakartaSans-Medium.ttf}[Path=../../../fonts/,BoldFont=PlusJakartaSans-Bold.ttf]
% Maths en sans-serif (Fira Math) avec chiffres et lettres droites de Plus
% Jakarta, pour que les formules se fondent dans le texte.
\usepackage[math-style=ISO,bold-style=ISO]{unicode-math}
\setmathfont{FiraMath-Regular.otf}[Path=../../../fonts/]
\setmathfont{PlusJakartaSans-Medium.ttf}[Path=../../../fonts/,range={up/num,up/{Latin,latin},\mathup}]
\usepackage{icomma} % virgule décimale sans espace en mode math
% Caractères mathématiques Unicode absents des polices de texte (Plus Jakarta,
% Archivo Black) : rendus par leur équivalent mathématique, en texte comme en
% formule — sinon ils s'affichent en carré vide (ex. « Arithmétique dans ℤ »).
\usepackage{newunicodechar}
\newunicodechar{ℝ}{\ensuremath{\mathbb{R}}}
\newunicodechar{ℤ}{\ensuremath{\mathbb{Z}}}
\newunicodechar{ℂ}{\ensuremath{\mathbb{C}}}
\newunicodechar{ℕ}{\ensuremath{\mathbb{N}}}
\newunicodechar{ℚ}{\ensuremath{\mathbb{Q}}}
\newunicodechar{𝔻}{\ensuremath{\mathbb{D}}}
\newunicodechar{α}{\ensuremath{\alpha}}
\newunicodechar{β}{\ensuremath{\beta}}
\newunicodechar{γ}{\ensuremath{\gamma}}
\newunicodechar{δ}{\ensuremath{\delta}}
\newunicodechar{ε}{\ensuremath{\varepsilon}}
\newunicodechar{θ}{\ensuremath{\theta}}
\newunicodechar{λ}{\ensuremath{\lambda}}
\newunicodechar{μ}{\ensuremath{\mu}}
\newunicodechar{σ}{\ensuremath{\sigma}}
\newunicodechar{τ}{\ensuremath{\tau}}
\newunicodechar{φ}{\ensuremath{\varphi}}
\newunicodechar{ψ}{\ensuremath{\psi}}
\newunicodechar{ω}{\ensuremath{\omega}}
\newunicodechar{Σ}{\ensuremath{\Sigma}}
% majuscules produites par \MakeUppercase dans les titres (α-aminés -> Α)
\newunicodechar{Α}{\ensuremath{\alpha}}
\newunicodechar{Β}{\ensuremath{\beta}}
\newunicodechar{Γ}{\ensuremath{\Gamma}}
\newunicodechar{Δ}{\ensuremath{\Delta}}
\newunicodechar{Ε}{\ensuremath{\varepsilon}}
\newunicodechar{Θ}{\ensuremath{\Theta}}
\newunicodechar{Λ}{\ensuremath{\Lambda}}
\newunicodechar{Μ}{\ensuremath{\mu}}
\newunicodechar{Τ}{\ensuremath{\tau}}
\newunicodechar{Φ}{\ensuremath{\Phi}}
\newunicodechar{Ψ}{\ensuremath{\Psi}}
\newunicodechar{Ω}{\ensuremath{\Omega}}
\newunicodechar{↦}{\ensuremath{\mapsto}}
\newunicodechar{∈}{\ensuremath{\in}}
\newunicodechar{∉}{\ensuremath{\notin}}
\newunicodechar{∧}{\ensuremath{\wedge}}
\newunicodechar{⇔}{\ensuremath{\Leftrightarrow}}
\newunicodechar{⟺}{\ensuremath{\Longleftrightarrow}}
\newunicodechar{⇒}{\ensuremath{\Rightarrow}}
\newunicodechar{⟹}{\ensuremath{\Longrightarrow}}
\newunicodechar{∘}{\ensuremath{\circ}}
\newunicodechar{∪}{\ensuremath{\cup}}
\newunicodechar{∩}{\ensuremath{\cap}}
\newunicodechar{≡}{\ensuremath{\equiv}}
\newunicodechar{⊂}{\ensuremath{\subset}}
\newunicodechar{⊥}{\ensuremath{\perp}}
\newunicodechar{∃}{\ensuremath{\exists}}
\newunicodechar{∀}{\ensuremath{\forall}}
\newunicodechar{∛}{\ensuremath{\sqrt[3]{\,}}}
\newunicodechar{∜}{\ensuremath{\sqrt[4]{\,}}}
\newunicodechar{⃗}{\ensuremath{{}^{\to}}}
\newunicodechar{ⁿ}{\ifmmode^{n}\else\textsuperscript{\ensuremath{n}}\fi}
\newunicodechar{ˣ}{\ifmmode^{x}\else\textsuperscript{\ensuremath{x}}\fi}
\newunicodechar{ᵇ}{\ifmmode^{b}\else\textsuperscript{\ensuremath{b}}\fi}
\newunicodechar{ʳ}{\ifmmode^{r}\else\textsuperscript{\ensuremath{r}}\fi}
\newunicodechar{ᵘ}{\ifmmode^{u}\else\textsuperscript{\ensuremath{u}}\fi}
\newunicodechar{ʸ}{\ifmmode^{y}\else\textsuperscript{\ensuremath{y}}\fi}
\newunicodechar{ᵃ}{\ifmmode^{a}\else\textsuperscript{\ensuremath{a}}\fi}
\newunicodechar{ᵉ}{\ifmmode^{e}\else\textsuperscript{\ensuremath{e}}\fi}
\newunicodechar{ᵏ}{\ifmmode^{k}\else\textsuperscript{\ensuremath{k}}\fi}
\newunicodechar{ᵖ}{\ifmmode^{p}\else\textsuperscript{\ensuremath{p}}\fi}
\newunicodechar{ᴺ}{\ifmmode^{N}\else\textsuperscript{\ensuremath{N}}\fi}
\newunicodechar{ᶜ}{\ifmmode^{c}\else\textsuperscript{\ensuremath{c}}\fi}
\newunicodechar{ʰ}{\ifmmode^{h}\else\textsuperscript{\ensuremath{h}}\fi}
\newunicodechar{ᵠ}{\ifmmode^{q}\else\textsuperscript{\ensuremath{q}}\fi}
\newunicodechar{ₙ}{\ifmmode_{n}\else\textsubscript{\ensuremath{n}}\fi}
\newunicodechar{ₐ}{\ifmmode_{a}\else\textsubscript{\ensuremath{a}}\fi}
\newunicodechar{ᵢ}{\ifmmode_{i}\else\textsubscript{\ensuremath{i}}\fi}
\newunicodechar{ᵣ}{\ifmmode_{r}\else\textsubscript{\ensuremath{r}}\fi}
\newunicodechar{ₖ}{\ifmmode_{k}\else\textsubscript{\ensuremath{k}}\fi}
\newunicodechar{ᵦ}{\ifmmode_{\beta}\else\textsubscript{\ensuremath{\beta}}\fi}
\newunicodechar{ₓ}{\ifmmode_{x}\else\textsubscript{\ensuremath{x}}\fi}
\newfontfamily\popdisplay{ArchivoBlack-Regular.ttf}[Path=../../../fonts/]
\newfontfamily\poplogo{SpaceGrotesk-Bold.ttf}[Path=../../../fonts/]
\newfontfamily\popsemi{PlusJakartaSans-SemiBold.ttf}[Path=../../../fonts/]
\newfontfamily\popxbold{PlusJakartaSans-ExtraBold.ttf}[Path=../../../fonts/]
\newfontfamily\popxboldls{PlusJakartaSans-ExtraBold.ttf}[Path=../../../fonts/,LetterSpace=3.0]

\setlist[enumerate]{leftmargin=1.7em,itemsep=5pt,topsep=4pt}
\setlist[itemize]{leftmargin=1.7em,itemsep=4pt,topsep=4pt,label={\color{poporange}\textbullet}}
\setlength{\parindent}{0pt}
\setlength{\parskip}{5pt}
\renewcommand{\arraystretch}{1.25}

% pgfplots : style Pop par défaut
\pgfplotsset{
  every axis/.append style={axis line style={popink,line width=1pt},tick style={popink},
    grid style={popline,line width=0.5pt},label style={font=\small},tick label style={font=\footnotesize}},
  popcurve/.style={poporange,line width=1.8pt,no markers,smooth},
}
% Même style disponible dans un \draw TikZ simple (hors pgfplots)
\tikzset{popcurve/.style={poporange,line width=1.8pt}}

% ============================================================
% Pill / squiggle
% ============================================================
\newcommand{\poppill}[3]{%
  \tikz[baseline=-0.55ex]\node[draw=popink,line width=1.2pt,rounded corners=2.6pt,%
    fill=#2,inner xsep=6.5pt,inner ysep=3pt]{%
    \popxboldls\fontsize{7}{7}\selectfont\color{#3}\MakeUppercase{#1}};%
}
\newcommand{\popsquiggle}[1]{%
  \tikz[baseline=-0.4ex]{
    \draw[#1,line width=2.6pt,line cap=round]
      plot [smooth] coordinates {(0,0.09) (0.34,0.01) (0.68,0.21) (1.02,0.01) (1.36,0.10)};
  }%
}
% Mot-clé surligné (marqueur orange sous le texte)
\newcommand{\cle}[1]{{\popxbold\color{popdark}#1}}

% ============================================================
% En-tête / pied
% ============================================================
\pagestyle{fancy}
\fancyhf{}
\renewcommand{\headrulewidth}{0pt}
\renewcommand{\footrulewidth}{0pt}
\lhead{\raisebox{-0.15ex}{\poplogo\fontsize{13}{13}\selectfont\color{popink}OnBuch\color{poporange}+}}
\rhead{\poppill{\DOCMATIERE\ \textbullet\ \DOCNIVEAU\ \textbullet\ Leçon \DOCLECON}{white}{popink}}
\lfoot{\popxbold\fontsize{7.6}{7.6}\selectfont\color{popmuted}\MakeUppercase{Cours signé OnBuch+}\ \textbullet\ {\popsemi\DOCTITRE}}
\rfoot{\tikz[baseline=-0.6ex]\node[rounded corners=8.5pt,fill=popink,minimum height=17pt,minimum width=17pt,inner xsep=5pt,inner sep=0]{\popxbold\fontsize{8}{8}\selectfont\color{popcream}\thepage\,/\,\pageref*{LastPage}};}

% ============================================================
% Titres
% ============================================================
\newcommand{\chaptitre}[2]{%
  \begin{tcolorbox}[enhanced,colback=poporange,colframe=popink,boxrule=2.6pt,arc=11pt,
      drop shadow=popink,left=15pt,right=15pt,top=12pt,bottom=11pt,before skip=3pt,after skip=18pt]
    \poppill{#2}{popink}{popcream}\par\vspace{9pt}
    {\raggedright\hyphenpenalty=10000\exhyphenpenalty=10000\popdisplay\fontsize{22}{24}\selectfont\color{popcream}\MakeUppercase{#1}\par}
  \end{tcolorbox}
}
% Section numérotée : pastille orange avec le numéro + titre Archivo
\newcounter{popsec}
\newcommand{\coursec}[1]{%
  \stepcounter{popsec}\setcounter{popsub}{0}%
  \par\vspace{16pt}\noindent
  \begin{minipage}{\linewidth}
  \tikz[baseline=-0.7ex]\node[draw=popink,line width=1.6pt,fill=poporange,rounded corners=4pt,minimum size=22pt,inner sep=0]{\popdisplay\fontsize{12}{12}\selectfont\color{popcream}\thepopsec};%
  \hspace{8pt}{\popdisplay\fontsize{15}{17}\selectfont\color{popink}\MakeUppercase{#1}}\par
  \vspace{4pt}\noindent{\color{poporange}\rule{36pt}{3.6pt}}
  \end{minipage}\par\nopagebreak\vspace{8pt}
}
\newcounter{popsub}
\newcommand{\courssub}[1]{%
  \stepcounter{popsub}%
  \par\vspace{9pt}\noindent{\popxbold\fontsize{11.5}{13}\selectfont\color{popink}{\color{poporange}\thepopsec.\thepopsub}\hspace{6pt}#1}\par\nopagebreak\vspace{4pt}
}

% ============================================================
% Boîtes pédagogiques — même recette : fond blanc, bordure épaisse colorée,
% ombre dure encre, badge plein. Titre optionnel [#1].
% ============================================================
\tcbset{popbase/.style={breakable,enhanced,colback=white,boxrule=2.3pt,arc=9pt,
  shadow={3pt}{-3pt}{0pt}{popink},left=14pt,right=14pt,top=11pt,bottom=11pt,before skip=11pt,after skip=15pt,
  fontupper=\popsemi\fontsize{10.6}{14.5}\selectfont\color{popink},
  fontlower=\popsemi\fontsize{10.6}{14.5}\selectfont\color{popink}}}
% Coche dessinée (le glyphe ✓ n'existe pas dans les polices du document)
\newcommand{\popcheck}[1]{\tikz[baseline=-0.5ex]\draw[#1,line width=1.6pt,line cap=round,line join=round](-0.13,0.0)--(-0.04,-0.09)--(0.13,0.1);}
\newcommand{\popboxhead}[3]{\poppill{#1}{#2}{white}\ifx\relax#3\relax\else\hspace{6pt}{\popxbold\fontsize{10.6}{12}\selectfont\color{popink}#3}\fi\par\vspace{7pt}}

\newtcolorbox{aretenir}[1][]{popbase,colframe=popgreen,before upper={\popboxhead{À retenir}{popgreen}{#1}}}
\newtcolorbox{exemplebox}[1][]{popbase,colframe=poporange,before upper={\popboxhead{Exemple}{poporange}{#1}}}
\newtcolorbox{definition}[1][]{popbase,colframe=popblue,before upper={\popboxhead{Définition}{popblue}{#1}}}
\newtcolorbox{propriete}[1][]{popbase,colframe=poppurple,before upper={\popboxhead{Propriété}{poppurple}{#1}}}
\newtcolorbox{methode}[1][]{popbase,colframe=popgold,before upper={\popboxhead{Méthode}{popgold}{#1}}}
\newtcolorbox{attention}[1][]{popbase,colframe=poppink,before upper={\popboxhead{Attention}{poppink}{#1}}}
\newtcolorbox{experience}[1][]{popbase,colframe=popblue,colback=popblueL,before upper={\popboxhead{Expérience}{popblue}{#1}}}
% « Le savais-tu ? » : carte violette pleine (contexte, culture, Cameroun)
\newtcolorbox{savaistu}[1][]{popbase,colback=poppurple,colframe=popink,
  fontupper=\popsemi\fontsize{10.4}{14.2}\selectfont\color{white},
  before upper={\poppill{Le savais-tu\,?}{white}{poppurple}\ifx\relax#1\relax\else\hspace{6pt}{\popxbold\color{white}#1}\fi\par\vspace{7pt}}}
% Exercice résolu : énoncé en haut, \tcblower puis la solution sur fond crème
\newcounter{popexo}\newcounter{popexr}
\newtcolorbox{exoresolu}[1][]{popbase,colframe=poporange,colbacklower=popcream,
  segmentation style={popink,line width=1.2pt,dashed},
  before upper={\stepcounter{popexr}\popboxhead{Exercice résolu \thepopexr}{poporange}{#1}},
  before lower={\poppill{Solution}{popink}{popcream}\par\vspace{6pt}}}
% Exercice (non résolu en place) — la correction est dans la partie Corrigés
\newtcolorbox{exercice}[1][]{popbase,colframe=popink,
  before upper={\stepcounter{popexo}\popboxhead{Exercice \thepopexo}{popink}{#1}}}
% Corrigé
\newtcolorbox{corrige}[1][]{popbase,colframe=popgreen,colback=popgreenL,
  before upper={\popboxhead{Corrigé}{popgreen}{#1}}}
% Objectifs / prérequis / bilan
\newtcolorbox{objectifs}{popbase,colframe=popink,colback=poporangeL,
  before upper={\popboxhead{Objectifs de la leçon}{popink}{}}}
\newtcolorbox{prerequis}{popbase,colframe=popink,colback=popblueL,
  before upper={\popboxhead{Prérequis}{popblue}{}}}
\newtcolorbox{bilan}{popbase,colframe=popink,colback=white,boxrule=2.8pt,
  before upper={\popboxhead{Fiche bilan}{poporange}{L'essentiel à connaître pour l'examen}}}
% Figure encadrée avec légende
\newtcolorbox{popfigure}[1][]{enhanced,breakable=false,colback=white,colframe=popline,boxrule=1.4pt,arc=8pt,
  left=10pt,right=10pt,top=10pt,bottom=8pt,before skip=10pt,after skip=12pt,halign=center,
  lower separated=false,halign lower=center,
  fontlower=\popsemi\fontsize{9}{11.5}\selectfont\color{popmuted},
  #1}
\newcommand{\legende}[1]{\par\vspace{6pt}{\popsemi\fontsize{9}{11.5}\selectfont\color{popmuted}#1}}

% ============================================================
% Signature OnBuch+
% ============================================================
\newcommand{\poplogoinline}[1]{{\poplogo\fontsize{#1}{#1}\selectfont\color{popink}OnBuch\color{poporange}+}}
\newcommand{\signatureonbuch}{%
  \begin{tcolorbox}[enhanced,colback=popink,colframe=popink,boxrule=2.2pt,arc=10pt,
      drop shadow=poporange,left=14pt,right=14pt,top=10pt,bottom=10pt,before skip=0pt,after skip=0pt]
    \tikz[baseline=-0.7ex]\node[circle,fill=popgreen,draw=popcream,line width=1.3pt,minimum size=22pt,inner sep=0]{\popcheck{white}};%
    \hspace{9pt}\begin{minipage}[c]{0.62\linewidth}
      {\popxbold\fontsize{12}{13}\selectfont\color{popcream}Signé\ \textbullet\ {\poplogo OnBuch\color{poporange}+}}\par\vspace{2pt}
      {\raggedright\popsemi\fontsize{8}{10.5}\selectfont\color{popline}\MakeUppercase{Équipe pédagogique OnBuch+}\\\MakeUppercase{Conforme au programme officiel MINESEC}\par}
    \end{minipage}\hfill
    {\poplogo\fontsize{22}{22}\selectfont\color{popcream}OnBuch\color{poporange}+}
  \end{tcolorbox}%
}

% ============================================================
% Page de garde
% #1 titre · #2 matière · #3 niveau · #4 module · #5 n° leçon · #6 durée ·
% #7 description / objectifs (texte ou itemize)
% ============================================================
\newcommand{\pagegarde}[7]{%
  \thispagestyle{empty}
  \newgeometry{a4paper,margin=1.7cm,top=1.6cm,bottom=1.4cm}%
  \begin{tikzpicture}[remember picture,overlay]
    \fill[poporange!18] ([xshift=-3.2cm,yshift=-3.6cm]current page.north east) circle (4.6cm);
    \fill[poppurple!14] ([xshift=2.4cm,yshift=7.6cm]current page.south west) circle (3.8cm);
  \end{tikzpicture}%
  {\poplogo\fontsize{30}{30}\selectfont\color{popink}OnBuch\color{poporange}+}\hfill
  \raisebox{0.6ex}{\poppill{Cours signé \textbullet\ Premium}{popink}{popcream}}\par
  \vspace{1pt}\popsquiggle{poppurple}\par
  \vspace{8pt}
  \begin{tcolorbox}[enhanced,colback=poporange,colframe=popink,boxrule=2.8pt,arc=13pt,
      drop shadow=popink,left=18pt,right=18pt,top=12pt,bottom=13pt,after skip=10pt]
    \poppill{#2 \textbullet\ #3}{popink}{popcream}\hspace{5pt}\poppill{#4}{white}{popink}\par\vspace{8pt}
    {\popdisplay\fontsize{14}{14}\selectfont\color{popink}LEÇON #5}\par\vspace{4pt}
    {\raggedright\hyphenpenalty=10000\exhyphenpenalty=10000\popdisplay\fontsize{25}{27}\selectfont\color{popcream}\MakeUppercase{#1}\par}
    \vspace{8pt}
    {\popxbold\fontsize{9.5}{9.5}\selectfont\color{popink}\MakeUppercase{Durée conseillée : #6}}
  \end{tcolorbox}
  \begin{tcolorbox}[enhanced,colback=white,colframe=popink,boxrule=2.2pt,arc=10pt,
      drop shadow=popink,left=16pt,right=16pt,top=10pt,bottom=10pt,after skip=8pt]
    \poppill{Dans cette leçon}{poppurple}{white}\par\vspace{5pt}
    \popsemi\fontsize{9.3}{12.6}\selectfont\color{popink}#7
  \end{tcolorbox}
  \noindent\begin{minipage}[t]{0.487\linewidth}
    \begin{tcolorbox}[enhanced,colback=popgreen,colframe=popink,boxrule=2.2pt,arc=10pt,
        drop shadow=popink,left=11pt,right=11pt,top=10pt,bottom=10pt,height=3.1cm,valign=center]
      \tcbox[colback=white,colframe=popink,boxrule=1.3pt,arc=5pt,boxsep=0pt,nobeforeafter,
        left=3pt,right=3pt,top=3pt,bottom=3pt,valign=center]{\qrcode[height=1.9cm]{https://play.google.com/store/apps/details?id=com.luvvix.onbuch&hl=fr}}%
      \hspace{8pt}\begin{minipage}[c]{0.5\linewidth}
        {\popdisplay\fontsize{11}{12.5}\selectfont\color{white}\MakeUppercase{Télécharge OnBuch}}\par\vspace{4pt}
        {\raggedright\popsemi\fontsize{8.4}{11}\selectfont\color{white}Gratuit \textbullet\ Google Play\\Tuteur IA, annales, concours, résultats\par}
      \end{minipage}
    \end{tcolorbox}
  \end{minipage}\hfill
  \begin{minipage}[t]{0.487\linewidth}
    \begin{tcolorbox}[enhanced,colback=poppurple,colframe=popink,boxrule=2.2pt,arc=10pt,
        drop shadow=popink,left=11pt,right=11pt,top=10pt,bottom=10pt,height=3.1cm,valign=center]
      \tikz[baseline=-0.7ex]\node[circle,fill=white,draw=popink,line width=1.3pt,minimum size=1.6cm,inner sep=0]{\popdisplay\fontsize{24}{24}\selectfont\color{poppurple}+};%
      \hspace{8pt}\begin{minipage}[c]{0.6\linewidth}
        {\popdisplay\fontsize{11}{12.5}\selectfont\color{white}\MakeUppercase{Passe à OnBuch+}}\par\vspace{4pt}
        {\raggedright\popsemi\fontsize{8.4}{11}\selectfont\color{white}Tous les cours signés, Tuteur IA illimité, fiches PDF — onglet OnBuch+ de l'app\par}
      \end{minipage}
    \end{tcolorbox}
  \end{minipage}
  \vspace{8pt}
  \popcontenu
  \vfill
  \signatureonbuch
  \restoregeometry
  \newpage
}

% Rangée « Ce cours contient » (page de garde)
\newcommand{\poptile}[3]{% #1 couleur #2 gros texte #3 légende
  \begin{tcolorbox}[enhanced,colback=#1,colframe=popink,boxrule=2pt,arc=9pt,shadow={2.5pt}{-2.5pt}{0pt}{popink},
      left=6pt,right=6pt,top=9pt,bottom=9pt,halign=center,valign=center,height=1.75cm,width=\linewidth,nobeforeafter]
    {\popdisplay\fontsize{12.5}{13}\selectfont\color{white}\MakeUppercase{#2}}\par\vspace{3pt}
    {\centering\popsemi\fontsize{7.8}{9.5}\selectfont\color{white}#3\par}
  \end{tcolorbox}}
\newcommand{\popcontenu}{%
  \par\noindent{\popxboldls\fontsize{7.5}{7.5}\selectfont\color{popmuted}CE COURS SIGNÉ CONTIENT}\par\vspace{7pt}
  \noindent\begin{minipage}[t]{0.235\linewidth}\poptile{popblue}{Cours}{complet, illustré, pas à pas}\end{minipage}\hfill
  \begin{minipage}[t]{0.235\linewidth}\poptile{poporange}{Méthodes}{et exercices résolus}\end{minipage}\hfill
  \begin{minipage}[t]{0.235\linewidth}\poptile{poppink}{Exercices}{type examen + corrigés}\end{minipage}\hfill
  \begin{minipage}[t]{0.235\linewidth}\poptile{popgreen}{Fiche bilan}{l'essentiel à retenir}\end{minipage}\par}

% Sommaire
\newtcolorbox{sommaire}{enhanced,colback=white,colframe=popink,boxrule=2.2pt,arc=10pt,
  drop shadow=popink,left=18pt,right=18pt,top=13pt,bottom=13pt,before skip=6pt,after skip=20pt,
  before upper={\poppill{Sommaire}{poppurple}{white}\par\vspace{9pt}\popsemi\fontsize{10.6}{14.5}\selectfont\color{popink}}}

% Signature de fin de cours — poussée tout en bas de la dernière page
\newcommand{\findecours}{%
  \par\vfill\nopagebreak
  \begin{center}\popsquiggle{poporange}\end{center}\vspace{4pt}
  \signatureonbuch
}
\AtBeginDocument{\def\llbracket{\mathopen{[\mkern-3.5mu[}}\def\rrbracket{\mathclose{]\mkern-3.5mu]}}} % intervalles d'entiers (glyphe ⟦ absent de la police)
% Glyphes absents de Fira Math : redéfinis à partir de glyphes disponibles
\AtBeginDocument{%
  \def\setminus{\mathbin{\backslash}}\let\smallsetminus\setminus
  \def\vdots{\mathord{\vbox{\baselineskip3.2pt\lineskiplimit0pt\kern4pt\hbox{.}\hbox{.}\hbox{.}}}}%
  \def\ddots{\mathinner{\mkern1mu\raise6.5pt\hbox{.}\mkern2mu\raise3.6pt\hbox{.}\mkern2mu\raise0.7pt\hbox{.}\mkern1mu}}%
  \def\triangle{\mathord{\scalebox{0.9}{$\symup{\Delta}$}}}%
  \def\nabla{\mathord{\raisebox{\depth}{\scalebox{1}[-1]{$\symup{\Delta}$}}}}%
  \def\checkmark{\popcheck{popdark}}%
  \def\star{\mathbin{\ast}}\def\bigstar{\mathord{\ast}}%
  \def\diamond{\mathbin{\scalebox{0.6}{$\Diamond$}}}%
}

%% FIGFIX-BEGIN v3 — correctifs globaux des figures (docs/onbuch-generation/tools/figfix)
\usepackage{adjustbox}\usepackage{etoolbox}
% toute figure TikZ/pgfplots plus large que la ligne est réduite à la largeur disponible
\BeforeBeginEnvironment{tikzpicture}{\begin{adjustbox}{max width=\linewidth}}
\AfterEndEnvironment{tikzpicture}{\end{adjustbox}}
% légendes pgfplots lisibles par-dessus les courbes
\pgfplotsset{every axis/.append style={legend style={fill=white,fill opacity=0.92,text opacity=1,draw=black!15,inner sep=3pt}}}
% étiquettes d'arêtes (syntaxe edge["..."]) avec fond blanc
\tikzset{every edge quotes/.append style={fill=white,inner sep=1.2pt}}
% bulles de cartes mentales : texte à la ligne dans la bulle, bulle un peu plus grande
\tikzset{every concept/.append style={text width=2.0cm,align=center,minimum size=2.3cm,inner sep=2pt}}
%% FIGFIX-END
\begin{document}

\pagegarde{Espaces vectoriels et applications linéaires en dimension 3}{Mathématiques}{Tle E}{Géométrie dans l'espace}{16}{16 h}{Cette leçon introduit la structure d'espace vectoriel réel et la notion de dimension à travers l'espace ℝ³, puis étudie les applications linéaires, leurs matrices, leur noyau et leur image. Elle culmine avec le théorème du rang, outil central du programme de Terminale C et E.
\vspace{4pt}
\begin{multicols}{2}\small
\begin{itemize}
\item À la fin de cette leçon, je sais reconnaître un espace vectoriel et un sous-espace vectoriel de ℝ³
\item À la fin de cette leçon, je sais montrer qu'une famille de trois vecteurs est libre, génératrice ou forme une base de ℝ³
\item À la fin de cette leçon, je sais déterminer les coordonnées d'un vecteur dans une base donnée en résolvant un système
\item À la fin de cette leçon, je sais vérifier qu'une application est linéaire
\item À la fin de cette leçon, je sais écrire la matrice d'une application linéaire dans une base donnée
\item À la fin de cette leçon, je sais déterminer le noyau et l'image d'une application linéaire
\end{itemize}\end{multicols}}

\begin{sommaire}
\begin{multicols}{2}
\begin{enumerate}
\item Espace vectoriel ℝ³ et sous-espaces vectoriels
\item Familles libres, familles génératrices, bases
\item Coordonnées d'un vecteur dans une base
\item Applications linéaires
\item Matrice d'une application linéaire
\item Noyau, image et théorème du rang
\item Activité d'intégration
\item Méthodes et exercices résolus
\item Exercices
\item Corrigés des exercices
\item Fiche bilan
\end{enumerate}
\end{multicols}
\end{sommaire}

\begin{objectifs}
\begin{itemize}
\item À la fin de cette leçon, je sais reconnaître un espace vectoriel et un sous-espace vectoriel de ℝ³
\item À la fin de cette leçon, je sais montrer qu'une famille de trois vecteurs est libre, génératrice ou forme une base de ℝ³
\item À la fin de cette leçon, je sais déterminer les coordonnées d'un vecteur dans une base donnée en résolvant un système
\item À la fin de cette leçon, je sais vérifier qu'une application est linéaire
\item À la fin de cette leçon, je sais écrire la matrice d'une application linéaire dans une base donnée
\item À la fin de cette leçon, je sais déterminer le noyau et l'image d'une application linéaire
\item À la fin de cette leçon, je sais utiliser le théorème du rang pour trouver la dimension d'un noyau ou d'une image
\item À la fin de cette leçon, je sais caractériser l'injectivité et la surjectivité d'un endomorphisme de ℝ³ à l'aide du rang
\end{itemize}
\end{objectifs}

\begin{prerequis}
\begin{itemize}
\item Vecteurs de l'espace, opérations (somme, produit par un réel) — Première C/E
\item Repérage dans l'espace et coordonnées d'un point, d'un vecteur
\item Résolution de systèmes linéaires 3×3 (méthode du pivot de Gauss)
\item Notion d'application, image et antécédent
\item Calcul matriciel élémentaire (produit d'une matrice par un vecteur colonne)
\end{itemize}
\end{prerequis}

\newpage

\coursec{Espace vectoriel $\mathbb{R}^3$ et sous-espaces vectoriels}

À Douala, un technicien de Camrail doit décrire la position d'un wagon dans un hangar en utilisant trois directions de référence : la longueur le long des rails, la largeur du hangar et la hauteur jusqu'au toit. À Yaoundé, un infographiste 3D applique des rotations et des projections à un modèle de case traditionnelle pour l'afficher sur un écran plat. Dans les deux cas, on manipule des déplacements de l'espace, des directions, des transformations. Comment décrire rigoureusement ces déplacements et transformations ? Comment savoir si trois directions suffisent pour repérer n'importe quel point ? Comment prévoir ce qu'une transformation conserve ou écrase ?

Les \cle{espaces vectoriels} et les \cle{applications linéaires} donnent les outils mathématiques pour répondre à ces questions. Cette première section pose les fondations : nous allons donner une structure algébrique précise à l'espace qui nous entoure, puis identifier les « sous-espaces » qui vivent à l'intérieur de lui — les droites et les plans passant par l'origine.

\courssub{Rappels : vecteurs de l'espace et calcul dans $\mathbb{R}^3$}

Depuis la classe de Première, tu sais qu'un \cle{vecteur} de l'espace est caractérisé par une direction, un sens et une norme. Si l'espace est muni d'un repère $(O\,;\vec{i},\vec{j},\vec{k})$, tout vecteur $\vec{u}$ s'écrit de manière unique :
\[
\vec{u} = x\,\vec{i} + y\,\vec{j} + z\,\vec{k},
\]
où $(x,y,z)$ sont les \cle{coordonnées} de $\vec{u}$ dans la base $(\vec{i},\vec{j},\vec{k})$. On note alors :
\[
\vec{u} = \begin{pmatrix} x \\ y \\ z \end{pmatrix} \qquad \text{ou simplement} \qquad \vec{u} = (x,y,z).
\]
L'ensemble de tous les triplets de réels est noté $\mathbb{R}^3$ :
\[
\mathbb{R}^3 = \left\{ (x,y,z) \mid x \in \mathbb{R},\ y \in \mathbb{R},\ z \in \mathbb{R} \right\}.
\]

Deux opérations naturelles agissent sur ces vecteurs :

\begin{itemize}
\item \textbf{L'addition} : si $\vec{u} = (x,y,z)$ et $\vec{v} = (x',y',z')$, alors
\[
\vec{u} + \vec{v} = (x+x',\ y+y',\ z+z').
\]
Géométriquement, c'est la règle du parallélogramme (ou de Chasles) : enchaîner deux déplacements.
\item \textbf{La multiplication par un scalaire} : si $\lambda \in \mathbb{R}$, alors
\[
\lambda \cdot \vec{u} = (\lambda x,\ \lambda y,\ \lambda z).
\]
Géométriquement, c'est dilater, contracter ou renverser le déplacement $\vec{u}$.
\end{itemize}

\begin{exemplebox}[Calculs dans $\mathbb{R}^3$]
Soient $\vec{u} = (2,-1,3)$ et $\vec{v} = (1,4,-2)$. Alors :
\begin{align*}
\vec{u} + \vec{v} &= (2+1,\ -1+4,\ 3-2) = (3,3,1),\\
3\vec{u} &= (6,-3,9),\\
2\vec{u} - \vec{v} &= (4-1,\ -2-4,\ 6+2) = (3,-6,8).
\end{align*}
\end{exemplebox}

Ces deux opérations ont des propriétés remarquables : l'addition est commutative, associative, possède un élément neutre (le vecteur nul $\vec{0} = (0,0,0)$) et chaque vecteur a un opposé ; la multiplication par un scalaire se distribue sur l'addition, etc. Plutôt que de vérifier ces propriétés au cas par cas, les mathématiciens ont dégagé une définition générale : celle d'\emph{espace vectoriel}.

\courssub{Définition d'un espace vectoriel sur $\mathbb{R}$}

\begin{definition}[Espace vectoriel sur $\mathbb{R}$]
Un \cle{espace vectoriel} sur $\mathbb{R}$ est un ensemble $E$ non vide muni de deux lois :
\begin{itemize}
\item une \textbf{addition interne} $+$ : à deux éléments $\vec{u}, \vec{v}$ de $E$, on associe $\vec{u}+\vec{v} \in E$ ;
\item une \textbf{multiplication externe} $\cdot$ : à un réel $\lambda$ et un élément $\vec{u}$ de $E$, on associe $\lambda \cdot \vec{u} \in E$ ;
\end{itemize}
vérifiant les huit axiomes suivants, pour tous $\vec{u}, \vec{v}, \vec{w} \in E$ et tous $\lambda, \mu \in \mathbb{R}$ :

\textbf{Famille 1 — l'addition se comporte bien :}
\begin{enumerate}
\item $\vec{u} + \vec{v} = \vec{v} + \vec{u}$ \hfill (commutativité)
\item $(\vec{u} + \vec{v}) + \vec{w} = \vec{u} + (\vec{v} + \vec{w})$ \hfill (associativité)
\item il existe $\vec{0} \in E$ tel que $\vec{u} + \vec{0} = \vec{u}$ \hfill (élément neutre)
\item tout $\vec{u}$ admet un opposé $-\vec{u}$ tel que $\vec{u} + (-\vec{u}) = \vec{0}$ \hfill (opposé)
\end{enumerate}

\textbf{Famille 2 — la multiplication par un scalaire se comporte bien :}
\begin{enumerate}
\setcounter{enumi}{4}
\item $1 \cdot \vec{u} = \vec{u}$ \hfill (neutre pour le scalaire $1$)
\item $\lambda \cdot (\mu \cdot \vec{u}) = (\lambda\mu) \cdot \vec{u}$ \hfill (associativité mixte)
\end{enumerate}

\textbf{Famille 3 — les deux lois sont compatibles :}
\begin{enumerate}
\setcounter{enumi}{6}
\item $\lambda \cdot (\vec{u} + \vec{v}) = \lambda\cdot\vec{u} + \lambda\cdot\vec{v}$ \hfill (distributivité sur les vecteurs)
\item $(\lambda + \mu) \cdot \vec{u} = \lambda\cdot\vec{u} + \mu\cdot\vec{u}$ \hfill (distributivité sur les scalaires)
\end{enumerate}

Les éléments de $E$ sont appelés \cle{vecteurs}, les réels sont appelés \cle{scalaires}.
\end{definition}

\begin{attention}[Pas de panique devant les axiomes]
Ces huit axiomes ne sont pas à apprendre par cœur pour le Bac : ils expriment simplement que l'on peut \emph{calculer comme d'habitude} avec les vecteurs. Retiens l'idée essentielle : un espace vectoriel est un ensemble où l'on peut \textbf{additionner} les éléments et les \textbf{multiplier par des réels}, en restant dans l'ensemble et avec les règles de calcul usuelles.
\end{attention}

\begin{propriete}[Exemple fondamental]
$(\mathbb{R}^3, +, \cdot)$, muni de l'addition et de la multiplication par un scalaire définies coordonnée par coordonnée, est un espace vectoriel sur $\mathbb{R}$.
\end{propriete}

\begin{experience}[Vérification guidée pour $\mathbb{R}^3$]
Vérifions quelques axiomes pour bien comprendre le mécanisme. Soient $\vec{u} = (x,y,z)$, $\vec{v} = (x',y',z')$ et $\lambda, \mu \in \mathbb{R}$.
\begin{enumerate}
\item \textbf{Commutativité} : $\vec{u} + \vec{v} = (x+x', y+y', z+z') = (x'+x, y'+y, z'+z) = \vec{v}+\vec{u}$, car l'addition des réels est commutative.
\item \textbf{Élément neutre} : le triplet $\vec{0} = (0,0,0)$ vérifie $\vec{u} + \vec{0} = (x+0, y+0, z+0) = \vec{u}$.
\item \textbf{Opposé} : $-\vec{u} = (-x,-y,-z)$ convient car $\vec{u} + (-\vec{u}) = (0,0,0) = \vec{0}$.
\item \textbf{Distributivité} : $\lambda(\vec{u}+\vec{v}) = \lambda(x+x', y+y', z+z') = (\lambda x + \lambda x', \ldots) = \lambda\vec{u} + \lambda\vec{v}$.
\end{enumerate}
Les quatre autres axiomes se vérifient de la même façon, coordonnée par coordonnée, en utilisant les propriétés du calcul dans $\mathbb{R}$.
\end{experience}

\begin{exemplebox}[D'autres espaces vectoriels]
\begin{itemize}
\item $(\mathbb{R}^2, +, \cdot)$ : les couples de réels, c'est-à-dire les vecteurs du plan ;
\item $(\mathbb{R}, +, \cdot)$ : les réels eux-mêmes forment un espace vectoriel sur $\mathbb{R}$ ;
\item l'ensemble des vecteurs du plan muni de l'addition géométrique (règle de Chasles) et de la multiplication par un réel ;
\item l'ensemble des vecteurs de l'espace tout entier.
\end{itemize}
Dans cette leçon, notre terrain de jeu principal sera $\mathbb{R}^3$.
\end{exemplebox}

\courssub{Combinaisons linéaires de vecteurs}

L'outil central de toute cette leçon est la notion de combinaison linéaire : c'est l'opération qui mélange addition et multiplication par des scalaires.

\begin{definition}[Combinaison linéaire]
Soient $\vec{u}_1, \vec{u}_2, \ldots, \vec{u}_p$ des vecteurs d'un espace vectoriel $E$. On appelle \cle{combinaison linéaire} de ces vecteurs tout vecteur de la forme :
\[
\lambda_1 \vec{u}_1 + \lambda_2 \vec{u}_2 + \cdots + \lambda_p \vec{u}_p,
\]
où $\lambda_1, \lambda_2, \ldots, \lambda_p$ sont des réels appelés \cle{coefficients} de la combinaison linéaire.
\end{definition}

\begin{exemplebox}[Combinaisons linéaires dans $\mathbb{R}^3$]
Soient $\vec{u} = (1,0,2)$, $\vec{v} = (-1,3,1)$ et $\vec{w} = (0,1,1)$.
\begin{itemize}
\item $2\vec{u} + \vec{v} = (2,0,4) + (-1,3,1) = (1,3,5)$ est une combinaison linéaire de $\vec{u}$ et $\vec{v}$ ;
\item $\vec{u} - \vec{v} + 2\vec{w} = (1,0,2) - (-1,3,1) + (0,2,2) = (2,-1,3)$ est une combinaison linéaire de $\vec{u}$, $\vec{v}$ et $\vec{w}$ ;
\item le vecteur nul $\vec{0} = 0\cdot\vec{u} + 0\cdot\vec{v}$ est toujours une combinaison linéaire (coefficients tous nuls) de n'importe quelle famille de vecteurs.
\end{itemize}
\end{exemplebox}

\begin{savaistu}[Le langage secret des infographistes]
Quand l'infographiste de Yaoundé déplace sa case traditionnelle à l'écran, son logiciel calcule des milliers de combinaisons linéaires par seconde : chaque sommet du modèle 3D est un vecteur de $\mathbb{R}^3$, et chaque déplacement, rotation ou mise à l'échelle s'exprime avec des coefficients réels appliqués à ces vecteurs. Les mathématiques de cette leçon sont littéralement le moteur des images de synthèse.
\end{savaistu}

\courssub{Sous-espaces vectoriels : définition et caractérisation}

\textbf{Intuition.} À l'intérieur de l'espace $\mathbb{R}^3$, certains sous-ensembles « héritent » de la structure d'espace vectoriel : on peut y additionner et y multiplier par des scalaires sans jamais en sortir. Pense à une feuille de papier posée sur une table et passant par l'origine du repère : la somme de deux points de la feuille reste sur la feuille, et doubler un point de la feuille le laisse sur la feuille. La feuille est un « petit espace vectoriel » à l'intérieur du grand. C'est exactement la notion de sous-espace vectoriel.

\begin{definition}[Sous-espace vectoriel]
Soit $E$ un espace vectoriel sur $\mathbb{R}$. Une partie $F$ de $E$ est un \cle{sous-espace vectoriel} de $E$ si :
\begin{enumerate}
\item $F$ est non vide ;
\item $F$ est \textbf{stable par addition} : pour tous $\vec{u}, \vec{v} \in F$, on a $\vec{u} + \vec{v} \in F$ ;
\item $F$ est \textbf{stable par multiplication par un scalaire} : pour tout $\vec{u} \in F$ et tout $\lambda \in \mathbb{R}$, on a $\lambda \vec{u} \in F$.
\end{enumerate}
Un sous-espace vectoriel de $E$ est alors lui-même un espace vectoriel pour les mêmes lois.
\end{definition}

En pratique, on regroupe ces conditions en un seul test, plus rapide à rédiger.

\begin{propriete}[Caractérisation pratique par les combinaisons linéaires]
Soit $E$ un espace vectoriel et $F$ une partie de $E$. Alors :
\[
F \text{ est un sous-espace vectoriel de } E \iff
\begin{cases}
\vec{0} \in F,\\[2pt]
\forall\, \vec{u}, \vec{v} \in F,\ \forall\, \lambda \in \mathbb{R},\quad \lambda\vec{u} + \vec{v} \in F.
\end{cases}
\]
Autrement dit : $F$ contient le vecteur nul et est \textbf{stable par combinaisons linéaires}.
\end{propriete}

\begin{experience}[Pourquoi cette caractérisation est-elle équivalente ?]
Montrons les deux implications.
\begin{enumerate}
\item \textbf{Sens direct.} Supposons $F$ sous-espace vectoriel. Comme $F$ est non vide, il contient un vecteur $\vec{u}$. Par stabilité par multiplication, $0\cdot\vec{u} = \vec{0} \in F$. Ensuite, pour $\vec{u}, \vec{v} \in F$ et $\lambda \in \mathbb{R}$ : $\lambda\vec{u} \in F$ (stabilité externe), puis $\lambda\vec{u} + \vec{v} \in F$ (stabilité par addition).
\item \textbf{Sens réciproque.} Supposons les deux conditions de droite. Alors $F$ est non vide car $\vec{0} \in F$. Pour la stabilité par addition : avec $\lambda = 1$, $\vec{u}+\vec{v} = 1\cdot\vec{u} + \vec{v} \in F$. Pour la stabilité externe : avec $\vec{v} = \vec{0}$, $\lambda\vec{u} = \lambda\vec{u} + \vec{0} \in F$.
\end{enumerate}
Les deux définitions disent donc exactement la même chose.
\end{experience}

\begin{methode}[Montrer qu'un ensemble est un sous-espace vectoriel]
Pour prouver qu'une partie $F$ de $\mathbb{R}^3$ est un sous-espace vectoriel :
\begin{enumerate}
\item \textbf{Étape 1 — le vecteur nul.} Vérifier que $(0,0,0) \in F$ (il satisfait la condition définissant $F$). Si ce n'est pas le cas, c'est terminé : $F$ n'est \emph{pas} un sous-espace vectoriel.
\item \textbf{Étape 2 — stabilité.} Prendre deux éléments \emph{quelconques} $\vec{u} = (x,y,z)$ et $\vec{v} = (x',y',z')$ de $F$, un réel $\lambda$ quelconque, et montrer que $\lambda\vec{u} + \vec{v}$ vérifie encore la condition définissant $F$.
\item \textbf{Étape 3 — conclure} par la caractérisation : « $F$ est un sous-espace vectoriel de $\mathbb{R}^3$ ».
\end{enumerate}
\end{methode}

\begin{exoresolu}[Un plan d'équation $x + 2y - z = 0$]
Montrer que $F = \{(x,y,z) \in \mathbb{R}^3 \mid x + 2y - z = 0\}$ est un sous-espace vectoriel de $\mathbb{R}^3$.
\tcblower
\textbf{Étape 1 — le vecteur nul.} Pour $(0,0,0)$ : $0 + 2\times 0 - 0 = 0$. La condition est vérifiée, donc $\vec{0} \in F$.

\textbf{Étape 2 — stabilité par combinaison linéaire.} Soient $\vec{u} = (x,y,z) \in F$ et $\vec{v} = (x',y',z') \in F$, et $\lambda \in \mathbb{R}$. Par hypothèse :
\[
x + 2y - z = 0 \qquad \text{et} \qquad x' + 2y' - z' = 0.
\]
Calculons $\lambda\vec{u} + \vec{v} = (\lambda x + x',\ \lambda y + y',\ \lambda z + z')$ et testons la condition :
\begin{align*}
(\lambda x + x') + 2(\lambda y + y') - (\lambda z + z')
&= \lambda(x + 2y - z) + (x' + 2y' - z')\\
&= \lambda \times 0 + 0 = 0.
\end{align*}
Donc $\lambda\vec{u} + \vec{v} \in F$.

\textbf{Étape 3 — conclusion.} $F$ contient $\vec{0}$ et est stable par combinaisons linéaires : $F$ est un sous-espace vectoriel de $\mathbb{R}^3$. Géométriquement, $F$ est le plan passant par l'origine et de vecteur normal $(1,2,-1)$.
\end{exoresolu}

\courssub{Exemples et contre-exemples fondamentaux dans $\mathbb{R}^3$}

\begin{propriete}[Les sous-espaces vectoriels de $\mathbb{R}^3$]
Dans $\mathbb{R}^3$, on rencontre les sous-espaces vectoriels suivants :
\begin{itemize}
\item $\{\vec{0}\}$ : le singleton réduit au vecteur nul (le plus petit de tous) ;
\item les \cle{droites vectorielles} : ensembles $\{\lambda\vec{u} \mid \lambda \in \mathbb{R}\}$ où $\vec{u} \neq \vec{0}$, c'est-à-dire les droites passant par l'origine ;
\item les \cle{plans vectoriels} : ensembles $\{\lambda\vec{u} + \mu\vec{v} \mid \lambda, \mu \in \mathbb{R}\}$ où $\vec{u}$ et $\vec{v}$ ne sont pas colinéaires, c'est-à-dire les plans passant par l'origine ;
\item $\mathbb{R}^3$ tout entier (le plus grand de tous).
\end{itemize}
Nous verrons dans la section suivante que ce sont les \emph{seuls} sous-espaces vectoriels de $\mathbb{R}^3$.
\end{propriete}

\begin{propriete}[Sous-espace engendré par une famille de vecteurs]
Soient $\vec{u}_1, \dots, \vec{u}_p \in \mathbb{R}^3$. L'ensemble des combinaisons linéaires de ces vecteurs :
\[
\operatorname{Vect}(\vec{u}_1, \dots, \vec{u}_p) = \left\{ \lambda_1\vec{u}_1 + \cdots + \lambda_p\vec{u}_p \mid \lambda_1, \dots, \lambda_p \in \mathbb{R} \right\}
\]
est un sous-espace vectoriel de $\mathbb{R}^3$, appelé \cle{sous-espace engendré} (ou \cle{span}) par la famille $(\vec{u}_1, \dots, \vec{u}_p)$. C'est le plus petit sous-espace vectoriel contenant tous les $\vec{u}_i$.
\end{propriete}

\begin{exemplebox}[Une droite vectorielle]
Soit $\vec{u} = (1,-2,1)$ et $D = \{\lambda\vec{u} \mid \lambda \in \mathbb{R}\} = \{(\lambda, -2\lambda, \lambda) \mid \lambda \in \mathbb{R}\}$.
\begin{itemize}
\item $\vec{0} = 0\cdot\vec{u} \in D$ (avec $\lambda = 0$) ;
\item si $\vec{a} = \lambda\vec{u}$ et $\vec{b} = \mu\vec{u}$ sont dans $D$ et $\alpha \in \mathbb{R}$, alors $\alpha\vec{a} + \vec{b} = (\alpha\lambda + \mu)\vec{u} \in D$.
\end{itemize}
$D$ est donc un sous-espace vectoriel de $\mathbb{R}^3$ : c'est la droite passant par $O$ et dirigée par $\vec{u}$. On note $D = \operatorname{Vect}(\vec{u})$.
\end{exemplebox}

\begin{popfigure}
\begin{tikzpicture}[>=Stealth, scale=1.15]
% Base pour projection cavalière (x droite, y haut-gauche, z haut)
% Coordonnées 2D : (x - 0.4*y, 0.5*y + 0.8*z)
% Axes
\draw[->, popmuted] (-2.5, 0) -- (3.2, 0) node[below right] {$x$};
\draw[->, popmuted] (0.8, -1.0) -- (-1.2, 1.5) node[above left] {$y$};
\draw[->, popmuted] (0, -1.6) -- (0, 2.4) node[above] {$z$};
% Plan vectoriel P (z=0, engendré par x et y)
\fill[popblue!25, opacity=0.85] (-1.6, -0.5) -- (2.4, -0.5) -- (1.4, 0.75) -- (-2.6, 0.75) -- cycle;
\draw[popblue, thick] (-1.6, -0.5) -- (2.4, -0.5) -- (1.4, 0.75) -- (-2.6, 0.75) -- cycle;
\node[popblue] at (2.2, 0.55) {$P$};
% Droite vectorielle D (axe z)
\draw[poporange, very thick] (0, -1.6) -- (0, 2.0);
\node[poporange] at (0.1, 2.2) {$D$};
% Origine
\fill[popdark] (0,0) circle (1.6pt) node[below left=1pt] {$O$};
% Vecteurs du plan
\draw[->, popblue, thick] (0,0) -- (1.5, 0) node[midway, below] {\footnotesize $\vec{u}$};
\draw[->, popblue, thick] (0,0) -- (-0.48, 0.6) node[midway, left] {\footnotesize $\vec{v}$};
% Vecteur de la droite
\draw[->, poporange!80!popdark, thick] (0,0) -- (0, 1.2) node[pos=0.7, right] {\footnotesize $\vec{w}$};
\end{tikzpicture}
\legende{Le plan vectoriel $P$ (engendré par $\vec{u}$ et $\vec{v}$, ici le plan $z=0$) et la droite vectorielle $D$ (engendrée par $\vec{w}$, ici l'axe $z$) passent tous deux par l'origine : ce sont des sous-espaces vectoriels de $\mathbb{R}^3$.}
\end{popfigure}

\begin{attention}[L'erreur classique : oublier le vecteur nul]
L'erreur la plus fréquente au Bac est d'oublier de vérifier que $\vec{0} \in F$, ou pire : de croire qu'un plan quelconque est un sous-espace vectoriel. \textbf{Un plan (ou une droite) qui ne passe pas par l'origine n'est JAMAIS un sous-espace vectoriel}, puisqu'il ne contient pas $\vec{0}$. Par exemple, l'ensemble
\[
G = \{(x,y,z) \in \mathbb{R}^3 \mid x + 2y - z = 3\}
\]
n'est pas un sous-espace vectoriel : pour $(0,0,0)$, on obtient $0 + 0 - 0 = 0 \neq 3$, donc $\vec{0} \notin G$. C'est un plan \emph{affine}, parallèle au plan vectoriel $F$ de l'exercice résolu, mais ce n'est pas un espace vectoriel. Autre symptôme du même mal : si $\vec{u} \in G$, alors $2\vec{u}$ vérifie $2(x+2y-z) = 6 \neq 3$, donc $2\vec{u} \notin G$ : la stabilité externe échoue immédiatement.
\end{attention}

\begin{popfigure}
\begin{tikzpicture}[>=Stealth, scale=1.15]
% Base projection cavalière
% Axes
\draw[->, popmuted] (-2.5, 0) -- (3.5, 0) node[below right] {$x$};
\draw[->, popmuted] (0.8, -1.0) -- (-1.2, 1.5) node[above left] {$y$};
\draw[->, popmuted] (0, -2.5) -- (0, 2.0) node[above] {$z$};
% Plan affine G : x + 2y - z = 3
% Intercepts : A(3,0,0), B(0,1.5,0), C(0,0,-3)
% Coords 2D : A(3,0), B(-0.6, 0.75), C(0, -2.4)
\fill[popink!18] (3,0) -- (-0.6, 0.75) -- (0, -2.4) -- cycle;
\draw[popink, thick] (3,0) -- (-0.6, 0.75) -- (0, -2.4) -- cycle;
\node[popink] at (1.5, -0.5) {$G$};
% Croix sur le plan (marqueur "interdit")
\draw[poppink, line width=1.6pt] (0.8, -0.6) -- (1.6, -1.4);
\draw[poppink, line width=1.6pt] (0.8, -1.4) -- (1.6, -0.6);
% Origine
\fill[popdark] (0,0) circle (1.6pt) node[below left=1pt] {$O$};
% Flèche O vers plan
\draw[->, poppink, thick, dashed] (0,0) -- (0.8, -0.6) node[midway, right] {\footnotesize $\vec{0} \notin G$};
% Mention
\node[poppink, align=center] at (0.4, -2.0) {\footnotesize \textbf{pas un sous-espace vectoriel}};
\end{tikzpicture}
\legende{Le plan $G$ d'équation $x+2y-z=3$ (représenté par son triangle d'intercepts) ne contient pas l'origine : il n'est pas stable par multiplication par un scalaire et n'est pas un sous-espace vectoriel de $\mathbb{R}^3$.}
\end{popfigure}

\begin{exoresolu}[À toi de juger]
Les ensembles suivants sont-ils des sous-espaces vectoriels de $\mathbb{R}^3$ ?
\begin{enumerate}
\item $H = \{(x,y,z) \in \mathbb{R}^3 \mid x = y\}$ ;
\item $K = \{(x,y,z) \in \mathbb{R}^3 \mid z = 1\}$ ;
\item $L = \{(x,y,z) \in \mathbb{R}^3 \mid x \geq 0\}$.
\end{enumerate}
\tcblower
\textbf{1. Pour $H$ :} $(0,0,0)$ vérifie $0 = 0$, donc $\vec{0} \in H$. Si $\vec{u} = (x,y,z)$ et $\vec{v} = (x',y',z')$ sont dans $H$ (donc $x = y$ et $x' = y'$) et $\lambda \in \mathbb{R}$, alors $\lambda\vec{u}+\vec{v} = (\lambda x + x', \lambda y + y', \lambda z + z')$ vérifie $\lambda x + x' = \lambda y + y'$. Donc $\lambda\vec{u}+\vec{v} \in H$ : \textbf{$H$ est un sous-espace vectoriel} (c'est un plan vectoriel).

\textbf{2. Pour $K$ :} $(0,0,0)$ vérifie $z = 0 \neq 1$, donc $\vec{0} \notin K$ : \textbf{$K$ n'est pas un sous-espace vectoriel} (c'est un plan affine horizontal).

\textbf{3. Pour $L$ :} $\vec{0} \in L$ et la somme de deux éléments de $L$ reste dans $L$. Mais la stabilité externe échoue : $\vec{u} = (1,0,0) \in L$ et pourtant $(-1)\vec{u} = (-1,0,0) \notin L$ car $-1 < 0$. \textbf{$L$ n'est pas un sous-espace vectoriel}. Moralité : une inégalité large ne définit jamais un sous-espace vectoriel (sauf cas dégénérés).
\end{exoresolu}

\begin{aretenir}[L'essentiel de la section]
\begin{itemize}
\item Un \cle{espace vectoriel} sur $\mathbb{R}$ est un ensemble muni d'une addition et d'une multiplication par un scalaire vérifiant 8 axiomes naturels. L'exemple de référence est $(\mathbb{R}^3, +, \cdot)$.
\item Une \cle{combinaison linéaire} de vecteurs est une expression $\lambda_1\vec{u}_1 + \cdots + \lambda_p\vec{u}_p$.
\item $F$ est un \cle{sous-espace vectoriel} de $E$ si et seulement si $\vec{0} \in F$ et $\lambda\vec{u} + \vec{v} \in F$ pour tous $\vec{u}, \vec{v} \in F$ et $\lambda \in \mathbb{R}$.
\item Dans $\mathbb{R}^3$, les sous-espaces vectoriels sont : $\{\vec{0}\}$, les droites passant par $O$, les plans passant par $O$, et $\mathbb{R}^3$.
\item Le \cle{sous-espace engendré} par une famille $(\vec{u}_1,\dots,\vec{u}_p)$ est l'ensemble de toutes leurs combinaisons linéaires ; c'est le plus petit sous-espace vectoriel les contenant.
\item Un plan ou une droite \textbf{ne passant pas par l'origine} n'est jamais un sous-espace vectoriel : toujours tester $\vec{0}$ en premier !
\end{itemize}
\end{aretenir}

\begin{exercice}[Entraînement]
\begin{enumerate}
\item Soit $F = \{(x,y,z) \in \mathbb{R}^3 \mid 2x - y + 3z = 0\}$. Montrer que $F$ est un sous-espace vectoriel de $\mathbb{R}^3$.
\item L'ensemble $A = \{(x,y,z) \in \mathbb{R}^3 \mid x + y = 1\}$ est-il un sous-espace vectoriel de $\mathbb{R}^3$ ? Justifier.
\item Soient $\vec{u} = (1,1,0)$ et $\vec{v} = (0,1,1)$. Le vecteur $\vec{w} = (2,5,3)$ est-il une combinaison linéaire de $\vec{u}$ et $\vec{v}$ ?
\end{enumerate}
\end{exercice}

\begin{corrige}[Exercice]
\textbf{1.} Pour $(0,0,0)$ : $2(0) - 0 + 3(0) = 0$, donc $\vec{0} \in F$. Soient $\vec{u} = (x,y,z) \in F$, $\vec{v} = (x',y',z') \in F$ et $\lambda \in \mathbb{R}$. On a $2x - y + 3z = 0$ et $2x' - y' + 3z' = 0$. Alors pour $\lambda\vec{u}+\vec{v} = (\lambda x + x', \lambda y + y', \lambda z + z')$ :
\[
2(\lambda x + x') - (\lambda y + y') + 3(\lambda z + z') = \lambda(2x - y + 3z) + (2x' - y' + 3z') = 0.
\]
Donc $\lambda\vec{u}+\vec{v} \in F$ : $F$ est un sous-espace vectoriel de $\mathbb{R}^3$.

\textbf{2.} Pour $(0,0,0)$ : $0 + 0 = 0 \neq 1$, donc $\vec{0} \notin A$. $A$ n'est pas un sous-espace vectoriel de $\mathbb{R}^3$.

\textbf{3.} On cherche $\lambda, \mu \in \mathbb{R}$ tels que $\lambda\vec{u} + \mu\vec{v} = \vec{w}$, c'est-à-dire $(\lambda, \lambda + \mu, \mu) = (2,5,3)$. D'où $\lambda = 2$, $\mu = 3$, et vérification : $\lambda + \mu = 5$ : c'est cohérent. Donc $\vec{w} = 2\vec{u} + 3\vec{v}$ : oui, $\vec{w}$ est une combinaison linéaire de $\vec{u}$ et $\vec{v}$.
\end{corrige}

Dans la section suivante, nous répondrons à la question du technicien de Camrail : quand trois directions suffisent-elles à décrire tout l'espace ? Ce sera la notion de \cle{base} de $\mathbb{R}^3$.

\coursec{Familles libres, familles génératrices, bases}

Dans la section précédente, nous avons découvert l'espace vectoriel $\mathbb{R}^3$ et ses sous-espaces vectoriels : droites vectorielles, plans vectoriels, et $\mathbb{R}^3$ tout entier. Nous avons vu qu'un plan vectoriel est l'ensemble des \cle{combinaisons linéaires} de deux vecteurs non colinéaires, et qu'une droite vectorielle est l'ensemble des multiples d'un seul vecteur non nul. Une question surgit alors naturellement : \emph{combien faut-il de vecteurs, au minimum, pour décrire tout l'espace $\mathbb{R}^3$ ?} Et comment reconnaître qu'une famille de vecteurs « suffit » pour tout décrire, ou au contraire qu'elle contient des vecteurs « inutiles » ? C'est exactement l'objet de cette section : les notions de \cle{famille libre}, de \cle{famille génératrice} et de \cle{base} sont les outils qui permettent de répondre rigoureusement à ces questions.

Imagine un technicien de CAMTEL qui doit repérer la position d'un pylône dans l'espace. Il lui faut trois directions de référence indépendantes : si deux de ses directions sont alignées, ou si les trois sont dans un même plan, il ne pourra jamais atteindre tous les points de l'espace. Cette intuition physique — trois directions vraiment indépendantes — est précisément ce que les mathématiques appellent une \cle{base} de $\mathbb{R}^3$.

\courssub{Familles libres et familles liées}

\textbf{Intuition.} Dire que des vecteurs sont « libres » (ou linéairement indépendants), c'est dire qu'aucun d'entre eux ne peut être « fabriqué » à partir des autres. Chacun apporte une direction réellement nouvelle. À l'inverse, une famille est « liée » lorsque l'un de ses vecteurs est combinaison linéaire des autres : il est redondant, il n'apporte rien de nouveau.

\begin{definition}[Famille libre]
Soit $(\vec{u}_1, \vec{u}_2, \ldots, \vec{u}_p)$ une famille de vecteurs de $\mathbb{R}^3$. On dit que cette famille est \cle{libre} (ou que les vecteurs sont \cle{linéairement indépendants}) lorsque la seule combinaison linéaire nulle est celle dont tous les coefficients sont nuls :
\[ \alpha_1 \vec{u}_1 + \alpha_2 \vec{u}_2 + \cdots + \alpha_p \vec{u}_p = \vec{0} \implies \alpha_1 = \alpha_2 = \cdots = \alpha_p = 0. \]
\end{definition}

\begin{definition}[Famille liée]
Une famille de vecteurs qui n'est pas libre est dite \cle{liée} : il existe des coefficients $\alpha_1, \ldots, \alpha_p$ \emph{non tous nuls} tels que
\[ \alpha_1 \vec{u}_1 + \alpha_2 \vec{u}_2 + \cdots + \alpha_p \vec{u}_p = \vec{0}. \]
De manière équivalente, une famille est liée si et seulement si l'un de ses vecteurs est combinaison linéaire des autres.
\end{definition}

\textbf{Interprétation géométrique dans $\mathbb{R}^3$.} C'est le point essentiel à visualiser :
\begin{itemize}
\item Deux vecteurs forment une famille libre si et seulement s'ils ne sont \emph{pas colinéaires} : ils déterminent alors un plan.
\item Trois vecteurs forment une famille libre si et seulement s'ils ne sont \emph{pas coplanaires} : ils « sortent » véritablement du plan, ils occupent tout le volume de l'espace.
\end{itemize}

\begin{popfigure}
\begin{center}
\begin{tikzpicture}[scale=1.05, >=Stealth]
% ---- Panneau gauche : famille liée (coplanaires) ----
\begin{scope}[shift={(0,0)}]
\fill[popblueL, opacity=0.55] (-1.6,-0.7) -- (1.8,-0.4) -- (1.4,1.1) -- (-1.9,0.8) -- cycle;
\draw[->, thick, popblue] (0,0) -- (1.3,0.35) node[right] {\small $\vec{u}$};
\draw[->, thick, poporange] (0,0) -- (-0.9,0.55) node[above left] {\small $\vec{v}$};
\draw[->, thick, poppurple] (0,0) -- (0.7,0.85) node[above] {\small $\vec{w}$};
\fill[popdark] (0,0) circle (1.5pt) node[below left=-2pt] {\small $O$};
\node[popdark, align=center] at (0,-1.3) {\small Trois vecteurs \textbf{coplanaires}\\ \small famille \textbf{liée}};
\end{scope}
% ---- Panneau droit : famille libre ----
\begin{scope}[shift={(6.2,0)}]
\fill[popgreenL, opacity=0.5] (-1.6,-0.7) -- (1.8,-0.4) -- (1.4,1.1) -- (-1.9,0.8) -- cycle;
\draw[->, thick, popblue] (0,0) -- (1.3,0.35) node[right] {\small $\vec{u}$};
\draw[->, thick, poporange] (0,0) -- (-0.9,0.55) node[above left] {\small $\vec{v}$};
\draw[->, thick, popgreen] (0,0) -- (0.25,1.7) node[above] {\small $\vec{w}$};
\fill[popdark] (0,0) circle (1.5pt) node[below left=-2pt] {\small $O$};
\node[popdark, align=center] at (0,-1.3) {\small Trois vecteurs \textbf{non coplanaires}\\ \small famille \textbf{libre} (une base)};
\end{scope}
\end{tikzpicture}
\legende{À gauche, $\vec{w}$ reste dans le plan engendré par $\vec{u}$ et $\vec{v}$ : la famille est liée. À droite, $\vec{w}$ sort du plan : la famille est libre.}
\end{center}
\end{popfigure}

\begin{exemplebox}[Une famille liée]
Considérons $\vec{u} = (1, 2, 0)$, $\vec{v} = (0, 1, 1)$ et $\vec{w} = (2, 5, 1)$. On remarque que
\[ \vec{w} = 2\vec{u} + \vec{v} \quad \text{car} \quad 2(1,2,0) + (0,1,1) = (2, 5, 1). \]
Autrement dit $2\vec{u} + \vec{v} - \vec{w} = \vec{0}$ avec des coefficients $(2, 1, -1)$ non tous nuls : la famille $(\vec{u}, \vec{v}, \vec{w})$ est \cle{liée}. Le vecteur $\vec{w}$ est « en trop » : il vit déjà dans le plan engendré par $\vec{u}$ et $\vec{v}$.
\end{exemplebox}

\begin{attention}[Pièges sur la définition de famille libre]
\begin{itemize}
\item Pour prouver qu'une famille est \textbf{libre}, il ne suffit pas de vérifier que $0\vec{u}_1 + 0\vec{u}_2 + 0\vec{u}_3 = \vec{0}$ — cela est toujours vrai ! Il faut partir de l'égalité $\alpha_1\vec{u}_1 + \alpha_2\vec{u}_2 + \alpha_3\vec{u}_3 = \vec{0}$ avec des coefficients \emph{inconnus} et \emph{démontrer} qu'ils sont forcément tous nuls.
\item Toute famille contenant le vecteur nul est automatiquement liée : par exemple $(\vec{0}, \vec{u}, \vec{v})$ vérifie $1\cdot\vec{0} + 0\vec{u} + 0\vec{v} = \vec{0}$ avec un coefficient non nul.
\end{itemize}
\end{attention}

\courssub{Familles génératrices}

\textbf{Intuition.} Une famille est génératrice lorsque ses vecteurs, par combinaisons linéaires, permettent d'atteindre \emph{tous} les vecteurs de l'espace. Rien n'est hors de portée. Un seul vecteur non nul engendre une droite ; deux vecteurs non colinéaires engendrent un plan ; il en faut donc \emph{au moins trois} pour engendrer tout $\mathbb{R}^3$.

\begin{definition}[Famille génératrice]
Soit $(\vec{u}_1, \vec{u}_2, \ldots, \vec{u}_p)$ une famille de vecteurs de $\mathbb{R}^3$. On dit que cette famille est \cle{génératrice} de $\mathbb{R}^3$ lorsque tout vecteur $\vec{x}$ de $\mathbb{R}^3$ peut s'écrire comme combinaison linéaire des vecteurs de la famille :
\[ \forall \vec{x} \in \mathbb{R}^3,\ \exists\, \alpha_1, \ldots, \alpha_p \in \mathbb{R},\quad \vec{x} = \alpha_1 \vec{u}_1 + \alpha_2 \vec{u}_2 + \cdots + \alpha_p \vec{u}_p. \]
On dit aussi que la famille \emph{engendre} $\mathbb{R}^3$, et on note $\mathbb{R}^3 = \mathrm{Vect}(\vec{u}_1, \ldots, \vec{u}_p)$.
\end{definition}

\begin{exemplebox}[Deux vecteurs ne suffisent jamais]
Soient $\vec{u} = (1, 0, 0)$ et $\vec{v} = (0, 1, 0)$. Toute combinaison linéaire s'écrit
\[ \alpha \vec{u} + \beta \vec{v} = (\alpha, \beta, 0). \]
Le vecteur $\vec{x} = (0, 0, 1)$ n'a aucune chance d'être atteint : sa troisième coordonnée est $1 \neq 0$. La famille $(\vec{u}, \vec{v})$ engendre le plan d'équation $z = 0$, mais \emph{pas} $\mathbb{R}^3$. De manière générale, \textbf{une famille de deux vecteurs ne peut jamais engendrer $\mathbb{R}^3$}.
\end{exemplebox}

\courssub{Bases de $\mathbb{R}^3$}

\textbf{Intuition.} Une base est une famille « parfaite » : ni trop, ni trop peu. Elle est libre (aucun vecteur inutile) et génératrice (aucun vecteur manquant). C'est le système de référence idéal pour décrire l'espace.

\begin{definition}[Base d'un espace vectoriel]
Une \cle{base} de $\mathbb{R}^3$ est une famille de vecteurs à la fois \textbf{libre} et \textbf{génératrice}. De manière équivalente, $(\vec{u}, \vec{v}, \vec{w})$ est une base de $\mathbb{R}^3$ si et seulement si tout vecteur $\vec{x}$ de $\mathbb{R}^3$ s'écrit de \emph{façon unique} :
\[ \vec{x} = \alpha \vec{u} + \beta \vec{v} + \gamma \vec{w}, \qquad (\alpha, \beta, \gamma) \in \mathbb{R}^3. \]
\end{definition}

L'unicité est capitale : c'est elle qui permettra, dans la section suivante, de définir sans ambiguïté les \cle{coordonnées} d'un vecteur dans une base.

\begin{definition}[Base canonique de $\mathbb{R}^3$]
La \cle{base canonique} de $\mathbb{R}^3$ est la famille $(\vec{e}_1, \vec{e}_2, \vec{e}_3)$ définie par :
\[ \vec{e}_1 = (1, 0, 0), \qquad \vec{e}_2 = (0, 1, 0), \qquad \vec{e}_3 = (0, 0, 1). \]
Elle est bien une base : tout vecteur $\vec{x} = (x, y, z)$ s'écrit de façon unique $\vec{x} = x\vec{e}_1 + y\vec{e}_2 + z\vec{e}_3$.
\end{definition}

\begin{popfigure}
\begin{center}
\begin{tikzpicture}[scale=1.5, >=Stealth]
% axes pointillés (prolongements)
\draw[dashed, popmuted] (0,0) -- (-1.1,-0.55);
\draw[dashed, popmuted] (0,0) -- (-1.3,0.45);
\draw[dashed, popmuted] (0,0) -- (0,-1.0);
% vecteurs de la base canonique
\draw[->, very thick, popblue] (0,0) -- (1.9,-0.95) node[below right] {\small $\vec{e}_1=(1,0,0)$};
\draw[->, very thick, poporange] (0,0) -- (2.1,0.75) node[right] {\small $\vec{e}_2=(0,1,0)$};
\draw[->, very thick, popgreen] (0,0) -- (0,1.9) node[above] {\small $\vec{e}_3=(0,0,1)$};
% cube unité en pointillés pour matérialiser le volume
\draw[dotted, popmuted] (1.0,-0.5) -- (1.0,1.05) -- (0,1.55);
\draw[dotted, popmuted] (1.0,-0.5) -- (2.05,-0.14) -- (2.05,1.4) -- (1.0,1.05);
\draw[dotted, popmuted] (0,1.55) -- (1.05,1.9) -- (2.05,1.4);
\draw[dotted, popmuted] (1.05,0.4) -- (1.05,1.9);
\draw[dotted, popmuted] (2.05,-0.14) -- (1.05,0.4) -- (1.0,-0.5);
\fill[popdark] (0,0) circle (1.6pt) node[below left=-2pt] {\small $O$};
\end{tikzpicture}
\legende{La base canonique $(\vec{e}_1, \vec{e}_2, \vec{e}_3)$ : trois flèches unitaires le long des trois axes. Elles délimitent le cube unité, témoin du « volume » de l'espace.}
\end{center}
\end{popfigure}

\begin{propriete}[Théorème de la dimension (admis)]
Toutes les bases d'un même espace vectoriel ont le \emph{même nombre} de vecteurs. Ce nombre commun est appelé la \cle{dimension} de l'espace. Comme la base canonique de $\mathbb{R}^3$ possède trois vecteurs :
\[ \dim \mathbb{R}^3 = 3. \]
De même, une droite vectorielle est de dimension $1$ et un plan vectoriel de dimension $2$.
\end{propriete}

Ce théorème a une conséquence pratique formidable en dimension 3 : il suffit de vérifier \emph{une seule} des deux conditions (libre \emph{ou} génératrice) pour conclure qu'une famille de 3 vecteurs est une base.

\begin{propriete}[Propriété clé en dimension 3 (admise)]
Soit $(\vec{u}, \vec{v}, \vec{w})$ une famille de \textbf{trois} vecteurs de $\mathbb{R}^3$. Alors :
\begin{itemize}
\item si la famille est \textbf{libre}, alors elle est automatiquement génératrice, donc c'est une \textbf{base} ;
\item si la famille est \textbf{génératrice}, alors elle est automatiquement libre, donc c'est une \textbf{base}.
\end{itemize}
\end{propriete}

\textbf{Justification par les systèmes (idée de la preuve).} Écrivons $\vec{u} = (a_1, b_1, c_1)$, $\vec{v} = (a_2, b_2, c_2)$, $\vec{w} = (a_3, b_3, c_3)$. Dire que la famille est libre, c'est dire que le système homogène
\[ \begin{cases} a_1 \alpha + a_2 \beta + a_3 \gamma = 0 \\ b_1 \alpha + b_2 \beta + b_3 \gamma = 0 \\ c_1 \alpha + c_2 \beta + c_3 \gamma = 0 \end{cases} \]
n'admet que la solution triviale $(0, 0, 0)$. Or un système linéaire homogène de 3 équations à 3 inconnues n'a que la solution nulle si et seulement si son déterminant est non nul — et dans ce cas, le même système avec \emph{n'importe quel} second membre $(x, y, z)$ admet une solution unique : c'est exactement dire que la famille est aussi génératrice. Liberté et génération sont donc deux visages d'une même réalité algébrique en dimension 3.

\begin{propriete}[Critère du déterminant]
Une famille de trois vecteurs $\vec{u} = (a_1, b_1, c_1)$, $\vec{v} = (a_2, b_2, c_2)$, $\vec{w} = (a_3, b_3, c_3)$ est une base de $\mathbb{R}^3$ si et seulement si
\[ \det(\vec{u}, \vec{v}, \vec{w}) = \begin{vmatrix} a_1 & a_2 & a_3 \\ b_1 & b_2 & b_3 \\ c_1 & c_2 & c_3 \end{vmatrix} \neq 0. \]
Géométriquement, ce déterminant est (au signe près) le \emph{volume} du parallélépipède construit sur les trois vecteurs : il est nul exactement quand les vecteurs sont coplanaires, c'est-à-dire quand le parallélépipède est « aplati ».
\end{propriete}

\courssub{Méthode de travail et exemple fondateur}

\begin{methode}[Montrer qu'une famille de 3 vecteurs est une base de $\mathbb{R}^3$]
Pour montrer que $(\vec{u}, \vec{v}, \vec{w})$ est une base de $\mathbb{R}^3$ :
\begin{enumerate}
\item \textbf{Poser} l'équation aux coefficients inconnus : $\alpha \vec{u} + \beta \vec{v} + \gamma \vec{w} = \vec{0}$.
\item \textbf{Traduire} en système de 3 équations à 3 inconnues $\alpha, \beta, \gamma$ en identifiant coordonnée par coordonnée.
\item \textbf{Résoudre} le système (substitution, combinaison ou pivot).
\item \textbf{Conclure} : si l'unique solution est $\alpha = \beta = \gamma = 0$, la famille est libre ; comme elle compte 3 vecteurs dans $\mathbb{R}^3$ (de dimension 3), c'est une \textbf{base}.
\end{enumerate}
Variante rapide : calculer le déterminant des trois vecteurs ; s'il est non nul, c'est une base.
\end{methode}

\begin{exoresolu}[Une base classique de $\mathbb{R}^3$]
Montrer que les vecteurs $\vec{u} = (1, 1, 0)$, $\vec{v} = (0, 1, 1)$ et $\vec{w} = (1, 0, 1)$ forment une base de $\mathbb{R}^3$.
\tcblower
\textbf{Solution détaillée.} La famille compte 3 vecteurs dans $\mathbb{R}^3$ : il suffit de montrer qu'elle est libre.

\emph{Étape 1.} Soient $\alpha, \beta, \gamma \in \mathbb{R}$ tels que $\alpha \vec{u} + \beta \vec{v} + \gamma \vec{w} = \vec{0}$.

\emph{Étape 2.} En coordonnées :
\[ \alpha(1,1,0) + \beta(0,1,1) + \gamma(1,0,1) = (0,0,0) \iff \begin{cases} \alpha + \gamma = 0 & (L_1) \\ \alpha + \beta = 0 & (L_2) \\ \beta + \gamma = 0 & (L_3) \end{cases} \]

\emph{Étape 3.} Résolution. De $(L_1)$ : $\gamma = -\alpha$. De $(L_2)$ : $\beta = -\alpha$. En reportant dans $(L_3)$ :
\[ \beta + \gamma = -\alpha - \alpha = -2\alpha = 0 \implies \alpha = 0, \]
puis $\beta = 0$ et $\gamma = 0$.

\emph{Étape 4.} Conclusion : la seule combinaison linéaire nulle est celle à coefficients nuls, donc la famille $(\vec{u}, \vec{v}, \vec{w})$ est \textbf{libre}. Étant une famille libre de 3 vecteurs dans $\mathbb{R}^3$ qui est de dimension 3, c'est une \textbf{base} de $\mathbb{R}^3$.

\emph{Vérification par le déterminant} (règle de Sarrus) :
\[ \begin{vmatrix} 1 & 0 & 1 \\ 1 & 1 & 0 \\ 0 & 1 & 1 \end{vmatrix} = (1 \times 1 \times 1 + 0 \times 0 \times 0 + 1 \times 1 \times 1) - (1 \times 1 \times 0 + 0 \times 1 \times 1 + 1 \times 0 \times 1) = 2 \neq 0. \]
Le déterminant vaut $2 \neq 0$ : on confirme que c'est une base (et le parallélépipède construit sur ces trois vecteurs a pour volume $2$).
\end{exoresolu}

\begin{attention}[Ne pas confondre « libre » et « génératrice »]
Ce sont deux propriétés \emph{distinctes}, et l'erreur la plus fréquente au Baccalauréat est de les mélanger :
\begin{itemize}
\item \textbf{Libre} concerne l'\emph{unicité} : aucun vecteur n'est redondant. Test : $\alpha\vec{u} + \beta\vec{v} + \gamma\vec{w} = \vec{0}$ force $\alpha = \beta = \gamma = 0$.
\item \textbf{Génératrice} concerne l'\emph{existence} : tout vecteur est atteignable. Test : l'équation $\alpha\vec{u} + \beta\vec{v} + \gamma\vec{w} = \vec{x}$ a une solution pour tout $\vec{x}$.
\end{itemize}
Retiens les contraintes de dimension : dans $\mathbb{R}^3$, une famille de \textbf{2 vecteurs ne peut jamais être génératrice} (elle engendre au plus un plan), et une famille de \textbf{4 vecteurs ou plus est toujours liée}. Seule une famille de \emph{exactement 3 vecteurs} peut prétendre être une base — et encore faut-il vérifier la liberté (ou le caractère générateur, ou le déterminant non nul).
\end{attention}

\begin{exercice}[À toi de jouer]
\begin{enumerate}
\item Montrer que la famille $\big(\vec{u}, \vec{v}, \vec{w}\big)$ avec $\vec{u} = (1, 0, 1)$, $\vec{v} = (1, 1, 0)$, $\vec{w} = (0, 1, 1)$ est une base de $\mathbb{R}^3$.
\item La famille $\big((1, 2, 3),\ (4, 5, 6),\ (7, 8, 9)\big)$ est-elle une base de $\mathbb{R}^3$ ? Justifier.
\end{enumerate}
\end{exercice}

\begin{corrige}[Exercice 1]
\textbf{1.} Soit $\alpha \vec{u} + \beta \vec{v} + \gamma \vec{w} = \vec{0}$. Le système associé est
\[ \begin{cases} \alpha + \beta = 0 \\ \beta + \gamma = 0 \\ \alpha + \gamma = 0 \end{cases} \]
De la première équation $\beta = -\alpha$ ; de la deuxième $\gamma = -\beta = \alpha$ ; la troisième donne $\alpha + \alpha = 2\alpha = 0$, donc $\alpha = \beta = \gamma = 0$. La famille est libre ; étant de cardinal 3 dans $\mathbb{R}^3$ de dimension 3, c'est une base.

\textbf{2.} On remarque que $(7, 8, 9) = 2(4, 5, 6) - (1, 2, 3)$ : le troisième vecteur est combinaison linéaire des deux premiers. La famille est donc \textbf{liée} (relation $-\vec{u} + 2\vec{v} - \vec{w} = \vec{0}$ à coefficients non tous nuls) : ce n'est \emph{pas} une base de $\mathbb{R}^3$. Les trois vecteurs sont coplanaires.
\end{corrige}

\begin{savaistu}[Pourquoi dit-on que l'espace a « trois dimensions » ?]
La notion de \cle{dimension} formalisée par les mathématiciens (Grassmann, puis Peano qui donna la définition axiomatique des espaces vectoriels en 1888) capture exactement l'idée physique que notre espace est « à 3 dimensions » : il faut \emph{exactement trois} nombres pour repérer un point — longitude, latitude et altitude, comme le fait le GPS pour localiser un taxi à Douala. Ni deux (il manquerait la hauteur), ni quatre (le quatrième serait superflu). C'est pourquoi $\dim \mathbb{R}^3 = 3$ n'est pas une convention arbitraire : c'est le reflet mathématique d'une propriété fondamentale du monde qui nous entoure. La relativité d'Einstein ajoutera plus tard une quatrième coordonnée — le temps — et l'algèbre linéaire en dimension 4 sera prête pour l'accueillir !
\end{savaistu}

\begin{aretenir}[L'essentiel de la section]
\begin{itemize}
\item Une famille est \cle{libre} si $\alpha_1\vec{u}_1 + \cdots + \alpha_p\vec{u}_p = \vec{0}$ impose que tous les coefficients sont nuls. Dans $\mathbb{R}^3$ : 2 vecteurs libres = non colinéaires ; 3 vecteurs libres = non coplanaires.
\item Une famille est \cle{génératrice} si tout vecteur de $\mathbb{R}^3$ est combinaison linéaire de ses vecteurs.
\item Une \cle{base} est une famille libre \emph{et} génératrice : tout vecteur s'écrit alors de façon \emph{unique} comme combinaison linéaire des vecteurs de la base.
\item La \cle{base canonique} de $\mathbb{R}^3$ est $(\vec{e}_1, \vec{e}_2, \vec{e}_3)$ avec $\vec{e}_1 = (1,0,0)$, $\vec{e}_2 = (0,1,0)$, $\vec{e}_3 = (0,0,1)$, et $\dim \mathbb{R}^3 = 3$.
\item En dimension 3, pour une famille de \textbf{3 vecteurs} : libre $\iff$ génératrice $\iff$ base $\iff$ déterminant non nul. Il suffit donc de vérifier \emph{une seule} de ces conditions.
\end{itemize}
\end{aretenir}

\coursec{Coordonnées d'un vecteur dans une base}

Dans la section précédente, nous avons découvert les \cle{bases} de $\mathbb{R}^3$ : ces familles de trois vecteurs à la fois \cle{libres} et \cle{génératrices}, qui forment en quelque sorte le « squelette minimal » permettant de fabriquer tous les vecteurs de l'espace. Nous allons maintenant exploiter cette double propriété pour répondre à une question fondamentale : comment \emph{repérer} un vecteur de l'espace à l'aide de nombres ? C'est toute la puissance de la notion de \cle{coordonnées}, que tu utilises déjà depuis la classe de Seconde sans forcément en connaître le fondement théorique.

\courssub{Intuition : une adresse unique pour chaque vecteur}

Imagine un livreur de l'entreprise \emph{Express Union} qui doit décrire la position d'un colis dans un grand entrepôt de Douala. Il peut dire : « avance de 2 pas vers l'est, puis 3 pas vers le nord, puis monte de 1 étage ». Le triplet $(2\,;\,3\,;\,1)$ décrit alors \emph{parfaitement} le déplacement, à condition que tout le monde soit d'accord sur les trois directions de référence « est », « nord », « haut ». Ces trois directions jouent le rôle d'une \emph{base}, et le triplet $(2\,;\,3\,;\,1)$ est l'\emph{adresse} du vecteur déplacement dans cette base.

Deux questions se posent alors naturellement :
\begin{itemize}
\item Tout déplacement peut-il être décrit ainsi ? (question d'\textbf{existence})
\item Un même déplacement peut-il avoir deux descriptions différentes dans les mêmes directions de référence ? (question d'\textbf{unicité})
\end{itemize}

C'est précisément le contenu du théorème central de cette section : la réponse est \emph{oui} à la première question grâce au caractère \textbf{générateur} de la base, et \emph{non} à la deuxième grâce à sa \textbf{liberté}. Une base, c'est exactement l'outil qui garantit les deux à la fois.

\courssub{Théorème fondamental : existence et unicité des coordonnées}

\begin{propriete}[Théorème-définition : coordonnées d'un vecteur dans une base]
Soit $\mathcal{B} = (\vec{u}, \vec{v}, \vec{w})$ une base de l'espace vectoriel $\mathbb{R}^3$. Alors, pour tout vecteur $\vec{x}$ de $\mathbb{R}^3$, il existe un \textbf{unique} triplet $(\alpha, \beta, \gamma)$ de réels tel que :
\[ \vec{x} = \alpha \vec{u} + \beta \vec{v} + \gamma \vec{w}. \]
Les réels $\alpha$, $\beta$, $\gamma$ sont appelés les \cle{coordonnées} (ou \emph{composantes}) du vecteur $\vec{x}$ dans la base $\mathcal{B}$. On note parfois :
\[ \vec{x}_{\mathcal{B}} = \begin{pmatrix} \alpha \\ \beta \\ \gamma \end{pmatrix} \quad \text{ou} \quad \vec{x}\,(\alpha\,;\,\beta\,;\,\gamma)_{\mathcal{B}}. \]
\end{propriete}

Ce théorème est \emph{exigible} au Baccalauréat : il faut savoir l'énoncer avec précision et en comprendre la démonstration. La beauté de ce résultat est qu'il repose entièrement sur les deux propriétés définissant une base, chacune servant exactement une fois.

\begin{experience}[Démonstration guidée : l'existence des coordonnées]
Soit $\mathcal{B} = (\vec{u}, \vec{v}, \vec{w})$ une base de $\mathbb{R}^3$ et soit $\vec{x}$ un vecteur quelconque de $\mathbb{R}^3$. Nous voulons montrer qu'il existe des réels $\alpha, \beta, \gamma$ tels que $\vec{x} = \alpha \vec{u} + \beta \vec{v} + \gamma \vec{w}$.

\textbf{Étape 1.} Rappelle la définition d'une famille génératrice : la famille $(\vec{u}, \vec{v}, \vec{w})$ est génératrice de $\mathbb{R}^3$ si tout vecteur de $\mathbb{R}^3$ est une combinaison linéaire de $\vec{u}$, $\vec{v}$, $\vec{w}$.

\textbf{Étape 2.} Or $\mathcal{B}$ est une base, donc une famille génératrice de $\mathbb{R}^3$.

\textbf{Étape 3.} Comme $\vec{x} \in \mathbb{R}^3$, il existe donc des réels $\alpha, \beta, \gamma$ tels que $\vec{x} = \alpha \vec{u} + \beta \vec{v} + \gamma \vec{w}$. \hfill $\blacksquare$

L'existence est donc une conséquence \emph{directe} du caractère générateur : il n'y a rien de plus à démontrer !
\end{experience}

\begin{experience}[Démonstration guidée : l'unicité des coordonnées]
Montrons maintenant que le triplet $(\alpha, \beta, \gamma)$ est unique. La méthode classique consiste à supposer qu'il existe \emph{deux} décompositions et à prouver qu'elles coïncident.

\textbf{Étape 1.} Supposons que $\vec{x}$ admette deux écritures :
\[ \vec{x} = \alpha \vec{u} + \beta \vec{v} + \gamma \vec{w} \quad \text{et} \quad \vec{x} = \alpha' \vec{u} + \beta' \vec{v} + \gamma' \vec{w}. \]

\textbf{Étape 2.} En soustrayant membre à membre ces deux égalités, on obtient :
\[ \vec{0} = (\alpha - \alpha')\,\vec{u} + (\beta - \beta')\,\vec{v} + (\gamma - \gamma')\,\vec{w}. \]

\textbf{Étape 3.} Or la famille $(\vec{u}, \vec{v}, \vec{w})$ est \textbf{libre} (car c'est une base). Par définition de la liberté, la seule combinaison linéaire de $\vec{u}, \vec{v}, \vec{w}$ égale au vecteur nul est celle dont tous les coefficients sont nuls. Donc :
\[ \alpha - \alpha' = 0, \qquad \beta - \beta' = 0, \qquad \gamma - \gamma' = 0, \]
c'est-à-dire $\alpha = \alpha'$, $\beta = \beta'$, $\gamma = \gamma'$.

\textbf{Étape 4.} Conclusion : les deux décompositions sont identiques, la décomposition de $\vec{x}$ est unique. \hfill $\blacksquare$

Retiens le schéma logique : \emph{existence} $\Leftarrow$ famille \textbf{génératrice} ; \emph{unicité} $\Leftarrow$ famille \textbf{libre}. Une base est précisément la famille qui offre les deux.
\end{experience}

\begin{aretenir}[L'idée à retenir]
Dans une base $\mathcal{B} = (\vec{u}, \vec{v}, \vec{w})$, chaque vecteur $\vec{x}$ possède une \emph{adresse unique} : le triplet de ses coordonnées. C'est ce qui permet de traduire toute la géométrie de l'espace en calculs sur des triplets de nombres réels.
\end{aretenir}

\courssub{Méthode de calcul : déterminer des coordonnées revient à résoudre un système}

Concrètement, comment trouver les coordonnées d'un vecteur $\vec{x}$ dans une base $\mathcal{B} = (\vec{u}, \vec{v}, \vec{w})$ ? On écrit l'égalité vectorielle $\vec{x} = \alpha \vec{u} + \beta \vec{v} + \gamma \vec{w}$, on la traduit composante par composante, et on obtient un \textbf{système linéaire de 3 équations à 3 inconnues} $\alpha, \beta, \gamma$, que l'on résout par substitution, combinaison ou pivot de Gauss.

\begin{methode}[Déterminer les coordonnées d'un vecteur dans une base]
\begin{enumerate}
\item \textbf{Écrire} l'égalité vectorielle $\vec{x} = \alpha \vec{u} + \beta \vec{v} + \gamma \vec{w}$ avec des inconnues $\alpha, \beta, \gamma$.
\item \textbf{Traduire} cette égalité en trois équations scalaires, une par composante (dans la base canonique).
\item \textbf{Résoudre} le système $3 \times 3$ obtenu (substitution, combinaisons linéaires ou pivot de Gauss).
\item \textbf{Vérifier} en recomposant : calculer $\alpha \vec{u} + \beta \vec{v} + \gamma \vec{w}$ et contrôler que l'on retombe bien sur $\vec{x}$.
\item \textbf{Conclure} en rédigeant proprement : « les coordonnées de $\vec{x}$ dans la base $\mathcal{B}$ sont $(\alpha\,;\,\beta\,;\,\gamma)$ ».
\end{enumerate}
\end{methode}

\begin{exemplebox}[Exemple chiffré détaillé]
On considère les vecteurs $\vec{u} = (1\,;\,1\,;\,0)$, $\vec{v} = (0\,;\,1\,;\,1)$, $\vec{w} = (1\,;\,0\,;\,1)$. On admet (on peut le vérifier) que $\mathcal{B} = (\vec{u}, \vec{v}, \vec{w})$ est une base de $\mathbb{R}^3$. Déterminons les coordonnées de $\vec{x} = (5\,;\,3\,;\,4)$ dans $\mathcal{B}$.

\textbf{Étape 1.} On écrit $\vec{x} = \alpha \vec{u} + \beta \vec{v} + \gamma \vec{w}$.

\textbf{Étape 2.} Composante par composante :
\[ \alpha \begin{pmatrix} 1 \\ 1 \\ 0 \end{pmatrix} + \beta \begin{pmatrix} 0 \\ 1 \\ 1 \end{pmatrix} + \gamma \begin{pmatrix} 1 \\ 0 \\ 1 \end{pmatrix} = \begin{pmatrix} 5 \\ 3 \\ 4 \end{pmatrix}
\iff
\begin{cases}
\alpha + \gamma = 5 & (L_1)\\
\alpha + \beta = 3 & (L_2)\\
\beta + \gamma = 4 & (L_3)
\end{cases} \]

\textbf{Étape 3.} Résolution par combinaison. Calculons $(L_1) + (L_2) - (L_3)$ :
\[ (\alpha + \gamma) + (\alpha + \beta) - (\beta + \gamma) = 5 + 3 - 4 \iff 2\alpha = 4 \iff \alpha = 2. \]
De $(L_2)$ : $\beta = 3 - \alpha = 1$. De $(L_1)$ : $\gamma = 5 - \alpha = 3$.

\textbf{Étape 4.} Vérification : $2\vec{u} + 1\vec{v} + 3\vec{w} = (2\,;\,2\,;\,0) + (0\,;\,1\,;\,1) + (3\,;\,0\,;\,3) = (5\,;\,3\,;\,4) = \vec{x}$. \checkmark

\textbf{Conclusion.} Les coordonnées de $\vec{x}$ dans la base $\mathcal{B}$ sont $(2\,;\,1\,;\,3)$, c'est-à-dire $\vec{x} = 2\vec{u} + \vec{v} + 3\vec{w}$.
\end{exemplebox}

Remarque essentielle : le \emph{même} vecteur $\vec{x}$ s'écrit $(5\,;\,3\,;\,4)$ dans la base canonique et $(2\,;\,1\,;\,3)$ dans la base $\mathcal{B}$. Le vecteur n'a pas changé ; c'est le \emph{système de référence} qui a changé.

\courssub{Interprétation géométrique : le parallélépipède des composantes}

Géométriquement, écrire $\vec{x} = \alpha \vec{u} + \beta \vec{v} + \gamma \vec{w}$ signifie que l'on peut atteindre l'extrémité de $\vec{x}$ en parcourant $\alpha$ fois le vecteur $\vec{u}$, puis $\beta$ fois $\vec{v}$, puis $\gamma$ fois $\vec{w}$. Les trois segments $\alpha\vec{u}$, $\beta\vec{v}$, $\gamma\vec{w}$ sont les arêtes d'un \emph{parallélépipède} dont $\vec{x}$ est la grande diagonale. C'est la généralisation naturelle du rectangle et de sa diagonale que tu connais dans le plan.

\begin{popfigure}
\begin{center}
\begin{tikzpicture}[scale=1.15, line join=round,
  axis/.style={-{Stealth[length=2.5mm]}, thick}]
% Origine
\coordinate (O) at (0,0);
% Vecteurs de base (directions)
\coordinate (U) at (2.2,0);      % u
\coordinate (V) at (0.9,1.5);    % v
\coordinate (W) at (-0.9,0.9);   % w
% Composantes : arêtes du parallélépipède
\coordinate (A) at (2.2,0);            % alpha u
\coordinate (B) at (0.9,1.5);          % beta v
\coordinate (C) at (-0.9,0.9);         % gamma w
% Sommet final x = A + B + C
\coordinate (X) at ($(A)+(B)+(C)$);
% Autres sommets du parallélépipède
\coordinate (AB) at ($(A)+(B)$);
\coordinate (AC) at ($(A)+(C)$);
\coordinate (BC) at ($(B)+(C)$);
% Faces (transparence légère)
\fill[popblueL, opacity=0.45] (O) -- (A) -- (AB) -- (B) -- cycle;
\fill[poporangeL, opacity=0.40] (O) -- (A) -- (AC) -- (C) -- cycle;
\fill[popgreenL, opacity=0.35] (A) -- (AB) -- (X) -- (AC) -- cycle;
% Arêtes du parallélépipède
\draw[popmuted, dashed] (O) -- (C);
\draw[popmuted, dashed] (C) -- (AC);
\draw[popmuted, dashed] (C) -- (BC);
\draw[popmuted] (O) -- (A) -- (AB) -- (B) -- cycle;
\draw[popmuted] (A) -- (AC);
\draw[popmuted] (AB) -- (X);
\draw[popmuted] (AC) -- (X);
\draw[popmuted] (B) -- (BC) -- (X);
% Vecteurs de base
\draw[axis, popblue] (O) -- ($0.55*(U)$) node[below right, scale=0.9] {$\vec{u}$};
\draw[axis, popgreen] (O) -- ($0.62*(V)$) node[above, scale=0.9] {$\vec{v}$};
\draw[axis, poppurple] (O) -- ($0.75*(W)$) node[left, scale=0.9] {$\vec{w}$};
% Composantes sur les arêtes
\draw[-{Stealth[length=2.5mm]}, very thick, popblue] (O) -- (A) node[midway, below, scale=0.9] {$\alpha\vec{u}$};
\draw[-{Stealth[length=2.5mm]}, very thick, popgreen] (A) -- (AB) node[midway, right, scale=0.9] {$\beta\vec{v}$};
\draw[-{Stealth[length=2.5mm]}, very thick, poppurple] (AB) -- (X) node[midway, right, scale=0.9] {$\gamma\vec{w}$};
% Vecteur x (grande diagonale)
\draw[-{Stealth[length=3mm]}, line width=1.6pt, popink] (O) -- (X) node[pos=0.55, above left, scale=0.95] {$\vec{x}$};
% Points
\fill[popdark] (O) circle (1.6pt) node[below left, scale=0.9] {$O$};
\fill[popink] (X) circle (1.8pt) node[above right, scale=0.9] {$M$};
\end{tikzpicture}
\end{center}
\legende{Décomposition d'un vecteur $\vec{x}$ selon les trois vecteurs d'une base $(\vec{u}, \vec{v}, \vec{w})$ : le point $M$ tel que $\overrightarrow{OM} = \vec{x}$ est atteint en parcourant successivement $\alpha\vec{u}$, $\beta\vec{v}$ et $\gamma\vec{w}$. Le vecteur $\vec{x}$ est la grande diagonale du parallélépipède construit sur ces trois composantes.}
\end{popfigure}

\courssub{Cas particulier : la base canonique}

\begin{definition}[Base canonique de $\mathbb{R}^3$]
On appelle \cle{base canonique} de $\mathbb{R}^3$ la famille $(\vec{i}, \vec{j}, \vec{k})$ définie par :
\[ \vec{i} = (1\,;\,0\,;\,0), \qquad \vec{j} = (0\,;\,1\,;\,0), \qquad \vec{k} = (0\,;\,0\,;\,1). \]
\end{definition}

\begin{propriete}[Coordonnées dans la base canonique]
Dans la base canonique, les coordonnées d'un vecteur sont exactement ses composantes usuelles : si $\vec{x} = (a\,;\,b\,;\,c)$, alors
\[ \vec{x} = a\,\vec{i} + b\,\vec{j} + c\,\vec{k}, \]
autrement dit les coordonnées de $\vec{x}$ dans la base canonique sont $(a\,;\,b\,;\,c)$.
\end{propriete}

La vérification est immédiate : $a\vec{i} + b\vec{j} + c\vec{k} = a(1\,;\,0\,;\,0) + b(0\,;\,1\,;\,0) + c(0\,;\,0\,;\,1) = (a\,;\,b\,;\,c)$. C'est pourquoi, depuis la Seconde, tu confonds spontanément « composantes » et « coordonnées » : dans la base canonique, les deux notions coïncident. Mais attention, cette coïncidence est un \emph{privilège} de la base canonique !

\begin{attention}[Erreur fréquente : confondre composantes canoniques et coordonnées dans une base]
Lorsqu'un énoncé demande les coordonnées d'un vecteur $\vec{x} = (5\,;\,3\,;\,4)$ \emph{dans une base $\mathcal{B}$ donnée}, il est \textbf{faux} de répondre $(5\,;\,3\,;\,4)$ : ce sont les coordonnées dans la base \emph{canonique}, pas dans $\mathcal{B}$ ! Dans notre exemple, la bonne réponse était $(2\,;\,1\,;\,3)$.

Deux réflexes pour éviter ce piège au Bac :
\begin{itemize}
\item Repère dans l'énoncé la phrase « dans la base $(\vec{u}, \vec{v}, \vec{w})$ » : elle t'oblige à poser et résoudre le système $\vec{x} = \alpha\vec{u} + \beta\vec{v} + \gamma\vec{w}$.
\item Vérifie toujours ta réponse en recomposant : si $\alpha\vec{u} + \beta\vec{v} + \gamma\vec{w} \neq \vec{x}$, ta réponse est fausse.
\end{itemize}
Autre piège de rédaction : l'\textbf{ordre} des vecteurs de la base compte. Dans la base $(\vec{v}, \vec{u}, \vec{w})$, les coordonnées de notre exemple seraient $(1\,;\,2\,;\,3)$ et non $(2\,;\,1\,;\,3)$. Une base est une famille \emph{ordonnée} de vecteurs.
\end{attention}

\courssub{Changement de point de vue : un vecteur, plusieurs adresses}

Un même vecteur possède des coordonnées différentes dans des bases différentes. C'est exactement comme la position d'un lieu à Yaoundé : on peut la décrire par rapport au carrefour Poste Centrale, ou par rapport au rond-point Nlongkak ; le lieu ne bouge pas, mais sa description change.

Illustrons cela avec notre exemple. Le vecteur $\vec{x}$ vérifie :
\begin{align*}
\vec{x} &= 5\vec{i} + 3\vec{j} + 4\vec{k} && \text{coordonnées } (5\,;\,3\,;\,4) \text{ dans } (\vec{i},\vec{j},\vec{k}),\\
\vec{x} &= 2\vec{u} + \vec{v} + 3\vec{w} && \text{coordonnées } (2\,;\,1\,;\,3) \text{ dans } (\vec{u},\vec{v},\vec{w}).
\end{align*}

Ce phénomène n'est pas un défaut, c'est une \emph{richesse} : dans de nombreux problèmes (symétries, rotations, projections, et plus tard les applications linéaires), un choix de base astucieux simplifie considérablement les calculs. Tout l'art du mathématicien consiste à choisir le « bon » système de coordonnées pour le problème donné.

\begin{savaistu}[Pourquoi changer de base ? Un exemple concret]
En géographie et en cartographie, les ingénieurs de l'Institut National de Cartographie du Cameroun utilisent des systèmes de coordonnées adaptés au terrain : la position d'un pylône peut être donnée en longitude-latitude, mais aussi dans un repère local aligné sur la vallée. De même, en mécanique, pour étudier le mouvement d'un véhicule sur une route de montagne comme l'axe Dschang--Bafoussam, on choisit une base dont un vecteur est tangent à la route : les équations deviennent bien plus simples. Changer de base, c'est changer de lunettes pour mieux voir le problème.
\end{savaistu}

\courssub{Lien avec la géométrie : coordonnées d'un point dans un repère de l'espace}

Tout ce qui précède se transpose aux \emph{points} de l'espace, via les vecteurs.

\begin{definition}[Repère de l'espace et coordonnées d'un point]
Un \cle{repère} de l'espace est un quadruplet $(O\,;\,\vec{u}, \vec{v}, \vec{w})$ où $O$ est un point de l'espace (l'\emph{origine}) et $(\vec{u}, \vec{v}, \vec{w})$ une base de $\mathbb{R}^3$. Pour tout point $M$ de l'espace, les \cle{coordonnées} de $M$ dans ce repère sont les coordonnées $(\alpha\,;\,\beta\,;\,\gamma)$ du vecteur $\overrightarrow{OM}$ dans la base $(\vec{u}, \vec{v}, \vec{w})$ :
\[ \overrightarrow{OM} = \alpha\,\vec{u} + \beta\,\vec{v} + \gamma\,\vec{w} \iff M\,(\alpha\,;\,\beta\,;\,\gamma) \text{ dans } (O\,;\,\vec{u}, \vec{v}, \vec{w}). \]
\end{definition}

Le repère \emph{canonique} $(O\,;\,\vec{i}, \vec{j}, \vec{k})$ redonne les coordonnées cartésiennes usuelles. Mais rien n'interdit de choisir une origine et des vecteurs de base adaptés à une situation : par exemple, pour décrire les positions des piliers d'un pont, on placera l'origine au pied du premier pilier.

\begin{exoresolu}[Coordonnées d'un point dans un repère]
L'espace est rapporté au repère canonique $(O\,;\,\vec{i}, \vec{j}, \vec{k})$. On considère les points $A(1\,;\,0\,;\,1)$, $B(2\,;\,1\,;\,1)$, $C(1\,;\,2\,;\,2)$ et le point $M(4\,;\,3\,;\,5)$. On admet que $(\overrightarrow{AB}, \overrightarrow{AC}, \vec{k})$ forme une base de $\mathbb{R}^3$. Déterminer les coordonnées du point $M$ dans le repère $(A\,;\,\overrightarrow{AB}, \overrightarrow{AC}, \vec{k})$.
\tcblower
\textbf{Étape 1 : calcul des vecteurs de la base.}
\[ \overrightarrow{AB} = (2-1\,;\,1-0\,;\,1-1) = (1\,;\,1\,;\,0), \qquad \overrightarrow{AC} = (1-1\,;\,2-0\,;\,2-1) = (0\,;\,2\,;\,1). \]
\textbf{Étape 2 : vecteur à décomposer.}
\[ \overrightarrow{AM} = (4-1\,;\,3-0\,;\,5-1) = (3\,;\,3\,;\,4). \]
\textbf{Étape 3 : on pose} $\overrightarrow{AM} = \alpha\,\overrightarrow{AB} + \beta\,\overrightarrow{AC} + \gamma\,\vec{k}$, ce qui donne composante par composante :
\[
\begin{cases}
\alpha = 3 & (L_1)\\
\alpha + 2\beta = 3 & (L_2)\\
\beta + \gamma = 4 & (L_3)
\end{cases}
\]
De $(L_1)$ : $\alpha = 3$. De $(L_2)$ : $2\beta = 3 - 3 = 0$, donc $\beta = 0$. De $(L_3)$ : $\gamma = 4 - 0 = 4$.

\textbf{Étape 4 : vérification.} $3\overrightarrow{AB} + 0\,\overrightarrow{AC} + 4\vec{k} = (3\,;\,3\,;\,0) + (0\,;\,0\,;\,4) = (3\,;\,3\,;\,4) = \overrightarrow{AM}$. \checkmark

\textbf{Conclusion.} Les coordonnées de $M$ dans le repère $(A\,;\,\overrightarrow{AB}, \overrightarrow{AC}, \vec{k})$ sont $(3\,;\,0\,;\,4)$. Remarque : dans le repère canonique, les coordonnées de $M$ étaient $(4\,;\,3\,;\,5)$ ; le changement d'origine et de base a entièrement modifié le triplet, alors que le point $M$ n'a pas bougé.
\end{exoresolu}

\begin{exercice}[À toi de jouer]
On considère la base $\mathcal{B} = (\vec{u}, \vec{v}, \vec{w})$ de $\mathbb{R}^3$ avec $\vec{u} = (1\,;\,0\,;\,1)$, $\vec{v} = (1\,;\,1\,;\,0)$, $\vec{w} = (0\,;\,1\,;\,1)$.
\begin{enumerate}
\item Déterminer les coordonnées du vecteur $\vec{x} = (4\,;\,5\,;\,3)$ dans la base $\mathcal{B}$.
\item Le vecteur $\vec{y}$ a pour coordonnées $(1\,;\,2\,;\,-1)$ dans la base $\mathcal{B}$. Quelles sont ses coordonnées dans la base canonique ?
\end{enumerate}
\end{exercice}

\begin{corrige}[Exercice 1]
\textbf{1.} On pose $\vec{x} = \alpha\vec{u} + \beta\vec{v} + \gamma\vec{w}$, d'où le système :
\[
\begin{cases}
\alpha + \beta = 4 & (L_1)\\
\beta + \gamma = 5 & (L_2)\\
\alpha + \gamma = 3 & (L_3)
\end{cases}
\]
$(L_1) + (L_3) - (L_2)$ donne $2\alpha = 4 + 3 - 5 = 2$, donc $\alpha = 1$. Alors $\beta = 4 - 1 = 3$ et $\gamma = 3 - 1 = 2$. Vérification : $\vec{u} + 3\vec{v} + 2\vec{w} = (1\,;\,0\,;\,1) + (3\,;\,3\,;\,0) + (0\,;\,2\,;\,2) = (4\,;\,5\,;\,3) = \vec{x}$. \checkmark{} Les coordonnées de $\vec{x}$ dans $\mathcal{B}$ sont $(1\,;\,3\,;\,2)$.

\textbf{2.} $\vec{y} = 1\vec{u} + 2\vec{v} - 1\vec{w} = (1\,;\,0\,;\,1) + (2\,;\,2\,;\,0) - (0\,;\,1\,;\,1) = (3\,;\,1\,;\,0)$. Dans la base canonique, $\vec{y}$ a pour coordonnées $(3\,;\,1\,;\,0)$.
\end{corrige}

\begin{aretenir}[Bilan de la section]
\begin{itemize}
\item Si $\mathcal{B} = (\vec{u}, \vec{v}, \vec{w})$ est une base de $\mathbb{R}^3$, tout vecteur $\vec{x}$ s'écrit de façon \textbf{unique} $\vec{x} = \alpha\vec{u} + \beta\vec{v} + \gamma\vec{w}$ : $(\alpha\,;\,\beta\,;\,\gamma)$ sont ses coordonnées dans $\mathcal{B}$.
\item L'\emph{existence} vient du caractère générateur, l'\emph{unicité} de la liberté.
\item Trouver des coordonnées = résoudre un système linéaire $3 \times 3$, puis \emph{vérifier}.
\item Dans la base canonique, coordonnées et composantes coïncident ; dans toute autre base, il faut calculer.
\item Les coordonnées d'un point $M$ dans le repère $(O\,;\,\vec{u},\vec{v},\vec{w})$ sont celles de $\overrightarrow{OM}$ dans la base $(\vec{u},\vec{v},\vec{w})$.
\end{itemize}
\end{aretenir}

Cette capacité à associer à tout vecteur un triplet de nombres — et réciproquement — est la clé qui va nous permettre, dans la prochaine section, d'étudier les \emph{applications linéaires} : des transformations de l'espace entièrement décrites, elles aussi, par des tableaux de nombres appelés \emph{matrices}.

\coursec{Applications linéaires}

Dans la section précédente, nous avons appris à \emph{repérer} un vecteur de l'espace grâce à ses coordonnées dans une base. Nous allons maintenant étudier comment l'espace peut être \emph{transformé} : projeter l'ombre d'un objet sur le sol, agrandir une maquette, faire pivoter une pièce mécanique autour d'un axe. Toutes ces transformations partagent une propriété remarquable : elles \textbf{respectent la structure d'espace vectoriel}, c'est-à-dire les additions et les multiplications par un scalaire. Ce sont les \cle{applications linéaires}, véritables « machines à transformer les vecteurs » qui seront au cœur de la suite du chapitre (matrices, noyau, image).

\courssub{Définition et premières propriétés}

\textbf{Une intuition pour démarrer.} Imagine une imprimante 3D à Douala qui agrandit un prototype : si tu doubles la taille de deux pièces séparément, puis que tu les assembles, tu obtiens le même résultat que si tu assembles d'abord puis doubles l'ensemble. L'agrandissement « commute » avec l'assemblage. C'est exactement l'esprit de la linéarité : \emph{l'image d'une combinaison est la combinaison des images}.

\begin{definition}[Application linéaire]
Soient $E$ et $F$ deux espaces vectoriels sur $\mathbb{R}$. Une application $f : E \to F$ est dite \cle{linéaire} si elle vérifie les deux conditions suivantes :
\begin{enumerate}
\item \textbf{Additivité :} pour tous vecteurs $\vec{u}, \vec{v}$ de $E$, $f(\vec{u} + \vec{v}) = f(\vec{u}) + f(\vec{v})$ ;
\item \textbf{Homogénéité :} pour tout vecteur $\vec{u}$ de $E$ et tout réel $\lambda$, $f(\lambda \vec{u}) = \lambda f(\vec{u})$.
\end{enumerate}
Lorsque $E = F$, on dit que $f$ est un \cle{endomorphisme} de $E$. Lorsque $F = \mathbb{R}$, on dit que $f$ est une \cle{forme linéaire}.
\end{definition}

Ces deux conditions peuvent être fusionnées en une seule, souvent plus pratique dans les démonstrations.

\begin{propriete}[Caractérisation équivalente]
Une application $f : E \to F$ est linéaire si et seulement si, pour tous vecteurs $\vec{u}, \vec{v}$ de $E$ et tous réels $\lambda, \mu$ :
\[ f(\lambda \vec{u} + \mu \vec{v}) = \lambda f(\vec{u}) + \mu f(\vec{v}). \]
Autrement dit : \emph{l'image d'une combinaison linéaire est la combinaison linéaire des images, avec les mêmes coefficients}.
\end{propriete}

\begin{experience}[Pourquoi les deux formulations sont équivalentes]
\textbf{Sens direct.} Supposons $f$ linéaire (les deux conditions de la définition). Alors, en appliquant d'abord l'additivité puis l'homogénéité :
\[ f(\lambda \vec{u} + \mu \vec{v}) = f(\lambda \vec{u}) + f(\mu \vec{v}) = \lambda f(\vec{u}) + \mu f(\vec{v}). \]
\textbf{Sens réciproque.} Supposons la condition combinée vraie pour tous $\lambda, \mu$. En prenant $\lambda = \mu = 1$, on obtient l'additivité ; en prenant $\mu = 0$ (et $\vec{v}$ quelconque), on obtient $f(\lambda \vec{u}) = \lambda f(\vec{u})$, c'est-à-dire l'homogénéité. Les deux formulations disent donc exactement la même chose.
\end{experience}

\begin{propriete}[Conséquences immédiates]
Si $f : E \to F$ est linéaire, alors :
\[ f(\vec{0}_E) = \vec{0}_F \qquad \text{et} \qquad f(-\vec{u}) = -f(\vec{u}) \ \text{pour tout } \vec{u} \in E. \]
\end{propriete}

\begin{experience}[Démonstration (exigible au Bac)]
\textbf{Image du vecteur nul.} Puisque $0 \cdot \vec{u} = \vec{0}_E$, l'homogénéité donne :
\[ f(\vec{0}_E) = f(0 \cdot \vec{u}) = 0 \cdot f(\vec{u}) = \vec{0}_F. \]
\textbf{Image de l'opposé.} En prenant $\lambda = -1$ dans l'homogénéité :
\[ f(-\vec{u}) = f((-1)\vec{u}) = (-1) f(\vec{u}) = -f(\vec{u}). \]
Une application linéaire envoie donc toujours l'origine sur l'origine, et l'opposé d'un vecteur sur l'opposé de son image.
\end{experience}

\begin{aretenir}[Le test rapide du vecteur nul]
Si $f(\vec{0}_E) \neq \vec{0}_F$, alors $f$ \textbf{n'est pas linéaire}. C'est le moyen le plus rapide de détecter une application non linéaire. Attention à la réciproque : $f(\vec{0}) = \vec{0}$ ne suffit \emph{pas} à garantir la linéarité (par exemple $f(x,y,z) = (x^2, y, z)$ envoie $\vec{0}$ sur $\vec{0}$ sans être linéaire).
\end{aretenir}

\courssub{Exemples fondamentaux dans l'espace}

\textbf{Exemples géométriques (endomorphismes de $\mathbb{R}^3$).} Les transformations que tu as rencontrées en géométrie de l'espace sont, pour la plupart, linéaires dès lors qu'elles fixent l'origine :

\begin{itemize}
\item \textbf{Homothétie de rapport $k$} (centrée en $O$) : $h(\vec{u}) = k\vec{u}$, soit en coordonnées $h(x, y, z) = (kx, ky, kz)$. Elle est linéaire car $k(\vec{u}+\vec{v}) = k\vec{u} + k\vec{v}$ et $k(\lambda \vec{u}) = \lambda(k\vec{u})$.
\item \textbf{Symétrie par rapport à un plan passant par $O$} : par exemple la symétrie par rapport au plan $(xOy)$, $s(x, y, z) = (x, y, -z)$.
\item \textbf{Projection orthogonale sur un plan passant par $O$} : par exemple la projection sur le plan $(xOy)$, $p(x, y, z) = (x, y, 0)$. Vérifions l'homogénéité : $p(\lambda x, \lambda y, \lambda z) = (\lambda x, \lambda y, 0) = \lambda\, p(x,y,z)$. L'additivité se vérifie de même.
\item \textbf{Rotation d'angle $\theta$ autour d'un axe passant par $O$} : par exemple autour de l'axe $(Oz)$, $r(x, y, z) = (x\cos\theta - y\sin\theta,\; x\sin\theta + y\cos\theta,\; z)$.
\end{itemize}

\begin{popfigure}
\begin{center}
\begin{tikzpicture}[scale=1.1, >=Stealth]
% Plan (xOy) en perspective
\fill[popblueL, opacity=0.55] (-2.6,-1.3) -- (2.2,-1.3) -- (3.4,0.5) -- (-1.4,0.5) -- cycle;
\draw[popblue] (-2.6,-1.3) -- (2.2,-1.3) -- (3.4,0.5) -- (-1.4,0.5) -- cycle;
\node[popblue] at (2.9,-0.9) {\small plan $(xOy)$};
% Origine
\coordinate (O) at (0,-0.4);
\fill (O) circle (1.5pt) node[below left] {\small $O$};
% Vecteur u et sa projection
\coordinate (U) at (1.6,2.2);
\coordinate (PU) at (1.6,-0.05);
\draw[->, very thick, popink] (O) -- (U) node[above] {\small $\vec{u}$};
\draw[->, very thick, poporange] (O) -- (PU) node[below right] {\small $p(\vec{u})$};
\draw[dashed, popmuted] (U) -- (PU);
% Vecteur v et sa projection
\coordinate (V) at (-1.5,1.5);
\coordinate (PV) at (-1.5,0.15);
\draw[->, very thick, poppurple] (O) -- (V) node[above left] {\small $\vec{v}$};
\draw[->, very thick, popgreen] (O) -- (PV) node[left] {\small $p(\vec{v})$};
\draw[dashed, popmuted] (V) -- (PV);
\end{tikzpicture}
\end{center}
\legende{Projection orthogonale $p$ de $\mathbb{R}^3$ sur le plan $(xOy)$ : chaque vecteur est « écrasé » verticalement sur le plan. On observe que $p(\vec{u} + \vec{v}) = p(\vec{u}) + p(\vec{v})$ : la projection est linéaire.}
\end{popfigure}

\begin{popfigure}
\begin{center}
\begin{tikzpicture}[scale=0.85, >=Stealth]
% Petit cube
\coordinate (A) at (0,0); \coordinate (B) at (1,0); \coordinate (C) at (1.35,0.45); \coordinate (D) at (0.35,0.45);
\coordinate (A') at (0,1); \coordinate (B') at (1,1); \coordinate (C') at (1.35,1.45); \coordinate (D') at (0.35,1.45);
\draw[thick, popblue] (A)--(B)--(C)--(D)--cycle (A')--(B')--(C')--(D')--cycle (A)--(A') (B)--(B') (C)--(C') (D)--(D');
% Grand cube (rapport 2)
\begin{scope}[shift={(3.2,0)}]
\coordinate (A2) at (0,0); \coordinate (B2) at (2,0); \coordinate (C2) at (2.7,0.9); \coordinate (D2) at (0.7,0.9);
\coordinate (A2') at (0,2); \coordinate (B2') at (2,2); \coordinate (C2') at (2.7,2.9); \coordinate (D2') at (0.7,2.9);
\draw[thick, poporange] (A2)--(B2)--(C2)--(D2)--cycle (A2')--(B2')--(C2')--(D2')--cycle (A2)--(A2') (B2)--(B2') (C2)--(C2') (D2)--(D2');
\end{scope}
% Flèche de transformation
\draw[->, very thick, popink] (1.9,1.2) -- (2.9,1.2) node[midway, above] {\small $h$, rapport $2$};
\node[popblue] at (0.7,-0.5) {\small cube initial};
\node[poporange] at (4.5,-0.5) {\small cube image};
\end{tikzpicture}
\end{center}
\legende{Homothétie $h$ de rapport $2$ appliquée à un cube : toutes les arêtes sont doublées. Chaque sommet $\vec{u}$ devient $2\vec{u}$ : c'est un endomorphisme de $\mathbb{R}^3$.}
\end{popfigure}

\textbf{Formes linéaires (applications de $\mathbb{R}^3$ dans $\mathbb{R}$).} Une forme linéaire associe à un vecteur un \emph{nombre}. Par exemple :
\[ f(x, y, z) = 2x - y + 3z. \]
Vérifions l'additivité avec $\vec{u} = (x, y, z)$ et $\vec{v} = (x', y', z')$ :
\begin{align*}
f(\vec{u} + \vec{v}) &= 2(x + x') - (y + y') + 3(z + z') \\
&= (2x - y + 3z) + (2x' - y' + 3z') = f(\vec{u}) + f(\vec{v}),
\end{align*}
et l'homogénéité : $f(\lambda x, \lambda y, \lambda z) = 2\lambda x - \lambda y + 3\lambda z = \lambda f(x, y, z)$. Donc $f$ est linéaire.

\begin{savaistu}[Les formes linéaires dans ta poche]
Quand ton téléphone calcule le coût d'une consommation mixte — $x$ minutes d'appels MTN, $y$ mégaoctets d'Internet, $z$ SMS — selon une formule du type $C = a x + b y + c z$, il applique une \emph{forme linéaire} ! En revanche, dès qu'il existe un « forfait fixe » (un terme constant), l'application devient affine et cesse d'être linéaire : c'est exactement le contre-exemple que nous allons voir.
\end{savaistu}

\textbf{Contre-exemples : ce qui n'est PAS linéaire.}
\begin{itemize}
\item \textbf{Les translations} : $t(\vec{u}) = \vec{u} + \vec{a}$ avec $\vec{a} \neq \vec{0}$. En effet $t(\vec{0}) = \vec{a} \neq \vec{0}$, donc $t$ n'est pas linéaire.
\item \textbf{Les applications affines avec terme constant} : $f(x, y, z) = (x + 2, y - 1, z)$, ou plus généralement $f(\vec{u}) = A\vec{u} + \vec{b}$ avec $\vec{b} \neq \vec{0}$. Là encore $f(\vec{0}) = \vec{b} \neq \vec{0}$.
\item \textbf{Les applications avec des puissances ou des produits de coordonnées} : $f(x, y, z) = (x^2, y, z)$ ou $g(x, y, z) = (xy, z, 0)$ ne sont pas linéaires (par exemple $f(2,0,0) = (4,0,0) \neq 2 f(1,0,0) = (2,0,0)$).
\end{itemize}

\begin{attention}[Le piège classique de l'application affine]
L'erreur la plus fréquente au Bac est de croire qu'une application du type $f(x, y, z) = (x + 1,\; y,\; z)$ est linéaire « parce qu'il n'y a que des $x$, $y$, $z$ à la puissance 1 ». \textbf{Faux !} Le terme constant $+1$ détruit la linéarité. Réflexe à acquérir : \emph{toute application linéaire s'écrit avec des coordonnées qui sont des combinaisons linéaires de $x, y, z$ \textbf{sans terme constant}}. Et en cas de doute, teste $f(\vec{0})$ : si ce n'est pas $\vec{0}$, c'est terminé.
\end{attention}

\courssub{Méthodes : prouver ou réfuter la linéarité}

\begin{methode}[Montrer qu'une application EST linéaire]
Pour prouver que $f : \mathbb{R}^3 \to \mathbb{R}^p$ est linéaire :
\begin{enumerate}
\item Prendre deux vecteurs \textbf{génériques} $\vec{u} = (x, y, z)$ et $\vec{v} = (x', y', z')$, et deux réels \textbf{génériques} $\lambda, \mu$.
\item Calculer $f(\lambda \vec{u} + \mu \vec{v})$ en développant complètement.
\item Calculer séparément $\lambda f(\vec{u}) + \mu f(\vec{v})$.
\item Conclure : si les deux expressions sont égales \emph{pour tous} $x, y, z, x', y', z', \lambda, \mu$, alors $f$ est linéaire.
\end{enumerate}
Le point essentiel est de travailler avec des \textbf{lettres} (vecteurs génériques), jamais avec des valeurs numériques particulières : un exemple numérique ne prouve rien.
\end{methode}

\begin{methode}[Montrer qu'une application N'EST PAS linéaire]
Trois stratégies, de la plus rapide à la plus générale :
\begin{enumerate}
\item \textbf{Test du vecteur nul} : calculer $f(\vec{0})$. Si $f(\vec{0}) \neq \vec{0}$, c'est fini, $f$ n'est pas linéaire.
\item \textbf{Contre-exemple à l'homogénéité} : trouver un vecteur $\vec{u}$ et un réel $\lambda$ tels que $f(\lambda \vec{u}) \neq \lambda f(\vec{u})$ (souvent $\lambda = 2$ suffit).
\item \textbf{Contre-exemple à l'additivité} : trouver $\vec{u}, \vec{v}$ tels que $f(\vec{u} + \vec{v}) \neq f(\vec{u}) + f(\vec{v})$.
\end{enumerate}
Un \textbf{seul} contre-exemple suffit à détruire la linéarité.
\end{methode}

\begin{exemplebox}[Une application linéaire et une qui ne l'est pas]
\textbf{1.} Soit $f : \mathbb{R}^3 \to \mathbb{R}^3$ définie par $f(x, y, z) = (x + y,\; y - z,\; 2x + z)$. Montrons que $f$ est linéaire.

Soient $\vec{u} = (x, y, z)$, $\vec{v} = (x', y', z')$ et $\lambda, \mu \in \mathbb{R}$. Alors $\lambda \vec{u} + \mu \vec{v} = (\lambda x + \mu x',\; \lambda y + \mu y',\; \lambda z + \mu z')$, donc :
\begin{align*}
f(\lambda \vec{u} + \mu \vec{v}) &= \big( (\lambda x + \mu x') + (\lambda y + \mu y'),\; (\lambda y + \mu y') - (\lambda z + \mu z'),\; 2(\lambda x + \mu x') + (\lambda z + \mu z') \big) \\
&= \lambda (x + y,\; y - z,\; 2x + z) + \mu (x' + y',\; y' - z',\; 2x' + z') \\
&= \lambda f(\vec{u}) + \mu f(\vec{v}).
\end{align*}
L'égalité vaut pour tous vecteurs et tous scalaires : $f$ est \textbf{linéaire}.

\textbf{2.} Soit $g : \mathbb{R}^3 \to \mathbb{R}^3$ définie par $g(x, y, z) = (x + 1,\; y,\; z)$. On calcule :
\[ g(0, 0, 0) = (1, 0, 0) \neq (0, 0, 0). \]
Donc $g$ \textbf{n'est pas linéaire}. (C'est une translation de vecteur $(1, 0, 0)$.)
\end{exemplebox}

\courssub{Détermination par les images d'une base}

Voici la propriété qui fait tout l'intérêt pratique des applications linéaires : il est \textbf{inutile de connaître l'image de chacun des vecteurs} (il y en a une infinité !). Il suffit de connaître les images des vecteurs d'une base.

\begin{propriete}[Une application linéaire est déterminée par les images d'une base]
Soit $(\vec{e}_1, \vec{e}_2, \vec{e}_3)$ une base de $\mathbb{R}^3$ et $f : \mathbb{R}^3 \to F$ une application linéaire. Alors $f$ est \textbf{entièrement déterminée} par les trois vecteurs $f(\vec{e}_1)$, $f(\vec{e}_2)$, $f(\vec{e}_3)$ : pour tout $\vec{u} = x\vec{e}_1 + y\vec{e}_2 + z\vec{e}_3$,
\[ f(\vec{u}) = x\, f(\vec{e}_1) + y\, f(\vec{e}_2) + z\, f(\vec{e}_3). \]
Réciproquement, si l'on se donne \emph{arbitrairement} trois vecteurs $\vec{a}_1, \vec{a}_2, \vec{a}_3$ de $F$, il existe une \textbf{unique} application linéaire $f$ telle que $f(\vec{e}_i) = \vec{a}_i$ pour $i = 1, 2, 3$.
\end{propriete}

\begin{experience}[Démonstration (formatrice)]
Soit $\vec{u} \in \mathbb{R}^3$. Puisque $(\vec{e}_1, \vec{e}_2, \vec{e}_3)$ est une base, $\vec{u}$ s'écrit de manière \emph{unique} $\vec{u} = x\vec{e}_1 + y\vec{e}_2 + z\vec{e}_3$ (c'est la section précédente !). Par linéarité de $f$ :
\[ f(\vec{u}) = f(x\vec{e}_1 + y\vec{e}_2 + z\vec{e}_3) = x f(\vec{e}_1) + y f(\vec{e}_2) + z f(\vec{e}_3). \]
Le résultat ne dépend donc que des coordonnées $(x, y, z)$ de $\vec{u}$ et des trois images $f(\vec{e}_i)$. D'où l'unicité. Pour l'existence, on \emph{définit} $f$ par la formule ci-dessus et on vérifie qu'elle est bien linéaire (le calcul est identique à celui de l'exemple précédent).
\end{experience}

C'est précisément cette propriété qui rendra possible, dans la prochaine section, le codage d'une application linéaire par un tableau de nombres : sa \cle{matrice}.

\begin{exoresolu}[Retrouver une application linéaire à partir des images d'une base]
Soit $f$ un endomorphisme de $\mathbb{R}^3$ tel que, pour la base canonique $(\vec{i}, \vec{j}, \vec{k})$ :
\[ f(\vec{i}) = (1, 0, 2), \qquad f(\vec{j}) = (1, 1, 0), \qquad f(\vec{k}) = (0, -1, 1). \]
Déterminer $f(x, y, z)$ pour un vecteur $(x, y, z)$ quelconque, puis calculer $f(2, 1, -1)$.
\tcblower
\textbf{Solution.} Le vecteur $(x, y, z)$ s'écrit $x\vec{i} + y\vec{j} + z\vec{k}$. Par linéarité :
\begin{align*}
f(x, y, z) &= x f(\vec{i}) + y f(\vec{j}) + z f(\vec{k}) \\
&= x(1, 0, 2) + y(1, 1, 0) + z(0, -1, 1) \\
&= (x + y,\; y - z,\; 2x + z).
\end{align*}
On retrouve l'application de l'exemple précédent ! Pour $(2, 1, -1)$ :
\[ f(2, 1, -1) = (2 + 1,\; 1 - (-1),\; 2 \times 2 + (-1)) = (3, 2, 3). \]
Vérification directe : $2 f(\vec{i}) + f(\vec{j}) - f(\vec{k}) = (2,0,4) + (1,1,0) - (0,-1,1) = (3, 2, 3)$. \checkmark
\end{exoresolu}

\courssub{Composition d'applications linéaires}

\begin{propriete}[La composée de deux applications linéaires est linéaire]
Soient $f : E \to F$ et $g : F \to G$ deux applications linéaires. Alors l'application composée $g \circ f : E \to G$, définie par $(g \circ f)(\vec{u}) = g\big(f(\vec{u})\big)$, est encore linéaire.
\end{propriete}

\begin{experience}[Démonstration (courte et exigible)]
Soient $\vec{u}, \vec{v} \in E$ et $\lambda, \mu \in \mathbb{R}$. On applique successivement la linéarité de $f$ puis celle de $g$ :
\begin{align*}
(g \circ f)(\lambda \vec{u} + \mu \vec{v}) &= g\big( f(\lambda \vec{u} + \mu \vec{v}) \big) \\
&= g\big( \lambda f(\vec{u}) + \mu f(\vec{v}) \big) \quad \text{(linéarité de } f\text{)} \\
&= \lambda\, g\big(f(\vec{u})\big) + \mu\, g\big(f(\vec{v})\big) \quad \text{(linéarité de } g\text{)} \\
&= \lambda\, (g \circ f)(\vec{u}) + \mu\, (g \circ f)(\vec{v}).
\end{align*}
Donc $g \circ f$ est linéaire.
\end{experience}

\begin{exemplebox}[Composer une rotation et une homothétie]
Soit $h$ l'homothétie de rapport $3$ et $r$ la rotation d'angle $\frac{\pi}{2}$ autour de l'axe $(Oz)$, définie par $r(x, y, z) = (-y, x, z)$. Alors :
\[ (r \circ h)(x, y, z) = r(3x, 3y, 3z) = (-3y, 3x, 3z), \]
qui est bien linéaire (chaque coordonnée est une combinaison linéaire de $x, y, z$ sans terme constant). Géométriquement : agrandir puis tourner revient à une seule transformation linéaire.
\end{exemplebox}

\begin{attention}[L'ordre de composition compte]
Contrairement à l'addition des vecteurs, la composition n'est \textbf{pas commutative} : en général $g \circ f \neq f \circ g$. Par exemple, avec la rotation $r$ ci-dessus et la symétrie $s(x, y, z) = (x, -y, z)$ : $(r \circ s)(1, 0, 0) = r(1, 0, 0) = (0, 1, 0)$, mais $(s \circ r)(1, 0, 0) = s(0, 1, 0) = (0, -1, 0)$. Rédige donc les composées dans le bon ordre : $g \circ f$ signifie « on applique $f$ \emph{d'abord}, puis $g$ ».
\end{attention}

\begin{exercice}[À toi de jouer]
On considère les applications de $\mathbb{R}^3$ dans $\mathbb{R}^3$ définies par :
\[ f(x, y, z) = (x - 2y + z,\; 3x + z,\; y), \qquad g(x, y, z) = (x,\; y,\; z^2), \qquad h(x, y, z) = (x,\; y + 5,\; z). \]
\begin{enumerate}
\item Parmi $f$, $g$ et $h$, lesquelles sont linéaires ? Justifier soigneusement.
\item Pour celle(s) qui le sont, déterminer les images des vecteurs $\vec{i}$, $\vec{j}$, $\vec{k}$ de la base canonique.
\end{enumerate}
\end{exercice}

\begin{corrige}[Exercice 1]
\textbf{1.} \textbf{Application $f$ :} soient $\vec{u} = (x, y, z)$, $\vec{v} = (x', y', z')$ et $\lambda, \mu \in \mathbb{R}$.
\begin{align*}
f(\lambda \vec{u} + \mu \vec{v}) &= \big( (\lambda x + \mu x') - 2(\lambda y + \mu y') + (\lambda z + \mu z'),\; 3(\lambda x + \mu x') + (\lambda z + \mu z'),\; \lambda y + \mu y' \big) \\
&= \lambda (x - 2y + z,\; 3x + z,\; y) + \mu (x' - 2y' + z',\; 3x' + z',\; y') = \lambda f(\vec{u}) + \mu f(\vec{v}).
\end{align*}
Donc $f$ est \textbf{linéaire}.

\textbf{Application $g$ :} testons l'homogénéité avec $\vec{u} = (0, 0, 1)$ et $\lambda = 2$ : $g(2\vec{u}) = g(0, 0, 2) = (0, 0, 4)$, tandis que $2 g(\vec{u}) = 2(0, 0, 1) = (0, 0, 2)$. Comme $(0,0,4) \neq (0,0,2)$, $g$ \textbf{n'est pas linéaire} (à cause du carré $z^2$).

\textbf{Application $h$ :} $h(0, 0, 0) = (0, 5, 0) \neq (0, 0, 0)$, donc $h$ \textbf{n'est pas linéaire} (terme constant $+5$ : c'est une application affine).

\textbf{2.} Pour $f$ : $f(\vec{i}) = f(1, 0, 0) = (1, 3, 0)$ ; $f(\vec{j}) = f(0, 1, 0) = (-2, 0, 1)$ ; $f(\vec{k}) = f(0, 0, 1) = (1, 1, 0)$.
\end{corrige}

\begin{aretenir}[L'essentiel de la section]
\begin{itemize}
\item $f$ est \cle{linéaire} si $f(\lambda \vec{u} + \mu \vec{v}) = \lambda f(\vec{u}) + \mu f(\vec{v})$ pour tous vecteurs et tous scalaires.
\item Toute application linéaire vérifie $f(\vec{0}) = \vec{0}$ : c'est le test de réfutation le plus rapide.
\item Homothéties, symétries, projections, rotations (centrées en/passant par $O$) et formes linéaires sont linéaires ; translations et applications affines avec terme constant ne le sont \textbf{jamais}.
\item Une application linéaire est entièrement déterminée par les images des vecteurs d'une base — c'est la clé de la notion de matrice, objet de la prochaine section.
\item La composée de deux applications linéaires est linéaire, mais l'ordre de composition compte.
\end{itemize}
\end{aretenir}

\coursec{Matrice d'une application linéaire}

Dans la section précédente, nous avons découvert les \cle{applications linéaires} : ces fonctions entre espaces vectoriels qui respectent l'addition et la multiplication par un scalaire, c'est-à-dire $f(u+v) = f(u) + f(v)$ et $f(\lambda u) = \lambda f(u)$. Nous avons vu qu'une application linéaire de $\mathbb{R}^3$ dans $\mathbb{R}^3$ s'écrit toujours sous la forme
\[ f(x,y,z) = (ax + by + cz,\; a'x + b'y + c'z,\; a''x + b''y + c''z), \]
où les neuf coefficients $a, b, c, a', \ldots$ sont des réels. Cette observation est capitale : une application linéaire est \emph{entièrement déterminée par un nombre fini de coefficients}. Toute la question est maintenant de les \textbf{organiser} de manière efficace. C'est exactement le rôle de la \cle{matrice} d'une application linéaire : un tableau de nombres qui code complètement la transformation, et qui permet de calculer les images par une simple multiplication. C'est ainsi, par exemple, qu'un logiciel de conception 3D calcule la position d'une pièce mécanique après rotation, ou qu'une application de cartographie transforme les coordonnées d'un point de Douala d'un repère à un autre.

\courssub{Principe fondamental : une application linéaire est déterminée par les images d'une base}

Commençons par l'intuition. Soit $f : E \to F$ une application linéaire, et $\mathcal{B} = (u_1, u_2, u_3)$ une base de $E$. Tout vecteur $u$ de $E$ se décompose de manière \emph{unique} :
\[ u = x\,u_1 + y\,u_2 + z\,u_3. \]
Par linéarité de $f$, on obtient :
\[ f(u) = x\,f(u_1) + y\,f(u_2) + z\,f(u_3). \]
Regarde bien cette formule : pour calculer $f(u)$ pour \emph{n'importe quel} vecteur $u$, il suffit de connaître les trois vecteurs $f(u_1)$, $f(u_2)$, $f(u_3)$. Toute l'information sur $f$ est contenue dans ces trois images.

\begin{aretenir}[Principe fondamental]
Une application linéaire $f : E \to F$ est \cle{entièrement déterminée} par les images des vecteurs d'une base de $E$. Si chacune de ces images est exprimée par ses coordonnées dans une base $\mathcal{B}'$ de $F$, alors $f$ est codée par un tableau de nombres : sa \textbf{matrice}.
\end{aretenir}

\begin{experience}[Activité de découverte : coder une transformation]
Considérons l'application linéaire $f : \mathbb{R}^3 \to \mathbb{R}^3$ définie par $f(x,y,z) = (x+y,\; y-z,\; 2x+z)$, et la base canonique $\mathcal{B} = (e_1, e_2, e_3)$ avec $e_1 = (1,0,0)$, $e_2 = (0,1,0)$, $e_3 = (0,0,1)$.
\begin{enumerate}
\item Calcule $f(e_1)$ : on remplace $(x,y,z)$ par $(1,0,0)$ : $f(e_1) = (1+0,\; 0-0,\; 2\times 1 + 0) = (1, 0, 2)$.
\item Calcule de même $f(e_2) = (0+1,\; 1-0,\; 0+0) = (1, 1, 0)$.
\item Calcule $f(e_3) = (0+0,\; 0-1,\; 0+1) = (0, -1, 1)$.
\item Range ces trois résultats \emph{en colonnes}, côte à côte, dans un tableau à 3 lignes et 3 colonnes.
\end{enumerate}
Tu viens de construire la matrice de $f$ ! Les sections suivantes formalisent cette démarche.
\end{experience}

\courssub{Définition de la matrice d'une application linéaire}

\begin{definition}[Matrice d'une application linéaire dans des bases données]
Soit $f : E \to F$ une application linéaire, $\mathcal{B} = (u_1, u_2, u_3)$ une base de $E$ (espace de \cle{départ}) et $\mathcal{B}' = (v_1, v_2, v_3)$ une base de $F$ (espace d'\cle{arrivée}). La \textbf{matrice de $f$ relativement aux bases $\mathcal{B}$ et $\mathcal{B}'$}, notée $\mathrm{Mat}_{\mathcal{B}, \mathcal{B}'}(f)$, est le tableau à 3 lignes et 3 colonnes dont :
\begin{itemize}
\item la \textbf{première colonne} contient les coordonnées de $f(u_1)$ dans la base $\mathcal{B}'$ ;
\item la \textbf{deuxième colonne} contient les coordonnées de $f(u_2)$ dans la base $\mathcal{B}'$ ;
\item la \textbf{troisième colonne} contient les coordonnées de $f(u_3)$ dans la base $\mathcal{B}'$.
\end{itemize}
Autrement dit, si $f(u_j) = a_{1j}\,v_1 + a_{2j}\,v_2 + a_{3j}\,v_3$ pour $j = 1, 2, 3$, alors
\[ \mathrm{Mat}_{\mathcal{B}, \mathcal{B}'}(f) = \begin{pmatrix} a_{11} & a_{12} & a_{13} \\ a_{21} & a_{22} & a_{23} \\ a_{31} & a_{32} & a_{33} \end{pmatrix}. \]
\end{definition}

Retiens la phrase-clé : \emph{les colonnes de la matrice sont les coordonnées des images des vecteurs de la base de départ, exprimées dans la base d'arrivée}.

\begin{center}
\begin{tabularx}{\linewidth}{|c|X|X|X|}
\hline
\rowcolor{popblueL} & \textbf{Colonne 1} & \textbf{Colonne 2} & \textbf{Colonne 3} \\
\hline
Contient & coordonnées de $f(u_1)$ dans $\mathcal{B}'$ & coordonnées de $f(u_2)$ dans $\mathcal{B}'$ & coordonnées de $f(u_3)$ dans $\mathcal{B}'$ \\
\hline
\rowcolor{popcream} Se lit & verticalement, de haut en bas & verticalement, de haut en bas & verticalement, de haut en bas \\
\hline
\end{tabularx}
\end{center}

Dans le cas fréquent où $E = F = \mathbb{R}^3$ et où l'on travaille dans la \textbf{même base} $\mathcal{B}$ au départ et à l'arrivée (on dit alors que $f$ est un \cle{endomorphisme}), on parle simplement de la \emph{matrice de $f$ dans la base $\mathcal{B}$}, notée $\mathrm{Mat}_{\mathcal{B}}(f)$. C'est une matrice carrée $3 \times 3$.

\begin{attention}[Le piège classique : colonnes, pas lignes !]
L'erreur la plus fréquente au Baccalauréat consiste à écrire les coordonnées des images $f(u_1), f(u_2), f(u_3)$ \textbf{en lignes} au lieu de les écrire \textbf{en colonnes}. La matrice obtenue est alors la \emph{transposée} de la bonne matrice, et tous les calculs qui suivent sont faux. Garde en mémoire le slogan : \textbf{« les images en colonnes »}. Une image mentale : chaque vecteur de base « tombe » verticalement dans sa colonne.
\end{attention}

\courssub{Méthode de construction et exemple complet}

\begin{methode}[Écrire la matrice d'une application linéaire dans la base canonique]
Soit $f : \mathbb{R}^3 \to \mathbb{R}^3$ une application linéaire donnée par ses formules. Pour obtenir sa matrice dans la base canonique $\mathcal{B} = (e_1, e_2, e_3)$ :
\begin{enumerate}
\item Calculer $f(e_1) = f(1,0,0)$, puis écrire ses coordonnées \textbf{en colonne} (colonne 1).
\item Calculer $f(e_2) = f(0,1,0)$, écrire ses coordonnées en colonne 2.
\item Calculer $f(e_3) = f(0,0,1)$, écrire ses coordonnées en colonne 3.
\item Vérifier : la matrice a 3 lignes et 3 colonnes, et la colonne $j$ redonne bien $f(e_j)$.
\end{enumerate}
\textbf{Astuce de lecture directe :} dans la base canonique, la colonne 1 contient les coefficients de $x$ dans les trois formules, la colonne 2 ceux de $y$, la colonne 3 ceux de $z$.
\end{methode}

\begin{exemplebox}[Construction pas à pas]
Soit $f : \mathbb{R}^3 \to \mathbb{R}^3$ définie par $f(x,y,z) = (x + y,\; y - z,\; 2x + z)$. Construisons sa matrice $A$ dans la base canonique.

\textbf{Étape 1.} $f(e_1) = f(1,0,0) = (1 + 0,\; 0 - 0,\; 2 \times 1 + 0) = (1, 0, 2)$.

\textbf{Étape 2.} $f(e_2) = f(0,1,0) = (0 + 1,\; 1 - 0,\; 0 + 0) = (1, 1, 0)$.

\textbf{Étape 3.} $f(e_3) = f(0,0,1) = (0 + 0,\; 0 - 1,\; 0 + 1) = (0, -1, 1)$.

\textbf{Étape 4.} On place ces trois vecteurs en colonnes :
\[ A = \begin{pmatrix} 1 & 1 & 0 \\ 0 & 1 & -1 \\ 2 & 0 & 1 \end{pmatrix}
\quad \begin{array}{l} \leftarrow \text{1\textsuperscript{re} coordonnée des images} \\ \leftarrow \text{2\textsuperscript{e} coordonnée des images} \\ \leftarrow \text{3\textsuperscript{e} coordonnée des images} \end{array} \]
Vérifions avec l'astuce : colonne 1 = coefficients de $x$ : $(1, 0, 2)$ ; colonne 2 = coefficients de $y$ : $(1, 1, 0)$ ; colonne 3 = coefficients de $z$ : $(0, -1, 1)$. C'est cohérent.
\end{exemplebox}

\begin{popfigure}
\begin{center}
\begin{tikzpicture}
\matrix[matrix of math nodes, left delimiter={(}, right delimiter={)}, row sep=6pt, column sep=10pt] (m) {
1 & 1 & 0 \\
0 & 1 & -1 \\
2 & 0 & 1 \\
};
\node[above=4pt of m-1-1, popblue] {\small $f(e_1)$};
\node[above=4pt of m-1-2, poporange] {\small $f(e_2)$};
\node[above=4pt of m-1-3, popgreen] {\small $f(e_3)$};
\draw[popblue, thick] (m-1-1.north west) rectangle (m-3-1.south east);
\draw[poporange, thick] (m-1-2.north west) rectangle (m-3-2.south east);
\draw[popgreen, thick] (m-1-3.north west) rectangle (m-3-3.south east);
\node[below=6pt of m-3-2, popdark] {\small Chaque colonne = coordonnées de l'image d'un vecteur de base};
\end{tikzpicture}
\end{center}
\legende{Matrice de $f(x,y,z) = (x+y,\, y-z,\, 2x+z)$ dans la base canonique : les trois colonnes sont les images de $e_1$, $e_2$, $e_3$.}
\end{popfigure}

\courssub{Le lien calculatoire : $Y = AX$}

Voici maintenant le résultat qui rend les matrices si puissantes : elles transforment le calcul d'une image en une simple multiplication.

\begin{propriete}[Calcul des coordonnées de l'image — admise]
Soit $f : E \to F$ une application linéaire, $A = \mathrm{Mat}_{\mathcal{B}, \mathcal{B}'}(f)$ sa matrice. Si un vecteur $u$ a pour colonne de coordonnées $X = \begin{pmatrix} x \\ y \\ z \end{pmatrix}$ dans la base $\mathcal{B}$, alors les coordonnées de $f(u)$ dans la base $\mathcal{B}'$ sont données par la colonne
\[ Y = AX, \]
où le produit se calcule ligne par ligne : la $i$-ème coordonnée de $Y$ est la somme des produits des coefficients de la $i$-ème ligne de $A$ par ceux de $X$ :
\[ \begin{pmatrix} a_{11} & a_{12} & a_{13} \\ a_{21} & a_{22} & a_{23} \\ a_{31} & a_{32} & a_{33} \end{pmatrix} \begin{pmatrix} x \\ y \\ z \end{pmatrix} = \begin{pmatrix} a_{11}x + a_{12}y + a_{13}z \\ a_{21}x + a_{22}y + a_{23}z \\ a_{31}x + a_{32}y + a_{33}z \end{pmatrix}. \]
Cette propriété est \textbf{admise} au programme de Terminale ; nous la vérifions ci-dessous sur un exemple.
\end{propriete}

\begin{experience}[Vérification de la formule $Y = AX$ sur un exemple]
Reprenons $f(x,y,z) = (x+y,\; y-z,\; 2x+z)$ de matrice $A = \begin{pmatrix} 1 & 1 & 0 \\ 0 & 1 & -1 \\ 2 & 0 & 1 \end{pmatrix}$, et le vecteur $u = (x, y, z)$ quelconque.
\begin{align*}
AX &= \begin{pmatrix} 1 & 1 & 0 \\ 0 & 1 & -1 \\ 2 & 0 & 1 \end{pmatrix} \begin{pmatrix} x \\ y \\ z \end{pmatrix}
= \begin{pmatrix} 1 \times x + 1 \times y + 0 \times z \\ 0 \times x + 1 \times y + (-1) \times z \\ 2 \times x + 0 \times y + 1 \times z \end{pmatrix}
= \begin{pmatrix} x + y \\ y - z \\ 2x + z \end{pmatrix}.
\end{align*}
Le produit matriciel redonne exactement les coordonnées de $f(u)$. La formule $Y = AX$ est vérifiée.
\end{experience}

\begin{attention}[Bien relire les formules avant de remplir la matrice]
Dans $f(x,y,z) = (x + y,\; y - z,\; 2x + z)$, le coefficient de $z$ dans la première formule vaut $0$ (le terme $z$ est absent). Oublier ce zéro fausse toute la troisième colonne. Prends l'habitude de réécrire chaque formule avec les trois variables : $x + y = 1\cdot x + 1\cdot y + 0\cdot z$.
\end{attention}

\begin{popfigure}
\begin{center}
\begin{tikzpicture}[node distance=1.6cm, >=Stealth]
\node[draw=popblue, thick, rounded corners, fill=popblueL, inner sep=8pt, align=center] (u) {$u$ \; de coordonnées \; $X = \begin{pmatrix} x \\ y \\ z \end{pmatrix}$};
\node[draw=popgreen, thick, rounded corners, fill=popgreenL, inner sep=8pt, right=3.2cm of u, align=center] (fu) {$f(u)$ \; de coordonnées \; $Y = \begin{pmatrix} x' \\ y' \\ z' \end{pmatrix}$};
\draw[->, very thick, popdark] (u) -- node[above, popdark]{$f$} node[below, popdark]{$Y = AX$} (fu);
\node[below=1.1cm of u, popmuted, align=center] {\small espace de départ $E$\\ \small base $\mathcal{B}$};
\node[below=1.1cm of fu, popmuted, align=center] {\small espace d'arrivée $F$\\ \small base $\mathcal{B}'$};
\end{tikzpicture}
\end{center}
\legende{La matrice $A$ agit sur les colonnes de coordonnées : les coordonnées de l'image s'obtiennent par le produit $Y = AX$.}
\end{popfigure}

\courssub{Deux matrices à connaître : identité et homothéties}

\begin{exemplebox}[Matrice de l'identité]
L'\textbf{identité} de $\mathbb{R}^3$ est l'application $\mathrm{id} : u \mapsto u$. Comme $\mathrm{id}(e_1) = e_1 = (1,0,0)$, $\mathrm{id}(e_2) = e_2$, $\mathrm{id}(e_3) = e_3$, sa matrice dans toute base est
\[ I_3 = \begin{pmatrix} 1 & 0 & 0 \\ 0 & 1 & 0 \\ 0 & 0 & 1 \end{pmatrix}, \]
appelée \cle{matrice identité} d'ordre 3. On vérifie que $I_3 X = X$ pour toute colonne $X$ : l'identité ne modifie rien, sa matrice non plus.
\end{exemplebox}

\begin{exemplebox}[Matrice d'une homothétie]
L'\textbf{homothétie de rapport $k$} est l'application linéaire $h : u \mapsto k\,u$, c'est-à-dire $h(x,y,z) = (kx, ky, kz)$. Ses images des vecteurs de base sont $h(e_1) = (k, 0, 0)$, $h(e_2) = (0, k, 0)$, $h(e_3) = (0, 0, k)$, d'où
\[ \mathrm{Mat}_{\mathcal{B}}(h) = \begin{pmatrix} k & 0 & 0 \\ 0 & k & 0 \\ 0 & 0 & k \end{pmatrix} = k\,I_3. \]
Par exemple, l'homothétie de rapport $3$ (qui triple toutes les distances, comme un agrandissement de plan) a pour matrice $3I_3$.
\end{exemplebox}

\courssub{Lecture inverse : retrouver les formules à partir de la matrice}

\begin{methode}[Retrouver les formules de $f$ connaissant sa matrice]
Si $A = \begin{pmatrix} a & b & c \\ a' & b' & c' \\ a'' & b'' & c'' \end{pmatrix}$ est la matrice de $f$ dans la base canonique, alors les formules de $f$ se lisent \textbf{en lignes} :
\begin{enumerate}
\item Première ligne : $f_1(x,y,z) = ax + by + cz$ ;
\item Deuxième ligne : $f_2(x,y,z) = a'x + b'y + c'z$ ;
\item Troisième ligne : $f_3(x,y,z) = a''x + b''y + c''z$.
\end{enumerate}
On peut aussi calculer directement $AX$ avec $X = \begin{pmatrix} x \\ y \\ z \end{pmatrix}$.
\end{methode}

\begin{exoresolu}[De la matrice aux formules, puis calcul d'une image]
\textbf{Énoncé.} Soit $f$ l'endomorphisme de $\mathbb{R}^3$ dont la matrice dans la base canonique est
\[ A = \begin{pmatrix} 1 & 2 & 0 \\ -1 & 0 & 3 \\ 0 & 1 & 1 \end{pmatrix}. \]
\begin{enumerate}
\item Déterminer les formules de $f(x, y, z)$.
\item Calculer $f(1, 2, -1)$ de deux manières : par les formules, puis par le produit matriciel.
\end{enumerate}
\tcblower
\textbf{Solution.}

\textbf{1.} On lit les lignes de $A$ :
\[ f(x, y, z) = (x + 2y,\; -x + 3z,\; y + z). \]

\textbf{2.} \emph{Par les formules :} en remplaçant $x = 1$, $y = 2$, $z = -1$ :
\begin{align*}
f(1, 2, -1) &= (1 + 2 \times 2,\; -1 + 3 \times (-1),\; 2 + (-1)) = (5,\; -4,\; 1).
\end{align*}
\emph{Par le produit matriciel :}
\[ AX = \begin{pmatrix} 1 & 2 & 0 \\ -1 & 0 & 3 \\ 0 & 1 & 1 \end{pmatrix} \begin{pmatrix} 1 \\ 2 \\ -1 \end{pmatrix} = \begin{pmatrix} 1 \times 1 + 2 \times 2 + 0 \times (-1) \\ -1 \times 1 + 0 \times 2 + 3 \times (-1) \\ 0 \times 1 + 1 \times 2 + 1 \times (-1) \end{pmatrix} = \begin{pmatrix} 5 \\ -4 \\ 1 \end{pmatrix}. \]
Les deux méthodes donnent le même résultat : $f(1,2,-1) = (5, -4, 1)$. Le produit matriciel est souvent plus rapide et plus sûr lorsqu'il y a beaucoup d'images à calculer.
\end{exoresolu}

\courssub{Complément culturel : la matrice comme photographie de la transformation}

\begin{savaistu}[Les matrices sont partout — hors-programme strict, utile pour la culture du Bac]
Géométriquement, les colonnes de la matrice $A$ disent \emph{où partent} les vecteurs de base : la transformation déforme l'espace en envoyant $e_1, e_2, e_3$ sur les trois colonnes de $A$. C'est exactement ce principe qu'utilisent les logiciels d'infographie 3D (jeux vidéo, conception de pièces mécaniques, animation) : chaque rotation, chaque agrandissement, chaque projection d'un objet est codé par une matrice $3 \times 3$, et le processeur graphique effectue des millions de produits $AX$ par seconde. Au Cameroun, les ingénieurs en télécommunications et en cartographie (par exemple pour les systèmes GPS utilisés dans les grandes villes comme Yaoundé et Douala) manipulent quotidiennement de telles matrices pour convertir des coordonnées d'un repère à un autre. La notion de matrice a été formalisée au XIX\textsuperscript{e} siècle par le mathématicien britannique Arthur Cayley ; le mot « matrice » lui-même fut introduit par James Joseph Sylvester en 1850.
\end{savaistu}

\begin{exercice}[À toi de jouer]
Soit $g : \mathbb{R}^3 \to \mathbb{R}^3$ définie par $g(x, y, z) = (2x - z,\; x + y + z,\; 3y)$.
\begin{enumerate}
\item Écrire la matrice $B$ de $g$ dans la base canonique.
\item Calculer $g(1, -1, 2)$ par le produit matriciel $BX$.
\item Soit $h$ l'homothétie de rapport $-2$. Écrire la matrice de $h$ dans la base canonique.
\end{enumerate}
\end{exercice}

\begin{corrige}[Exercice : À toi de jouer]
\textbf{1.} On calcule les images des vecteurs de base : $g(e_1) = g(1,0,0) = (2, 1, 0)$ ; $g(e_2) = g(0,1,0) = (0, 1, 3)$ ; $g(e_3) = g(0,0,1) = (-1, 1, 0)$. On place ces coordonnées \textbf{en colonnes} :
\[ B = \begin{pmatrix} 2 & 0 & -1 \\ 1 & 1 & 1 \\ 0 & 3 & 0 \end{pmatrix}. \]
\textbf{2.} Avec $X = \begin{pmatrix} 1 \\ -1 \\ 2 \end{pmatrix}$ :
\[ BX = \begin{pmatrix} 2 \times 1 + 0 \times (-1) + (-1) \times 2 \\ 1 \times 1 + 1 \times (-1) + 1 \times 2 \\ 0 \times 1 + 3 \times (-1) + 0 \times 2 \end{pmatrix} = \begin{pmatrix} 0 \\ 2 \\ -3 \end{pmatrix}, \]
donc $g(1, -1, 2) = (0, 2, -3)$. Vérification par les formules : $(2 - 2,\; 1 - 1 + 2,\; -3) = (0, 2, -3)$. Correct.

\textbf{3.} La matrice de l'homothétie de rapport $-2$ est $-2\,I_3 = \begin{pmatrix} -2 & 0 & 0 \\ 0 & -2 & 0 \\ 0 & 0 & -2 \end{pmatrix}$.
\end{corrige}

\begin{aretenir}[L'essentiel de la section]
\begin{itemize}
\item La matrice de $f$ dans les bases $\mathcal{B}$ (départ) et $\mathcal{B}'$ (arrivée) a pour \textbf{colonnes} les coordonnées de $f(u_1), f(u_2), f(u_3)$ dans $\mathcal{B}'$.
\item Si $X$ est la colonne des coordonnées de $u$, alors les coordonnées de $f(u)$ sont données par $Y = AX$.
\item Matrice de l'identité : $I_3$ ; matrice de l'homothétie de rapport $k$ : $kI_3$.
\item Les formules de $f$ se lisent sur les \textbf{lignes} de $A$ ; les images des vecteurs de base se lisent sur les \textbf{colonnes}.
\end{itemize}
\end{aretenir}

Nous sommes désormais capables de coder une application linéaire par une matrice et de calculer des images efficacement. Dans la prochaine section, nous exploiterons cette matrice pour étudier deux ensembles fondamentaux attachés à $f$ : son \cle{noyau} et son \cle{image}, couronnés par le célèbre théorème du rang.

\coursec{Noyau, image et théorème du rang}

Dans la section précédente, nous avons appris à traduire une application linéaire $f$ par sa \cle{matrice} $A$ dans une base donnée : calculer $f(\vec{u})$ revient alors à multiplier $A$ par la colonne des coordonnées de $\vec{u}$. Cette traduction algébrique va maintenant nous permettre de répondre à deux questions naturelles et fondamentales :

\begin{itemize}
\item Quels sont les vecteurs que $f$ \og écrase \fg{} sur le vecteur nul ? Autrement dit, que vaut l'ensemble des antécédents de $\vec{0}$ ?
\item Quelle portion de l'espace d'arrivée est réellement \og atteinte \fg{} par $f$ ?
\end{itemize}

Ces deux ensembles portent des noms : le \cle{noyau} et l'\cle{image}. Ils mesurent respectivement ce que $f$ \emph{perd} et ce que $f$ \emph{couvre}. Le théorème du rang, point culminant de cette leçon, affirme que ces deux quantités se compensent exactement : ce qui est perdu d'un côté manque de l'autre. C'est l'un des plus beaux théorèmes de l'année, et l'un des plus utiles au Baccalauréat.

\courssub{Le noyau d'une application linéaire}

\textbf{Intuition.} Imagine un projecteur qui éclaire un mur : tous les points de l'espace situés sur un même rayon lumineux sont envoyés sur le même point du mur. En particulier, tous les points alignés avec la source dans la direction perpendiculaire au mur sont écrasés en un seul point. L'ensemble des vecteurs que l'application \og oublie \fg{} en les envoyant sur $\vec{0}$ forme son noyau.

\begin{definition}[Noyau d'une application linéaire]
Soit $f : E \to F$ une application linéaire entre deux espaces vectoriels. On appelle \cle{noyau} de $f$, noté $\ker f$, l'ensemble des vecteurs de $E$ dont l'image est le vecteur nul de $F$ :
\[
\ker f = \{\vec{u} \in E \mid f(\vec{u}) = \vec{0}_F\}.
\]
\end{definition}

\begin{propriete}[Le noyau est un sous-espace vectoriel]
Pour toute application linéaire $f : E \to F$, $\ker f$ est un \cle{sous-espace vectoriel} de $E$.
\end{propriete}

\begin{experience}[Démonstration : le noyau est un sous-espace vectoriel]
Il faut vérifier les trois points de la caractérisation des sous-espaces vectoriels.
\begin{enumerate}
\item \textbf{Non vide :} $f(\vec{0}_E) = \vec{0}_F$ (propriété des applications linéaires), donc $\vec{0}_E \in \ker f$.
\item \textbf{Stabilité par addition :} soient $\vec{u}, \vec{v} \in \ker f$. Alors $f(\vec{u}) = \vec{0}_F$ et $f(\vec{v}) = \vec{0}_F$. Par linéarité :
\[
f(\vec{u} + \vec{v}) = f(\vec{u}) + f(\vec{v}) = \vec{0}_F + \vec{0}_F = \vec{0}_F,
\]
donc $\vec{u} + \vec{v} \in \ker f$.
\item \textbf{Stabilité par multiplication par un scalaire :} soient $\vec{u} \in \ker f$ et $\lambda \in \mathbb{R}$. Alors :
\[
f(\lambda \vec{u}) = \lambda f(\vec{u}) = \lambda \vec{0}_F = \vec{0}_F,
\]
donc $\lambda \vec{u} \in \ker f$.
\end{enumerate}
Conclusion : $\ker f$ est un sous-espace vectoriel de $E$. \hfill$\blacksquare$
\end{experience}

Ce résultat est précieux : il garantit que le noyau possède une \emph{dimension}, une \emph{base}, et qu'on peut le décrire par des vecteurs générateurs. Dire \og le noyau est une droite \fg{} ou \og le noyau est un plan \fg{} a alors un sens précis.

\begin{methode}[Déterminer le noyau en pratique]
Pour déterminer $\ker f$ lorsque $f : \mathbb{R}^3 \to \mathbb{R}^3$ est donnée par son expression analytique ou par sa matrice $A$ :
\begin{enumerate}
\item Écrire l'équation $f(x, y, z) = (0, 0, 0)$, c'est-à-dire résoudre le système homogène $AX = 0$.
\item Résoudre ce système (méthode du pivot ou substitutions).
\item Exprimer l'ensemble des solutions sous forme de combinaison linéaire : $(x, y, z) = t\,\vec{u}$ (droite vectorielle) ou $(x,y,z) = t\,\vec{u} + s\,\vec{v}$ (plan vectoriel).
\item Conclure : $\ker f = \mathrm{Vect}(\vec{u})$ ou $\mathrm{Vect}(\vec{u}, \vec{v})$, et en donner la dimension.
\end{enumerate}
\end{methode}

\begin{attention}[Pièges fréquents sur le noyau]
\begin{itemize}
\item Le noyau n'est \textbf{jamais vide} : il contient toujours au moins $\vec{0}_E$. Écrire $\ker f = \emptyset$ est une erreur grave.
\item Ne confonds pas $\ker f = \{\vec{0}_E\}$ (noyau \emph{réduit au vecteur nul}) avec $\ker f = \emptyset$. Le premier cas est possible et même souhaitable ; le second est impossible.
\item Le noyau vit dans l'espace de \textbf{départ} $E$, pas dans l'espace d'arrivée.
\end{itemize}
\end{attention}

\courssub{L'image d'une application linéaire}

\textbf{Intuition.} Si le noyau est ce que $f$ écrase, l'image est ce que $f$ atteint. Quand on projette l'espace tout entier sur un mur, on n'obtient que le mur : l'image est un plan, pas tout l'espace. L'image mesure la \og portée \fg{} de l'application.

\begin{definition}[Image d'une application linéaire]
Soit $f : E \to F$ une application linéaire. On appelle \cle{image} de $f$, notée $\mathrm{Im}\, f$, l'ensemble des images de tous les vecteurs de $E$ :
\[
\mathrm{Im}\, f = \{f(\vec{u}) \mid \vec{u} \in E\} = \{\vec{w} \in F \mid \exists\, \vec{u} \in E,\ f(\vec{u}) = \vec{w}\}.
\]
\end{definition}

\begin{propriete}[L'image est un sous-espace vectoriel]
Pour toute application linéaire $f : E \to F$, $\mathrm{Im}\, f$ est un \cle{sous-espace vectoriel} de $F$.
\end{propriete}

\begin{experience}[Démonstration : l'image est un sous-espace vectoriel]
\begin{enumerate}
\item \textbf{Non vide :} $\vec{0}_F = f(\vec{0}_E) \in \mathrm{Im}\, f$.
\item \textbf{Stabilité par addition :} soient $\vec{w}_1, \vec{w}_2 \in \mathrm{Im}\, f$. Il existe $\vec{u}_1, \vec{u}_2 \in E$ tels que $\vec{w}_1 = f(\vec{u}_1)$ et $\vec{w}_2 = f(\vec{u}_2)$. Alors :
\[
\vec{w}_1 + \vec{w}_2 = f(\vec{u}_1) + f(\vec{u}_2) = f(\vec{u}_1 + \vec{u}_2) \in \mathrm{Im}\, f.
\]
\item \textbf{Stabilité par multiplication par un scalaire :} soient $\vec{w} = f(\vec{u}) \in \mathrm{Im}\, f$ et $\lambda \in \mathbb{R}$. Alors :
\[
\lambda \vec{w} = \lambda f(\vec{u}) = f(\lambda \vec{u}) \in \mathrm{Im}\, f.
\]
\end{enumerate}
Conclusion : $\mathrm{Im}\, f$ est un sous-espace vectoriel de $F$. \hfill$\blacksquare$
\end{experience}

La propriété suivante est la clé pratique pour déterminer une image : il suffit de regarder les images des vecteurs d'une base.

\begin{propriete}[L'image est engendrée par les images d'une base]
Soit $f : E \to F$ une application linéaire et $(\vec{e}_1, \vec{e}_2, \vec{e}_3)$ une base de $E$. Alors :
\[
\mathrm{Im}\, f = \mathrm{Vect}\big(f(\vec{e}_1),\ f(\vec{e}_2),\ f(\vec{e}_3)\big).
\]
Autrement dit, si $A$ est la matrice de $f$, alors $\mathrm{Im}\, f$ est engendrée par les \textbf{colonnes} de $A$.
\end{propriete}

\textbf{Justification.} Tout vecteur $\vec{u} \in E$ s'écrit $\vec{u} = x\vec{e}_1 + y\vec{e}_2 + z\vec{e}_3$. Par linéarité :
\[
f(\vec{u}) = x\,f(\vec{e}_1) + y\,f(\vec{e}_2) + z\,f(\vec{e}_3),
\]
donc toute image est combinaison linéaire de $f(\vec{e}_1), f(\vec{e}_2), f(\vec{e}_3)$, et réciproquement toute telle combinaison est une image. \hfill$\blacksquare$

\begin{methode}[Déterminer l'image en pratique]
\begin{enumerate}
\item Calculer les images des vecteurs de la base canonique : ce sont les \textbf{colonnes} de la matrice $A$ de $f$.
\item Écrire $\mathrm{Im}\, f = \mathrm{Vect}(C_1, C_2, C_3)$ où $C_1, C_2, C_3$ sont ces colonnes.
\item \textbf{Extraire une famille libre} : repérer les colonnes liées (par exemple une colonne somme des deux autres) et les éliminer.
\item Conclure : les colonnes restantes forment une base de $\mathrm{Im}\, f$, et leur nombre est $\dim \mathrm{Im}\, f$.
\end{enumerate}
\end{methode}

\courssub{Injectivité, surjectivité : le rôle du noyau et de l'image}

Rappelons qu'une application est \cle{injective} si deux vecteurs distincts ont toujours des images distinctes, et \cle{surjective} si tout vecteur de l'arrivée possède au moins un antécédent. Pour une application \emph{linéaire}, ces propriétés se lisent entièrement sur le noyau et l'image.

\begin{propriete}[Critère d'injectivité — démonstration exigible]
Soit $f : E \to F$ une application linéaire. Alors :
\[
f \text{ est injective} \iff \ker f = \{\vec{0}_E\}.
\]
\end{propriete}

\begin{experience}[Démonstration complète du critère d'injectivité]
On démontre les deux implications.

\textbf{Sens direct ($\Rightarrow$).} Supposons $f$ injective. Soit $\vec{u} \in \ker f$, c'est-à-dire $f(\vec{u}) = \vec{0}_F$. Or on sait que $f(\vec{0}_E) = \vec{0}_F$. Donc $f(\vec{u}) = f(\vec{0}_E)$, et par injectivité, $\vec{u} = \vec{0}_E$. Ainsi $\ker f = \{\vec{0}_E\}$.

\textbf{Sens réciproque ($\Leftarrow$).} Supposons $\ker f = \{\vec{0}_E\}$. Soient $\vec{u}, \vec{v} \in E$ tels que $f(\vec{u}) = f(\vec{v})$. Par linéarité :
\[
f(\vec{u} - \vec{v}) = f(\vec{u}) - f(\vec{v}) = \vec{0}_F,
\]
donc $\vec{u} - \vec{v} \in \ker f = \{\vec{0}_E\}$, d'où $\vec{u} - \vec{v} = \vec{0}_E$, c'est-à-dire $\vec{u} = \vec{v}$. L'application $f$ est donc injective. \hfill$\blacksquare$
\end{experience}

\begin{aretenir}[L'idée à retenir]
Pour une application linéaire, tester l'injectivité ne demande pas de comparer toutes les paires de vecteurs : il suffit de vérifier que \textbf{seul le vecteur nul} est envoyé sur le vecteur nul. La linéarité ramène une propriété globale à un test en un seul point.
\end{aretenir}

\begin{propriete}[Critère de surjectivité]
Soit $f : E \to F$ une application linéaire. Alors :
\[
f \text{ est surjective} \iff \mathrm{Im}\, f = F.
\]
En dimension finie, cela équivaut à : $\dim \mathrm{Im}\, f = \dim F$.
\end{propriete}

\textbf{Justification.} C'est presque la définition : $f$ surjective signifie que tout vecteur de $F$ est atteint, c'est-à-dire appartient à $\mathrm{Im}\, f$. Or $\mathrm{Im}\, f$ est un sous-espace vectoriel de $F$, et un sous-espace vectoriel de $F$ qui a la même dimension que $F$ est égal à $F$ tout entier. \hfill$\blacksquare$

\courssub{Le théorème du rang}

Nous voici au sommet de la leçon. Le théorème du rang affirme que la dimension de l'espace de départ se répartit \emph{exactement} entre ce que $f$ écrase (le noyau) et ce que $f$ atteint (l'image).

\begin{propriete}[Théorème du rang — admis]
Soit $f : E \to F$ une application linéaire, avec $E$ de dimension finie. Alors :
\[
\dim E = \dim \ker f + \dim \mathrm{Im}\, f.
\]
En particulier, pour $f : \mathbb{R}^3 \to \mathbb{R}^3$ : $\dim \ker f + \dim \mathrm{Im}\, f = 3$.
\end{propriete}

\textbf{Justification intuitive.} Décomposons mentalement l'espace de départ $E$ en deux morceaux : le noyau $\ker f$ (les vecteurs écrasés sur $\vec{0}$) et un \emph{complément} $S$ (un sous-espace qui le complète, de sorte que $E = \ker f \oplus S$). Sur le noyau, $f$ ne fait rien d'intéressant : tout part à zéro. Sur le complément $S$, en revanche, $f$ est injective (aucun vecteur non nul de $S$ ne peut être écrasé, sinon il serait dans le noyau) : elle transporte donc $S$ \emph{sans perte} sur $\mathrm{Im}\, f$. D'où $\dim S = \dim \mathrm{Im}\, f$, et comme $\dim E = \dim \ker f + \dim S$, on obtient la formule.

\begin{popfigure}
\begin{center}
\begin{tikzpicture}[scale=1.0, >=Stealth]
% Bloc E
\draw[thick, popblue, fill=popblueL, rounded corners=8pt] (0,0) rectangle (4.2,3.2);
\node[popblue, font=\bfseries] at (2.1,3.55) {$E$ (départ), $\dim E = 3$};
% Ker f
\draw[thick, poppink, fill=poppinkL, rounded corners=6pt] (0.25,0.25) rectangle (1.9,2.95);
\node[poppink, font=\bfseries] at (1.07,2.55) {$\ker f$};
\node[poppink] at (1.07,1.4) {\small écrasé};
\node[poppink] at (1.07,1.0) {\small sur $\vec{0}$};
% Complément
\draw[thick, popgreen, fill=popgreenL, rounded corners=6pt] (2.3,0.25) rectangle (3.95,2.95);
\node[popgreen!60!black, font=\bfseries] at (3.12,2.55) {complément $S$};
\node[popgreen!60!black] at (3.12,1.2) {\small transporté};
\node[popgreen!60!black] at (3.12,0.8) {\small sans perte};
% Bloc F
\draw[thick, poporange, fill=poporangeL, rounded corners=8pt] (7.2,0) rectangle (11.4,3.2);
\node[poporange, font=\bfseries] at (9.3,3.55) {$F$ (arrivée)};
% Im f
\draw[thick, poppurple, fill=poppurpleL, rounded corners=6pt] (7.45,0.25) rectangle (9.6,2.95);
\node[poppurple, font=\bfseries] at (8.52,2.55) {$\mathrm{Im}\, f$};
\node[poppurple] at (8.52,1.3) {\small $\dim \mathrm{Im}\, f = \dim S$};
% Point 0
\fill[poppink] (10.6,0.6) circle (2pt);
\node[below, poppink] at (10.6,0.55) {$\vec{0}_F$};
% Flèches
\draw[->, very thick, poppink] (1.07,0.6) .. controls (3.5,-0.9) and (8,-0.9) .. (10.5,0.55);
\draw[->, very thick, popgreen!60!black] (3.12,1.6) .. controls (5,2.3) and (6,2.3) .. (8.52,1.9);
\node[popgreen!60!black] at (5.7,2.75) {\small $f$ injective sur $S$};
% Formule
\node[popdark, font=\bfseries] at (5.7,-1.5) {$\dim E = \underbrace{\dim \ker f}_{\text{perdu}} + \underbrace{\dim \mathrm{Im}\, f}_{\text{conservé}}$};
\end{tikzpicture}
\end{center}
\legende{Le théorème du rang : l'espace de départ se partage entre le noyau (écrasé sur $\vec{0}_F$) et un complément transporté fidèlement sur l'image.}
\end{popfigure}

\begin{exemplebox}[La projection sur un plan : l'exemple modèle]
Soit $p : \mathbb{R}^3 \to \mathbb{R}^3$ la projection orthogonale sur le plan horizontal $(Oxy)$, définie par $p(x, y, z) = (x, y, 0)$.
\begin{itemize}
\item \textbf{Noyau :} $p(x,y,z) = (0,0,0) \iff x = 0$ et $y = 0$. Donc $\ker p = \{(0, 0, z) \mid z \in \mathbb{R}\} = \mathrm{Vect}(\vec{k})$ : c'est l'axe $(Oz)$, une \textbf{droite}, $\dim \ker p = 1$.
\item \textbf{Image :} $\mathrm{Im}\, p = \{(x, y, 0) \mid x, y \in \mathbb{R}\} = \mathrm{Vect}(\vec{i}, \vec{j})$ : c'est le plan $(Oxy)$, $\dim \mathrm{Im}\, p = 2$.
\item \textbf{Vérification :} $1 + 2 = 3 = \dim \mathbb{R}^3$. Le théorème du rang est vérifié.
\end{itemize}
\end{exemplebox}

\begin{popfigure}
\begin{center}
\begin{tikzpicture}[scale=1.1, >=Stealth]
% Plan (Oxy)
\draw[thick, popblue, fill=popblueL, opacity=0.85] (-2.4,-0.9) -- (0.6,-1.7) -- (2.9,0.1) -- (-0.1,0.9) -- cycle;
\node[popblue] at (2.2,-0.9) {\small plan image $(Oxy)$};
% Axe z
\draw[->, thick, poppink] (0,0) -- (0,2.9) node[above] {$z$ : $\ker p$};
\draw[->, thick, popdark] (0,0) -- (1.9,-0.55) node[right] {$x$};
\draw[->, thick, popdark] (0,0) -- (-1.35,0.62) node[left] {$y$};
% Point et projection
\fill[popgreen!60!black] (1.0,2.2) circle (2pt) node[above right, popgreen!60!black] {$M(x,y,z)$};
\draw[->, dashed, thick, popgreen!60!black] (1.0,2.2) -- (1.0,-0.28);
\fill[popgreen!60!black] (1.0,-0.28) circle (2pt) node[below right, popgreen!60!black] {$p(M) = (x, y, 0)$};
% Vecteur du noyau
\draw[->, very thick, poppink] (0,0) -- (0,1.6) node[midway, right, poppink] {\small $\vec{u} \in \ker p$};
\draw[->, dashed, poppink] (0,1.6) -- (0.12,0.15);
\node[poppink] at (-1.9,2.3) {\small $p(\vec{u}) = \vec{0}$};
% Formule
\node[popdark, font=\bfseries] at (0.3,-2.3) {$\dim \ker p + \dim \mathrm{Im}\, p = 1 + 2 = 3$};
\end{tikzpicture}
\end{center}
\legende{Projection sur le plan $(Oxy)$ : le noyau est la droite $(Oz)$ (dimension 1), l'image est le plan (dimension 2), et $1 + 2 = 3$.}
\end{popfigure}

\begin{propriete}[Conséquence majeure en dimension 3]
Soit $f$ un \cle{endomorphisme} de $\mathbb{R}^3$ (application linéaire de $\mathbb{R}^3$ dans $\mathbb{R}^3$). Alors les trois propriétés suivantes sont \textbf{équivalentes} :
\[
f \text{ injective} \iff f \text{ surjective} \iff f \text{ bijective}.
\]
\end{propriete}

\textbf{Justification.} Si $f$ est injective, alors $\ker f = \{\vec{0}\}$, donc $\dim \ker f = 0$, et le théorème du rang donne $\dim \mathrm{Im}\, f = 3 = \dim \mathbb{R}^3$, donc $\mathrm{Im}\, f = \mathbb{R}^3$ : $f$ est surjective. Réciproquement, si $f$ est surjective, $\dim \mathrm{Im}\, f = 3$, donc $\dim \ker f = 0$, donc $\ker f = \{\vec{0}\}$ : $f$ est injective. \hfill$\blacksquare$

\begin{aretenir}[Réflexe Bac]
Pour un endomorphisme de $\mathbb{R}^3$, il suffit de montrer \textbf{une seule} des deux propriétés (injectivité \emph{ou} surjectivité) pour conclure à la bijectivité. En pratique, on montre presque toujours que $\ker f = \{\vec{0}\}$, car cela revient à vérifier que le système $AX = 0$ n'a que la solution nulle.
\end{aretenir}

\courssub{Exemple chiffré complet et rang d'une application linéaire}

\begin{exoresolu}[Étude complète d'un endomorphisme de $\mathbb{R}^3$]
Soit $f : \mathbb{R}^3 \to \mathbb{R}^3$ définie par $f(x, y, z) = (x + y,\ y - z,\ x + 2y - z)$.
\begin{enumerate}
\item Déterminer $\ker f$ et donner sa dimension.
\item Déterminer $\mathrm{Im}\, f$ et donner sa dimension.
\item Vérifier le théorème du rang. L'application $f$ est-elle bijective ?
\end{enumerate}
\tcblower
\textbf{1. Noyau.} On résout $f(x, y, z) = (0, 0, 0)$ :
\[
\begin{cases}
x + y = 0 \\
y - z = 0 \\
x + 2y - z = 0
\end{cases}
\]
La première équation donne $x = -y$ ; la deuxième donne $z = y$. En remplaçant dans la troisième : $-y + 2y - y = 0$, ce qui est toujours vrai. Le système se réduit donc à $x = -y$ et $z = y$, avec $y$ libre. En posant $y = t$ :
\[
(x, y, z) = (-t,\ t,\ t) = t\,(-1, 1, 1).
\]
Donc $\ker f = \mathrm{Vect}\big((-1, 1, 1)\big)$ : c'est une \textbf{droite vectorielle}, et $\dim \ker f = 1$.

\textbf{2. Image.} Les images des vecteurs de la base canonique sont les colonnes de la matrice de $f$ :
\[
f(\vec{i}) = (1, 0, 1), \quad f(\vec{j}) = (1, 1, 2), \quad f(\vec{k}) = (0, -1, -1).
\]
Donc $\mathrm{Im}\, f = \mathrm{Vect}\big((1,0,1),\ (1,1,2),\ (0,-1,-1)\big)$. Cherchons les liaisons : on remarque que
\[
(1, 1, 2) = (1, 0, 1) - (0, -1, -1),
\]
c'est-à-dire $f(\vec{j}) = f(\vec{i}) - f(\vec{k})$. La troisième colonne est liée aux deux autres ; on l'élimine. Reste la famille $\big((1, 0, 1),\ (0, -1, -1)\big)$, qui est libre (ces deux vecteurs ne sont pas colinéaires). Donc :
\[
\mathrm{Im}\, f = \mathrm{Vect}\big((1, 0, 1),\ (0, -1, -1)\big) \text{ est un \textbf{plan vectoriel}, } \dim \mathrm{Im}\, f = 2.
\]

\textbf{3. Théorème du rang.} On a bien $\dim \ker f + \dim \mathrm{Im}\, f = 1 + 2 = 3 = \dim \mathbb{R}^3$. Comme $\ker f \neq \{\vec{0}\}$, $f$ n'est pas injective, donc (par la conséquence majeure) pas bijective non plus.
\end{exoresolu}

\begin{attention}[Vérifie toujours la cohérence]
Si tu trouves $\dim \ker f = 2$ et $\dim \mathrm{Im}\, f = 2$ pour un endomorphisme de $\mathbb{R}^3$, il y a forcément une erreur quelque part : la somme doit valoir $3$. Le théorème du rang est un excellent outil d'auto-contrôle au Baccalauréat.
\end{attention}

\begin{definition}[Rang d'une application linéaire]
Soit $f : E \to F$ une application linéaire. On appelle \cle{rang} de $f$, noté $\mathrm{rg}(f)$, la dimension de son image :
\[
\mathrm{rg}(f) = \dim \mathrm{Im}\, f.
\]
Si $A$ est la matrice de $f$ dans des bases données, le rang de $f$ est aussi le \textbf{nombre maximal de colonnes libres} de $A$ (on parle aussi de rang de la matrice $A$).
\end{definition}

Le théorème du rang se réécrit alors, pour $f : \mathbb{R}^3 \to \mathbb{R}^3$ :
\[
\dim \ker f + \mathrm{rg}(f) = 3.
\]

\begin{center}
\begin{tabularx}{\linewidth}{|l|X|X|X|}
\hline
\rowcolor{popblueL} \textbf{Situation} & \textbf{Noyau} & \textbf{Image} & \textbf{Conclusion} \\
\hline
$\mathrm{rg}(f) = 3$ & $\ker f = \{\vec{0}\}$ ($\dim 0$) & $\mathrm{Im}\, f = \mathbb{R}^3$ & $f$ bijective \\
\hline
$\mathrm{rg}(f) = 2$ & droite ($\dim 1$) & plan ($\dim 2$) & ni injective ni surjective \\
\hline
$\mathrm{rg}(f) = 1$ & plan ($\dim 2$) & droite ($\dim 1$) & ni injective ni surjective \\
\hline
$\mathrm{rg}(f) = 0$ & $\ker f = \mathbb{R}^3$ & $\{\vec{0}\}$ & $f$ est l'application nulle \\
\hline
\end{tabularx}
\end{center}

\begin{savaistu}[Pourquoi ce théorème porte-t-il ce nom ?]
Le \og rang \fg{} d'une matrice est un concept central de l'algèbre linéaire moderne, développé au XIX\textsuperscript{e} siècle notamment par Frobenius et Sylvester. Aujourd'hui, le théorème du rang est utilisé partout : en compression d'images (une image est une grande matrice, et on l'approche par des matrices de petit rang pour la stocker avec moins de données), en traitement du signal chez les ingénieurs télécoms, et dans les algorithmes de recommandation des plateformes numériques. Quand ton téléphone compresse une photo avant de l'envoyer par WhatsApp, c'est une version géante du théorème que tu viens d'apprendre qui travaille en coulisses !
\end{savaistu}

\begin{exercice}[À toi de jouer]
Soit $g : \mathbb{R}^3 \to \mathbb{R}^3$ définie par $g(x, y, z) = (x - y + z,\ 2x + y - z,\ x + 2y - 2z)$.
\begin{enumerate}
\item Déterminer $\ker g$ et sa dimension.
\item En déduire $\mathrm{rg}(g)$ sans calculer l'image, puis déterminer une base de $\mathrm{Im}\, g$.
\item L'application $g$ est-elle bijective ?
\end{enumerate}
\end{exercice}

\begin{corrige}[Exercice : étude de $g$]
\textbf{1.} On résout $g(x, y, z) = (0,0,0)$ :
\[
\begin{cases}
x - y + z = 0 \\
2x + y - z = 0 \\
x + 2y - 2z = 0
\end{cases}
\]
En additionnant les deux premières équations : $3x = 0$, donc $x = 0$. Le système devient $-y + z = 0$ et $2y - 2z = 0$, soit $z = y$. Avec $y = t$ : $(x, y, z) = t\,(0, 1, 1)$. Donc $\ker g = \mathrm{Vect}\big((0,1,1)\big)$, droite vectorielle, $\dim \ker g = 1$.

\textbf{2.} Par le théorème du rang : $\mathrm{rg}(g) = 3 - 1 = 2$. Les colonnes de la matrice sont $g(\vec{i}) = (1, 2, 1)$, $g(\vec{j}) = (-1, 1, 2)$, $g(\vec{k}) = (1, -1, -2)$. On remarque $g(\vec{k}) = -g(\vec{j})$, et les deux premières colonnes ne sont pas colinéaires : une base de $\mathrm{Im}\, g$ est $\big((1, 2, 1),\ (-1, 1, 2)\big)$.

\textbf{3.} Comme $\ker g \neq \{\vec{0}\}$, $g$ n'est pas injective, donc pas bijective.
\end{corrige}

\begin{aretenir}[Bilan de la section]
\begin{itemize}
\item $\ker f = \{\vec{u} \in E \mid f(\vec{u}) = \vec{0}\}$ est un sous-espace vectoriel de $E$ ; on le détermine en résolvant $AX = 0$.
\item $\mathrm{Im}\, f = \{f(\vec{u}) \mid \vec{u} \in E\}$ est un sous-espace vectoriel de $F$, engendré par les \textbf{colonnes} de la matrice de $f$ ; on en extrait une famille libre.
\item $f$ injective $\iff \ker f = \{\vec{0}\}$ ; $f$ surjective $\iff \mathrm{Im}\, f = F$.
\item \textbf{Théorème du rang :} $\dim E = \dim \ker f + \dim \mathrm{Im}\, f$.
\item Pour un endomorphisme de $\mathbb{R}^3$ : injectif $\iff$ surjectif $\iff$ bijectif.
\item Le \textbf{rang} de $f$ est $\dim \mathrm{Im}\, f$, c'est-à-dire le nombre de colonnes libres de sa matrice.
\end{itemize}
\end{aretenir}

\newpage

\coursec{Activité d'intégration}

\courssub{Objectifs de l'activité}

À l'issue de cette activité d'intégration, tu seras capable de mobiliser \emph{ensemble} les savoir-faire de la leçon sur une situation concrète unique. Plus précisément, tu devras :

\begin{itemize}
\item reconnaître et démontrer qu'une application de $\mathbb{R}^3$ dans $\mathbb{R}^3$ est \cle{linéaire} ;
\item écrire la \cle{matrice} d'une application linéaire dans la base canonique de $\mathbb{R}^3$ ;
\item déterminer le \cle{noyau} et l'\cle{image} d'une application linéaire, et en donner une interprétation géométrique ;
\item vérifier le \cle{théorème du rang} : $\dim E = \dim \ker f + \dim \operatorname{Im} f$ ;
\item montrer qu'une famille de trois vecteurs est une \cle{base} de $\mathbb{R}^3$ et déterminer les \cle{coordonnées} d'un vecteur dans cette base.
\end{itemize}

L'organisation retenue est le travail en groupes de 3 ou 4 élèves, suivi d'une restitution au tableau et d'une mise en commun dirigée par l'enseignant. Durée indicative : 2 heures.

\courssub{Situation}

\textbf{Le décorateur d'intérieur et la projection des ombres au sol.}

Monsieur Kamdem est décorateur d'intérieur à Bafoussam, dans la région de l'Ouest. Pour l'inauguration d'un restaurant du quartier Tougang, il conçoit un \emph{mobile décoratif} : un assemblage de petits objets suspendus au plafond par des fils de nylon. La nuit venue, une lampe halogène placée très haut, au-dessus du mobile, projette l'ombre des objets sur le sol carrelé de la salle, créant un dessin lumineux très apprécié des clients.

Pour préparer son installation, Monsieur Kamdem modélise la situation à l'aide des mathématiques. Il munit l'espace de la salle d'un repère $(O\,;\vec{i},\vec{j},\vec{k})$ où :

\begin{itemize}
\item l'origine $O$ est le coin de la salle situé au pied du mur, à l'aplomb de la lampe ;
\item le sol carrelé est le plan d'équation $z = 0$ ;
\item l'axe vertical est dirigé par le vecteur $\vec{k}$, l'unité étant le mètre.
\end{itemize}

La lampe étant placée très haut sur la verticale de $O$, les rayons lumineux sont supposés verticaux : l'ombre d'un point $M(x\,;y\,;z)$ du mobile sur le sol est donc le point $M'(x\,;y\,;0)$. On note $f$ l'application qui, à tout point de l'espace, associe son ombre sur le sol :
\[
f : \mathbb{R}^3 \longrightarrow \mathbb{R}^3, \qquad f(x,y,z) = (x,\,y,\,0).
\]

Le plafond de la salle est soutenu par trois poutres métalliques dont les directions, mesurées par Monsieur Kamdem, sont données par les vecteurs :
\[
\vec{u} = (1,1,0), \qquad \vec{v} = (0,1,1), \qquad \vec{w} = (1,0,1).
\]

Un des éléments du mobile, une petite statuette en bois de Foumban, est accrochée au point $A(3\,;2\,;2)$. Monsieur Kamdem souhaite exprimer la position de $A$ en fonction des directions de ses poutres, afin de régler les longueurs des fils de suspension.

\textbf{Problème posé :} modéliser mathématiquement la projection de l'ombre, décrire l'ensemble des points dont l'ombre tombe exactement au pied du mur (point $O$), décrire la forme de l'ombre, et exprimer la position de la statuette dans la base des poutres.

\begin{popfigure}
\begin{center}
\begin{tikzpicture}[scale=1.05, >=Stealth]
% Sol (plan z=0) vu en perspective
\fill[popblueL, opacity=0.55] (0,0) -- (4.6,0) -- (6.1,1.7) -- (1.5,1.7) -- cycle;
\draw[popline] (0,0) -- (4.6,0) -- (6.1,1.7) -- (1.5,1.7) -- cycle;
\node[popmuted, font=\small] at (5.3,0.45) {sol : plan $z=0$};
% Origine et axes
\coordinate (O) at (1.5,1.7);
\fill[popdark] (O) circle (1.6pt) node[below left, popdark] {$O$};
\draw[->, popdark, thick] (O) -- (3.1,1.7) node[right, popdark] {$\vec{i}$};
\draw[->, popdark, thick] (O) -- (2.3,2.5) node[right, popdark] {$\vec{j}$};
\draw[->, popdark, thick] (O) -- (1.5,3.6) node[above, popdark] {$\vec{k}$};
% Lampe
\fill[popgold] (1.5,4.6) circle (2.4pt) node[above, popgold!70!popdark] {lampe};
\draw[popgold, dashed] (1.5,4.6) -- (O);
% Point A et son ombre
\coordinate (A) at (4.3,3.4);
\coordinate (Aprime) at (4.3,0.85);
\fill[popink] (A) circle (2pt) node[above right, popink] {$A(3\,;2\,;2)$};
\draw[popink, dashed, thick] (A) -- (Aprime);
\fill[poporange] (Aprime) circle (2pt) node[below, poporange!80!popdark] {$A'(3\,;2\,;0)$};
\node[popmuted, font=\small, align=center] at (0.6,4.3) {rayons\\ verticaux};
\end{tikzpicture}
\end{center}
\legende{La projection verticale $f$ : l'ombre $A'$ du point $A$ sur le sol s'obtient en « écrasant » la hauteur $z$.}
\end{popfigure}

\courssub{Consignes}

\textbf{Partie A : étude de la projection.}

\begin{enumerate}
\item Montrer que l'application $f$ est une application linéaire de $\mathbb{R}^3$ dans $\mathbb{R}^3$.
\item Déterminer les images par $f$ des vecteurs $\vec{i}$, $\vec{j}$, $\vec{k}$ de la base canonique, puis écrire la matrice $M$ de $f$ dans la base canonique.
\item Déterminer le noyau $\ker f$. En donner une base et sa dimension. Interpréter géométriquement : quels sont les points de la salle dont l'ombre tombe exactement au point $O$ ?
\item Déterminer l'image $\operatorname{Im} f$. En donner une base et sa dimension. Interpréter géométriquement : quelle est la forme de l'ombre du mobile sur le sol ?
\item Vérifier que le théorème du rang est satisfait pour $f$.
\end{enumerate}

\textbf{Partie B : changement de base.}

\begin{enumerate}
\setcounter{enumi}{5}
\item Montrer que la famille $\mathcal{B}' = (\vec{u}, \vec{v}, \vec{w})$ est une base de $\mathbb{R}^3$.
\item Déterminer les coordonnées du vecteur $\overrightarrow{OA}$ dans la base $\mathcal{B}'$.
\item Calculer $f(\vec{u})$, $f(\vec{v})$, $f(\vec{w})$. Les ombres des trois directions de poutres forment-elles une famille libre du plan du sol ? Commenter.
\end{enumerate}

\courssub{Ressources à mobiliser}

\begin{aretenir}[Boîte à outils de la leçon]
\begin{itemize}
\item \textbf{Linéarité :} $f$ est linéaire si pour tous vecteurs $\vec{a}, \vec{b}$ et tout réel $\lambda$ : $f(\vec{a}+\vec{b}) = f(\vec{a}) + f(\vec{b})$ et $f(\lambda \vec{a}) = \lambda f(\vec{a})$.
\item \textbf{Matrice dans une base :} les colonnes de la matrice de $f$ dans la base canonique sont les coordonnées de $f(\vec{i})$, $f(\vec{j})$, $f(\vec{k})$.
\item \textbf{Noyau :} $\ker f = \{\vec{a} \in \mathbb{R}^3 \mid f(\vec{a}) = \vec{0}\}$ ; c'est un sous-espace vectoriel de $\mathbb{R}^3$.
\item \textbf{Image :} $\operatorname{Im} f = \{f(\vec{a}) \mid \vec{a} \in \mathbb{R}^3\} = \operatorname{Vect}\big(f(\vec{i}), f(\vec{j}), f(\vec{k})\big)$ ; c'est aussi un sous-espace vectoriel de $\mathbb{R}^3$.
\item \textbf{Théorème du rang :} $\dim \mathbb{R}^3 = \dim \ker f + \dim \operatorname{Im} f$.
\item \textbf{Base de $\mathbb{R}^3$ :} trois vecteurs forment une base si et seulement s'ils sont libres, c'est-à-dire si l'égalité $\alpha \vec{u} + \beta \vec{v} + \gamma \vec{w} = \vec{0}$ impose $\alpha = \beta = \gamma = 0$.
\item \textbf{Coordonnées :} les coordonnées de $\vec{a}$ dans $(\vec{u},\vec{v},\vec{w})$ sont les réels $(\alpha,\beta,\gamma)$ tels que $\vec{a} = \alpha \vec{u} + \beta \vec{v} + \gamma \vec{w}$, obtenus en résolvant un système de trois équations à trois inconnues.
\end{itemize}
\end{aretenir}

\courssub{Démarche et corrigé commenté}

\begin{exoresolu}[Le décorateur d'intérieur et la projection des ombres]
On considère $f : \mathbb{R}^3 \to \mathbb{R}^3$ définie par $f(x,y,z) = (x,y,0)$, les vecteurs $\vec{u} = (1,1,0)$, $\vec{v} = (0,1,1)$, $\vec{w} = (1,0,1)$ et le point $A(3\,;2\,;2)$.
\tcblower
\textbf{1. Linéarité de $f$.} Soient $\vec{a} = (x,y,z)$, $\vec{b} = (x',y',z')$ deux vecteurs quelconques de $\mathbb{R}^3$ et $\lambda \in \mathbb{R}$.
\begin{align*}
f(\vec{a} + \vec{b}) &= f(x+x',\, y+y',\, z+z') = (x+x',\, y+y',\, 0) \\
&= (x,y,0) + (x',y',0) = f(\vec{a}) + f(\vec{b}).
\end{align*}
De même :
\[
f(\lambda \vec{a}) = f(\lambda x, \lambda y, \lambda z) = (\lambda x, \lambda y, 0) = \lambda (x,y,0) = \lambda f(\vec{a}).
\]
Les deux conditions sont satisfaites pour tous vecteurs et tout réel $\lambda$ : $f$ est donc une \cle{application linéaire} de $\mathbb{R}^3$ dans $\mathbb{R}^3$ (c'est un \emph{endomorphisme} de $\mathbb{R}^3$).

\textbf{2. Matrice de $f$ dans la base canonique.} On calcule les images des vecteurs de base :
\[
f(\vec{i}) = f(1,0,0) = (1,0,0) = \vec{i}, \quad f(\vec{j}) = f(0,1,0) = (0,1,0) = \vec{j}, \quad f(\vec{k}) = f(0,0,1) = (0,0,0) = \vec{0}.
\]
En plaçant ces coordonnées \emph{en colonnes}, on obtient :
\[
M = \begin{pmatrix} 1 & 0 & 0 \\ 0 & 1 & 0 \\ 0 & 0 & 0 \end{pmatrix}.
\]
Contrôle : le produit matriciel redonne bien la définition de $f$ :
\[
M \begin{pmatrix} x \\ y \\ z \end{pmatrix} = \begin{pmatrix} x \\ y \\ 0 \end{pmatrix} = f(x,y,z).
\]

\textbf{3. Noyau de $f$.} On résout $f(x,y,z) = (0,0,0)$, c'est-à-dire $(x,y,0) = (0,0,0)$, soit $x = 0$ et $y = 0$, \emph{sans condition sur $z$}. Donc :
\[
\ker f = \{(0,0,z) \mid z \in \mathbb{R}\} = \operatorname{Vect}(\vec{k}).
\]
Le noyau est la droite vectorielle dirigée par $\vec{k}$ : une base de $\ker f$ est $(\vec{k})$ et $\dim \ker f = 1$.

\emph{Interprétation :} les points dont l'ombre tombe exactement au pied du mur $O$ sont tous les points de la verticale passant par $O$, c'est-à-dire la colonne d'air située juste sous la lampe. Tout objet suspendu sur cette verticale, quelle que soit sa hauteur, a son ombre en $O$.

\textbf{4. Image de $f$.} On a $\operatorname{Im} f = \operatorname{Vect}\big(f(\vec{i}), f(\vec{j}), f(\vec{k})\big) = \operatorname{Vect}(\vec{i}, \vec{j}, \vec{0}) = \operatorname{Vect}(\vec{i}, \vec{j})$. Or la famille $(\vec{i}, \vec{j})$ est libre, donc c'est une base de $\operatorname{Im} f$ et :
\[
\operatorname{Im} f = \{(x,y,0) \mid x, y \in \mathbb{R}\}, \qquad \dim \operatorname{Im} f = 2.
\]
\emph{Interprétation :} l'image est le plan vectoriel d'équation $z = 0$, c'est-à-dire le sol. Toute ombre est un dessin plan sur le carrelage : l'application $f$ « écrase » l'espace de dimension 3 sur un plan de dimension 2.

\textbf{5. Vérification du théorème du rang.}
\[
\dim \ker f + \dim \operatorname{Im} f = 1 + 2 = 3 = \dim \mathbb{R}^3.
\]
Le théorème du rang est bien vérifié. La dimension « perdue » par la projection (la hauteur) correspond exactement au noyau.

\textbf{6. La famille $\mathcal{B}' = (\vec{u},\vec{v},\vec{w})$ est une base.} Soient $\alpha, \beta, \gamma \in \mathbb{R}$ tels que $\alpha \vec{u} + \beta \vec{v} + \gamma \vec{w} = \vec{0}$. En coordonnées :
\[
\alpha(1,1,0) + \beta(0,1,1) + \gamma(1,0,1) = (0,0,0)
\iff
\begin{cases} \alpha + \gamma = 0 \\ \alpha + \beta = 0 \\ \beta + \gamma = 0 \end{cases}
\]
La première équation donne $\gamma = -\alpha$, la deuxième $\beta = -\alpha$ ; en reportant dans la troisième : $-\alpha + (-\alpha) = 0$, soit $-2\alpha = 0$, d'où $\alpha = 0$, puis $\beta = \gamma = 0$. La famille est donc \cle{libre}. Comme elle comporte $3$ vecteurs dans $\mathbb{R}^3$ qui est de dimension $3$, c'est une \cle{base} de $\mathbb{R}^3$.

\textbf{7. Coordonnées de $\overrightarrow{OA}$ dans $\mathcal{B}'$.} On cherche $(\alpha, \beta, \gamma)$ tels que $\overrightarrow{OA} = \alpha \vec{u} + \beta \vec{v} + \gamma \vec{w}$, c'est-à-dire :
\[
\begin{cases} \alpha + \gamma = 3 \\ \alpha + \beta = 2 \\ \beta + \gamma = 2 \end{cases}
\]
De la première équation : $\gamma = 3 - \alpha$ ; de la deuxième : $\beta = 2 - \alpha$. En reportant dans la troisième :
\[
(2 - \alpha) + (3 - \alpha) = 2 \iff 5 - 2\alpha = 2 \iff \alpha = \dfrac{3}{2}.
\]
Alors $\beta = 2 - \dfrac{3}{2} = \dfrac{1}{2}$ et $\gamma = 3 - \dfrac{3}{2} = \dfrac{3}{2}$. Vérification :
\[
\dfrac{3}{2}(1,1,0) + \dfrac{1}{2}(0,1,1) + \dfrac{3}{2}(1,0,1) = \left(\dfrac{3}{2} + \dfrac{3}{2},\, \dfrac{3}{2} + \dfrac{1}{2},\, \dfrac{1}{2} + \dfrac{3}{2}\right) = (3, 2, 2).
\]
Les coordonnées de $\overrightarrow{OA}$ dans la base $\mathcal{B}'$ sont donc :
\[
\left(\dfrac{3}{2}\,;\, \dfrac{1}{2}\,;\, \dfrac{3}{2}\right).
\]
Monsieur Kamdem peut ainsi régler ses fils selon les trois directions de poutres : par exemple, en prenant le mètre comme unité le long de chaque poutre, les trois « quotas » de décomposition sont $1{,}5$ m, $0{,}5$ m et $1{,}5$ m.

\textbf{8. Ombres des directions de poutres.}
\[
f(\vec{u}) = (1,1,0), \qquad f(\vec{v}) = (0,1,0), \qquad f(\vec{w}) = (1,0,0).
\]
Ces trois vecteurs appartiennent tous au plan $\operatorname{Im} f$, qui est de dimension $2$ : ils ne peuvent donc pas former une famille libre. Explicitement :
\[
f(\vec{v}) + f(\vec{w}) = (0,1,0) + (1,0,0) = (1,1,0) = f(\vec{u}),
\]
c'est-à-dire $f(\vec{u}) - f(\vec{v}) - f(\vec{w}) = \vec{0}$ : combinaison linéaire nulle à coefficients non tous nuls. La famille des ombres est \cle{liée}.

C'est cohérent avec la linéarité : $\vec{u} - \vec{v} - \vec{w} = (1-0-1,\, 1-1-0,\, 0-1-1) = (0,0,-2) = -2\vec{k} \in \ker f$, donc $f(\vec{u} - \vec{v} - \vec{w}) = \vec{0}$, c'est-à-dire $f(\vec{u}) = f(\vec{v}) + f(\vec{w})$. La projection a « écrasé » la composante verticale qui distinguait les trois poutres.
\end{exoresolu}

\begin{attention}[Pièges à éviter dans cette activité]
\begin{itemize}
\item Ne pas confondre $\ker f$ (ensemble des vecteurs \emph{de départ} envoyés sur $\vec{0}$) avec un plan : ici le noyau est une \emph{droite} vectorielle, la verticale dirigée par $\vec{k}$.
\item La matrice $M$ n'est pas inversible : son déterminant est nul, ce qui est normal puisque $\ker f \neq \{\vec{0}\}$ (une application linéaire est injective si et seulement si son noyau est réduit au vecteur nul).
\item Une famille libre de $\mathbb{R}^3$ peut avoir une image \emph{liée} par une application linéaire non injective : la linéarité ne conserve la liberté que si $f$ est injective.
\item Dans la résolution des systèmes des questions 6 et 7, ne pas oublier de \emph{vérifier} la solution obtenue en la reportant dans les trois équations : c'est la meilleure protection contre les erreurs de signe.
\end{itemize}
\end{attention}

\courssub{Bilan de l'activité}

Cette activité a permis de mettre en évidence les points essentiels de la leçon :

\begin{itemize}
\item une \textbf{projection} est le prototype géométrique d'une application linéaire non injective et non surjective : elle « perd de l'information » (la hauteur $z$) ;
\item le \textbf{noyau} mesure précisément cette perte : c'est l'ensemble des directions que la projection écrase sur $\vec{0}$ ; l'\textbf{image} est le « support » du résultat, ici le plan du sol ;
\item le \textbf{théorème du rang} traduit un bilan de dimensions : ce qui est écrasé ($\dim \ker f = 1$) et ce qui subsiste ($\dim \operatorname{Im} f = 2$) se recomposent en la dimension totale de l'espace ($3$) ;
\item la \textbf{matrice} d'une application linéaire se lit directement sur les images des vecteurs de base, et le calcul matriciel remplace le calcul vectoriel ;
\item le \textbf{changement de base} n'est pas un exercice artificiel : il répond à un besoin concret de l'artisan (décomposer une position selon les directions réelles de ses poutres), et il se ramène toujours à la résolution d'un système linéaire ;
\item enfin, une application linéaire non injective peut transformer une base en famille liée : la liberté n'est préservée que lorsque le noyau est réduit au vecteur nul.
\end{itemize}

Tu as ainsi mobilisé, sur une même situation authentique, la quasi-totalité des savoir-faire exigibles au Baccalauréat sur les espaces vectoriels et les applications linéaires en dimension 3.

\newpage

\coursec{Méthodes et exercices résolus}

\courssub{Montrer qu'une famille de vecteurs est une base de $\mathbb{R}^3$}

\begin{methode}[Prouver qu'une famille est une base de $\mathbb{R}^3$]
Pour montrer qu'une famille $(\vec{u},\vec{v},\vec{w})$ de trois vecteurs de $\mathbb{R}^3$ est une base, il suffit de montrer qu'elle est \cle{libre} (car toute famille libre de 3 vecteurs dans un espace de dimension 3 est une base).
\begin{enumerate}
\item Écrire l'équation aux scalaires : soient $\alpha,\beta,\gamma\in\mathbb{R}$ tels que $\alpha\vec{u}+\beta\vec{v}+\gamma\vec{w}=\vec{0}$.
\item Traduire cette égalité vectorielle en un \textbf{système de 3 équations à 3 inconnues} $\alpha,\beta,\gamma$ (une équation par coordonnée).
\item Résoudre le système (substitution ou pivot de Gauss).
\item Conclure : si l'unique solution est $\alpha=\beta=\gamma=0$, la famille est libre ; comme elle comporte 3 vecteurs dans $\mathbb{R}^3$ (de dimension 3), c'est une \textbf{base}.
\end{enumerate}
\textbf{Réflexe Bac :} on peut aussi montrer que la famille est \emph{génératrice} de $\mathbb{R}^3$ (tout vecteur se décompose), mais la liberté est en général plus rapide. Une famille de 3 vecteurs liée (par exemple avec un vecteur nul ou deux vecteurs colinéaires) n'est jamais une base.
\end{methode}

\begin{exoresolu}[Type Bac — Liberté d'une famille]
Dans l'espace muni de la base canonique $(\vec{i},\vec{j},\vec{k})$, on considère les vecteurs
\[
\vec{u}=(1,1,0),\qquad \vec{v}=(0,1,1),\qquad \vec{w}=(1,0,1).
\]
Montrer que $(\vec{u},\vec{v},\vec{w})$ est une base de $\mathbb{R}^3$.
\tcblower
\textbf{Étape 1 : équation aux scalaires.} Soient $\alpha,\beta,\gamma$ trois réels tels que
\[
\alpha\vec{u}+\beta\vec{v}+\gamma\vec{w}=\vec{0}.
\]
\textbf{Étape 2 : traduction en système.} En calculant coordonnée par coordonnée :
\[
\alpha(1,1,0)+\beta(0,1,1)+\gamma(1,0,1)=(\alpha+\gamma,\ \alpha+\beta,\ \beta+\gamma)=(0,0,0),
\]
d'où le système
\[
\begin{cases}
\alpha+\gamma=0 & (L_1)\\
\alpha+\beta=0 & (L_2)\\
\beta+\gamma=0 & (L_3)
\end{cases}
\]
\textbf{Étape 3 : résolution.} De $(L_1)$ : $\gamma=-\alpha$. De $(L_2)$ : $\beta=-\alpha$. En reportant dans $(L_3)$ :
\[
\beta+\gamma=-\alpha-\alpha=-2\alpha=0 \quad\Longrightarrow\quad \alpha=0,
\]
puis $\beta=0$ et $\gamma=0$. L'unique solution est $(\alpha,\beta,\gamma)=(0,0,0)$.
\textbf{Étape 4 : conclusion.} La famille $(\vec{u},\vec{v},\vec{w})$ est \textbf{libre}. Or elle comporte $3$ vecteurs et $\dim\mathbb{R}^3=3$ : toute famille libre de 3 vecteurs de $\mathbb{R}^3$ est une base. Donc $(\vec{u},\vec{v},\vec{w})$ est une \textbf{base de $\mathbb{R}^3$}.
\end{exoresolu}

\courssub{Déterminer les coordonnées d'un vecteur dans une base}

\begin{methode}[Coordonnées d'un vecteur dans une base de $\mathbb{R}^3$]
Soit $\mathcal{B}=(\vec{u},\vec{v},\vec{w})$ une base de $\mathbb{R}^3$ et $\vec{a}=(a_1,a_2,a_3)$ un vecteur donné dans la base canonique. Pour trouver les coordonnées $(x,y,z)$ de $\vec{a}$ dans $\mathcal{B}$ :
\begin{enumerate}
\item Poser la décomposition $\vec{a}=x\vec{u}+y\vec{v}+z\vec{w}$.
\item Identifier les coordonnées dans la base canonique : on obtient un système linéaire de 3 équations d'inconnues $x,y,z$.
\item Résoudre le système par substitution ou combinaisons linéaires.
\item Vérifier le résultat en recalculant $x\vec{u}+y\vec{v}+z\vec{w}$, puis conclure : $\vec{a}$ a pour coordonnées $(x,y,z)$ dans $\mathcal{B}$, noté $\vec{a}\,(x,y,z)_{\mathcal{B}}$.
\end{enumerate}
\textbf{À retenir :} les coordonnées dépendent de la base choisie ; un même vecteur a des coordonnées différentes dans des bases différentes.
\end{methode}

\begin{exoresolu}[Type Bac — Changement de base]
On reprend la base $\mathcal{B}=(\vec{u},\vec{v},\vec{w})$ avec $\vec{u}=(1,1,0)$, $\vec{v}=(0,1,1)$, $\vec{w}=(1,0,1)$. Déterminer les coordonnées du vecteur $\vec{a}=(2,3,1)$ dans la base $\mathcal{B}$.
\tcblower
\textbf{Étape 1 :} on cherche $(x,y,z)$ réels tels que $\vec{a}=x\vec{u}+y\vec{v}+z\vec{w}$.
\textbf{Étape 2 :} en coordonnées canoniques :
\[
x(1,1,0)+y(0,1,1)+z(1,0,1)=(x+z,\ x+y,\ y+z)=(2,3,1),
\]
soit le système
\[
\begin{cases}
x+z=2 & (L_1)\\
x+y=3 & (L_2)\\
y+z=1 & (L_3)
\end{cases}
\]
\textbf{Étape 3 : résolution.} De $(L_1)$ : $z=2-x$. De $(L_2)$ : $y=3-x$. En reportant dans $(L_3)$ :
\[
(3-x)+(2-x)=1 \;\Longleftrightarrow\; 5-2x=1 \;\Longleftrightarrow\; x=2.
\]
Alors $y=3-2=1$ et $z=2-2=0$.
\textbf{Étape 4 : vérification.} $2\vec{u}+1\vec{v}+0\vec{w}=(2,2,0)+(0,1,1)=(2,3,1)=\vec{a}$. \checkmark
\textbf{Conclusion :} $\vec{a}$ a pour coordonnées $(2,1,0)$ dans la base $\mathcal{B}$, c'est-à-dire $\vec{a}=2\vec{u}+\vec{v}$.
\end{exoresolu}

\courssub{Reconnaître une application linéaire}

\begin{methode}[Montrer qu'une application est linéaire (ou ne l'est pas)]
Pour montrer que $f:\mathbb{R}^3\to\mathbb{R}^3$ est \cle{linéaire}, deux stratégies :
\begin{enumerate}
\item \textbf{Stratégie directe (définition) :} montrer que pour tous vecteurs $\vec{x},\vec{y}$ et tout réel $\lambda$ :
\[
f(\vec{x}+\vec{y})=f(\vec{x})+f(\vec{y})\quad\text{et}\quad f(\lambda\vec{x})=\lambda f(\vec{x}).
\]
\item \textbf{Stratégie par la formule :} si $f(x,y,z)$ est donné par des expressions \emph{linéaires} en $x,y,z$ (somme de termes de degré 1, sans constante), alors $f$ est linéaire — mais au Bac, il faut le \emph{démontrer} par la définition ou le justifier proprement.
\end{enumerate}
Pour montrer que $f$ \textbf{n'est pas} linéaire, il suffit d'un contre-exemple : le plus rapide est souvent de vérifier $f(\vec{0})$. Si $f(\vec{0})\neq\vec{0}$, $f$ n'est pas linéaire. Attention : $f(\vec{0})=\vec{0}$ ne suffit \emph{pas} à prouver la linéarité !
\end{methode}

\begin{exoresolu}[Type Bac — Linéaire ou pas ?]
On considère les applications de $\mathbb{R}^3$ dans $\mathbb{R}^3$ définies par :
\[
f(x,y,z)=(x+2y-z,\; y+z,\; x),\qquad g(x,y,z)=(x+y,\; z^2,\; 1).
\]
\begin{enumerate}
\item Montrer que $f$ est une application linéaire.
\item Montrer que $g$ n'est pas une application linéaire.
\end{enumerate}
\tcblower
\textbf{1. Linéarité de $f$.} Soient $\vec{x}=(x,y,z)$, $\vec{x}'=(x',y',z')$ deux vecteurs de $\mathbb{R}^3$ et $\lambda\in\mathbb{R}$.
\emph{Additivité :} $\vec{x}+\vec{x}'=(x+x',\,y+y',\,z+z')$, donc
\begin{align*}
f(\vec{x}+\vec{x}')&=\big((x+x')+2(y+y')-(z+z'),\;(y+y')+(z+z'),\;x+x'\big)\\
&=(x+2y-z,\;y+z,\;x)+(x'+2y'-z',\;y'+z',\;x')\\
&=f(\vec{x})+f(\vec{x}').
\end{align*}
\emph{Homogénéité :} $\lambda\vec{x}=(\lambda x,\lambda y,\lambda z)$, donc
\[
f(\lambda\vec{x})=(\lambda x+2\lambda y-\lambda z,\;\lambda y+\lambda z,\;\lambda x)=\lambda(x+2y-z,\;y+z,\;x)=\lambda f(\vec{x}).
\]
$f$ vérifie les deux axiomes : $f$ est une \textbf{application linéaire} (endomorphisme de $\mathbb{R}^3$).

\textbf{2. Non-linéarité de $g$.} Calculons $g(\vec{0})=g(0,0,0)=(0+0,\;0^2,\;1)=(0,0,1)\neq\vec{0}$.
Or toute application linéaire vérifie $g(\vec{0})=\vec{0}$. Donc $g$ \textbf{n'est pas linéaire}.
\emph{(On pouvait aussi remarquer le terme $z^2$, non linéaire : $g(0,0,2)=(0,4,1)$ alors que $2\,g(0,0,1)=(0,2,2)$ : l'homogénéité est mise en défaut.)}
\end{exoresolu}

\courssub{Écrire la matrice d'une application linéaire dans une base}

\begin{methode}[Construire la matrice de $f$ dans une base]
Soit $f:\mathbb{R}^3\to\mathbb{R}^3$ linéaire et $\mathcal{B}=(\vec{e_1},\vec{e_2},\vec{e_3})$ une base. La matrice $M$ de $f$ dans $\mathcal{B}$ s'obtient ainsi :
\begin{enumerate}
\item Calculer les images $f(\vec{e_1})$, $f(\vec{e_2})$, $f(\vec{e_3})$.
\item Exprimer chaque image \textbf{dans la base $\mathcal{B}$} (dans la base canonique, les coordonnées se lisent directement).
\item Placer les coordonnées de $f(\vec{e_j})$ en \textbf{colonne $j$} de la matrice :
\[
M=\begin{pmatrix} | & | & | \\ f(\vec{e_1}) & f(\vec{e_2}) & f(\vec{e_3}) \\ | & | & | \end{pmatrix}.
\]
\item Retenir la règle d'usage : si $\vec{x}$ a pour coordonnées la colonne $X$ dans $\mathcal{B}$, alors $f(\vec{x})$ a pour coordonnées $MX$ dans $\mathcal{B}$.
\end{enumerate}
\textbf{Piège :} les images s'écrivent en \emph{colonnes}, jamais en lignes.
\end{methode}

\begin{exoresolu}[Type Bac — Matrice dans la base canonique]
Soit $f:\mathbb{R}^3\to\mathbb{R}^3$ l'application linéaire définie par
\[
f(x,y,z)=(x+2y-z,\;y+z,\;x).
\]
\begin{enumerate}
\item Déterminer la matrice $A$ de $f$ dans la base canonique $\mathcal{B}_0=(\vec{i},\vec{j},\vec{k})$.
\item En utilisant $A$, calculer l'image du vecteur $\vec{a}=(1,-1,2)$.
\end{enumerate}
\tcblower
\textbf{1. Construction de $A$.} On calcule les images des vecteurs de base :
\[
f(\vec{i})=f(1,0,0)=(1,0,1),\quad f(\vec{j})=f(0,1,0)=(2,1,0),\quad f(\vec{k})=f(0,0,1)=(-1,1,0).
\]
On place ces vecteurs en colonnes :
\[
A=\begin{pmatrix} 1 & 2 & -1 \\ 0 & 1 & 1 \\ 1 & 0 & 0 \end{pmatrix}.
\]
\textbf{2. Image de $\vec{a}$.} Les coordonnées de $f(\vec{a})$ dans $\mathcal{B}_0$ sont données par le produit $AX$ avec $X=\begin{pmatrix}1\\-1\\2\end{pmatrix}$ :
\[
AX=\begin{pmatrix} 1 & 2 & -1 \\ 0 & 1 & 1 \\ 1 & 0 & 0 \end{pmatrix}\begin{pmatrix}1\\-1\\2\end{pmatrix}
=\begin{pmatrix} 1\times1+2\times(-1)+(-1)\times2 \\ 0\times1+1\times(-1)+1\times2 \\ 1\times1+0\times(-1)+0\times2 \end{pmatrix}
=\begin{pmatrix} -3 \\ 1 \\ 1 \end{pmatrix}.
\]
\emph{Vérification directe :} $f(1,-1,2)=(1+2(-1)-2,\;-1+2,\;1)=(-3,1,1)$. \checkmark
\textbf{Conclusion :} $f(\vec{a})=(-3,1,1)$.
\end{exoresolu}

\courssub{Déterminer le noyau et l'image d'une application linéaire}

\begin{methode}[Noyau et image : méthodes de calcul]
\textbf{Noyau.} $\ker f=\{\vec{x}\in\mathbb{R}^3 \mid f(\vec{x})=\vec{0}\}$. Pour le déterminer :
\begin{enumerate}
\item Traduire $f(x,y,z)=(0,0,0)$ en un système homogène de 3 équations.
\item Résoudre le système ; exprimer les solutions à l'aide de paramètres.
\item Écrire $\ker f$ comme ensemble des combinaisons linéaires des vecteurs obtenus : $\ker f=\mathrm{Vect}(\vec{u_1},\dots)$, et donner une base et la dimension.
\end{enumerate}
\textbf{Image.} $\mathrm{Im}\,f=\{f(\vec{x})\mid \vec{x}\in\mathbb{R}^3\}$. Deux présentations :
\begin{enumerate}
\item $\mathrm{Im}\,f=\mathrm{Vect}\big(f(\vec{e_1}),f(\vec{e_2}),f(\vec{e_3})\big)$ : l'image est engendrée par les colonnes de la matrice.
\item Extraire de cette famille génératrice une sous-famille \textbf{libre} : c'est une base de $\mathrm{Im}\,f$ ; son nombre d'éléments est le \cle{rang} de $f$.
\end{enumerate}
\textbf{Réflexes :} $f$ injective $\Leftrightarrow \ker f=\{\vec{0}\}$ ; $f$ surjective $\Leftrightarrow \mathrm{Im}\,f=\mathbb{R}^3$ ; pour un endomorphisme de $\mathbb{R}^3$, $f$ bijective $\Leftrightarrow$ injective $\Leftrightarrow$ surjective.
\end{methode}

\begin{exoresolu}[Type Bac — Noyau et image]
Soit $g:\mathbb{R}^3\to\mathbb{R}^3$ l'application linéaire définie par
\[
g(x,y,z)=(x+y+z,\;x-y,\;2x+z).
\]
\begin{enumerate}
\item Déterminer $\ker g$ ; en donner une base et sa dimension. $g$ est-elle injective ?
\item Déterminer une base de $\mathrm{Im}\,g$ et le rang de $g$. $g$ est-elle surjective ?
\end{enumerate}
\tcblower
\textbf{1. Noyau.} $(x,y,z)\in\ker g \Leftrightarrow g(x,y,z)=(0,0,0)$, soit :
\[
\begin{cases}
x+y+z=0 & (L_1)\\
x-y=0 & (L_2)\\
2x+z=0 & (L_3)
\end{cases}
\]
De $(L_2)$ : $y=x$. De $(L_3)$ : $z=-2x$. Vérifions dans $(L_1)$ : $x+x-2x=0$ : toujours vrai. Les solutions sont donc les triplets
\[
(x,y,z)=(x,\;x,\;-2x)=x\,(1,1,-2),\qquad x\in\mathbb{R}.
\]
Ainsi $\ker g=\mathrm{Vect}(\vec{u})$ avec $\vec{u}=(1,1,-2)$. Le vecteur $\vec{u}$ étant non nul, il forme une famille libre : $(\vec{u})$ est une \textbf{base de $\ker g$} et $\dim\ker g=1$.
Comme $\ker g\neq\{\vec{0}\}$, $g$ \textbf{n'est pas injective}.

\textbf{2. Image.} L'image est engendrée par les images des vecteurs de la base canonique :
\[
g(\vec{i})=(1,1,2),\quad g(\vec{j})=(1,-1,0),\quad g(\vec{k})=(1,0,1).
\]
Donc $\mathrm{Im}\,g=\mathrm{Vect}\big((1,1,2),(1,-1,0),(1,0,1)\big)$. Cherchons une sous-famille libre.
Les vecteurs $\vec{a}=(1,1,2)$ et $\vec{b}=(1,-1,0)$ ne sont pas colinéaires (leurs coordonnées ne sont pas proportionnelles : $\frac{1}{1}=1$ mais $\frac{1}{-1}=-1$), donc $(\vec{a},\vec{b})$ est libre.
De plus, remarquons que $(1,0,1)=\tfrac{1}{2}\big[(1,1,2)+(1,-1,0)\big]=\tfrac12\vec{a}+\tfrac12\vec{b}$ : en effet $\tfrac12(1+1,\,1-1,\,2+0)=(1,0,1)$. Le troisième vecteur est donc combinaison linéaire des deux premiers.
Par conséquent $\mathrm{Im}\,g=\mathrm{Vect}(\vec{a},\vec{b})$, la famille $(\vec{a},\vec{b})$ est une \textbf{base de $\mathrm{Im}\,g$}, et le rang de $g$ vaut $\mathrm{rg}(g)=\dim\mathrm{Im}\,g=2$.
Comme $\mathrm{Im}\,g\neq\mathbb{R}^3$ (dimension 2 au lieu de 3), $g$ \textbf{n'est pas surjective}.
\end{exoresolu}

\courssub{Utiliser le théorème du rang}

\begin{methode}[Exploiter $\dim E=\dim\ker f+\dim\mathrm{Im}\,f$]
Le \cle{théorème du rang} relie les dimensions : pour $f:\mathbb{R}^3\to\mathbb{R}^3$ linéaire,
\[
\dim\mathbb{R}^3=\dim\ker f+\dim\mathrm{Im}\,f,\qquad\text{soit}\qquad 3=\dim\ker f+\mathrm{rg}(f).
\]
Mode d'emploi :
\begin{enumerate}
\item Calculer \textbf{une seule} des deux dimensions (celle qui est la plus facile : souvent le rang via la matrice, ou le noyau si le système est simple).
\item En déduire l'autre par la formule, \emph{sans refaire de calcul}.
\item Conclure sur la bijectivité : $f$ bijective $\Leftrightarrow \ker f=\{\vec{0}\} \Leftrightarrow \mathrm{rg}(f)=3$.
\end{enumerate}
\textbf{Réflexe Bac :} si l'énoncé donne $\dim\ker f$, le rang s'obtient immédiatement ; réciproquement. Ne jamais additionner les dimensions sans vérifier que la somme vaut 3.
\end{methode}

\begin{exoresolu}[Type Bac — Théorème du rang]
Soit $h:\mathbb{R}^3\to\mathbb{R}^3$ l'application linéaire dont la matrice dans la base canonique est
\[
B=\begin{pmatrix} 1 & -1 & 2 \\ 2 & 1 & 1 \\ 1 & 2 & -1 \end{pmatrix}.
\]
\begin{enumerate}
\item Montrer que les vecteurs colonnes $\vec{c_1}=(1,2,1)$ et $\vec{c_2}=(-1,1,2)$ sont libres, puis que $\vec{c_3}=(2,1,-1)$ est combinaison linéaire de $\vec{c_1}$ et $\vec{c_2}$.
\item En déduire le rang de $h$, puis, sans aucun calcul supplémentaire, la dimension de $\ker h$.
\item Donner une base de $\ker h$.
\end{enumerate}
\tcblower
\textbf{1.} \emph{Liberté de $(\vec{c_1},\vec{c_2})$ :} ces deux vecteurs ne sont pas colinéaires, car leurs coordonnées ne sont pas proportionnelles : $\frac{1}{-1}=-1$ alors que $\frac{2}{1}=2$. Deux vecteurs non colinéaires forment une famille libre.

\emph{Combinaison linéaire :} cherchons $\alpha,\beta$ tels que $\vec{c_3}=\alpha\vec{c_1}+\beta\vec{c_2}$ :
\[
\begin{cases}
\alpha-\beta=2\\
2\alpha+\beta=1\\
\alpha+2\beta=-1
\end{cases}
\]
En additionnant les deux premières équations : $3\alpha=3$, donc $\alpha=1$, puis $\beta=\alpha-2=-1$. Vérification dans la troisième : $1+2(-1)=-1$. \checkmark
Donc $\vec{c_3}=\vec{c_1}-\vec{c_2}$.

\textbf{2. Rang de $h$.} On sait que $\mathrm{Im}\,h=\mathrm{Vect}(\vec{c_1},\vec{c_2},\vec{c_3})$. Comme $\vec{c_3}$ est combinaison linéaire de $\vec{c_1}$ et $\vec{c_2}$, on a $\mathrm{Im}\,h=\mathrm{Vect}(\vec{c_1},\vec{c_2})$. La famille $(\vec{c_1},\vec{c_2})$ est génératrice de $\mathrm{Im}\,h$ et libre : c'est une base de $\mathrm{Im}\,h$, donc
\[
\mathrm{rg}(h)=\dim\mathrm{Im}\,h=2.
\]
\emph{Dimension du noyau par le théorème du rang :}
\[
\dim\mathbb{R}^3=\dim\ker h+\mathrm{rg}(h)\;\Longrightarrow\; \dim\ker h=3-2=1.
\]

\textbf{3. Base de $\ker h$.} La relation $\vec{c_3}=\vec{c_1}-\vec{c_2}$ signifie $h(\vec{k})=h(\vec{i})-h(\vec{j})$, c'est-à-dire, par linéarité, $h(\vec{i}-\vec{j}-\vec{k})=\vec{0}$. Le vecteur $\vec{v}=(1,-1,-1)$ appartient donc à $\ker h$.
Vérification directe : $B\begin{pmatrix}1\\-1\\-1\end{pmatrix}=\begin{pmatrix}1+1-2\\2-1-1\\1-2+1\end{pmatrix}=\begin{pmatrix}0\\0\\0\end{pmatrix}$. \checkmark
Comme $\vec{v}\neq\vec{0}$ et $\dim\ker h=1$, la famille $(\vec{v})$ est une \textbf{base de $\ker h$} : $\ker h=\mathrm{Vect}\big((1,-1,-1)\big)$.
\end{exoresolu}

\courssub{Étude complète : bijectivité et application au contexte}

\begin{methode}[Conduire une étude complète d'un endomorphisme]
Face à une application linéaire $f$ de $\mathbb{R}^3$ donnée par sa matrice ou sa formule, le plan d'étude type Bac est :
\begin{enumerate}
\item Justifier la linéarité (souvent admise si $f$ est donnée par sa matrice).
\item Écrire la matrice $A$ dans la base canonique.
\item Résoudre $AX=0$ pour obtenir $\ker f$.
\item Appliquer le théorème du rang pour obtenir $\mathrm{rg}(f)$.
\item Conclure : si $\ker f=\{\vec{0}\}$ (ou $\mathrm{rg}(f)=3$), alors $f$ est un \textbf{isomorphisme} (bijective) ; sinon préciser base du noyau et base de l'image.
\end{enumerate}
\end{methode}

\begin{exoresolu}[Type Bac — Synthèse complète]
Une entreprise de télécommunications modélise un codage de signaux par l'application $f:\mathbb{R}^3\to\mathbb{R}^3$ définie par
\[
f(x,y,z)=(x-y+z,\;2x+y,\;x+2y-z).
\]
\begin{enumerate}
\item Déterminer $\ker f$ et en donner une base.
\item En déduire le rang de $f$ grâce au théorème du rang, puis une base de $\mathrm{Im}\,f$.
\item Le codage $f$ est-il inversible (c'est-à-dire $f$ est-elle bijective) ? Justifier.
\end{enumerate}
\tcblower
\textbf{1. Noyau.} $(x,y,z)\in\ker f\Leftrightarrow \begin{cases} x-y+z=0 & (L_1)\\ 2x+y=0 & (L_2)\\ x+2y-z=0 & (L_3)\end{cases}$

De $(L_2)$ : $y=-2x$. En reportant dans $(L_1)$ : $x-(-2x)+z=0$, donc $z=-3x$.
Vérification dans $(L_3)$ : $x+2(-2x)-(-3x)=x-4x+3x=0$ : toujours vérifié.
Les solutions sont $(x,y,z)=(x,-2x,-3x)=x(1,-2,-3)$, $x\in\mathbb{R}$.
Donc $\ker f=\mathrm{Vect}\big((1,-2,-3)\big)$, de base le vecteur $(1,-2,-3)$ (non nul, donc famille libre), et $\dim\ker f=1$.

\textbf{2. Rang et image.} Par le théorème du rang :
\[
\mathrm{rg}(f)=\dim\mathbb{R}^3-\dim\ker f=3-1=2.
\]
L'image est engendrée par les colonnes de la matrice de $f$ :
\[
f(\vec{i})=(1,2,1),\quad f(\vec{j})=(-1,1,2),\quad f(\vec{k})=(1,0,-1).
\]
Les deux premiers vecteurs ne sont pas colinéaires ($\frac{1}{-1}=-1\neq\frac{2}{1}=2$), donc la famille $\big((1,2,1),(-1,1,2)\big)$ est libre ; elle comporte $2=\mathrm{rg}(f)$ vecteurs de $\mathrm{Im}\,f$ : c'est une \textbf{base de $\mathrm{Im}\,f$}.
\emph{(Cohérence :} la relation $f(\vec{i})-2f(\vec{j})-3f(\vec{k})=\vec{0}$, qui traduit $(1,-2,-3)\in\ker f$, donne $f(\vec{k})=\tfrac13\big[f(\vec{i})-2f(\vec{j})\big]$ ; en effet $\tfrac13\big[(1,2,1)-2(-1,1,2)\big]=\tfrac13(3,0,-3)=(1,0,-1)$ : le troisième vecteur colonne est bien combinaison linéaire des deux premiers.\emph{)}

\textbf{3. Bijectivité.} $f$ est bijective si et seulement si $\ker f=\{\vec{0}\}$. Or ici $\dim\ker f=1$, donc $\ker f\neq\{\vec{0}\}$ : $f$ \textbf{n'est pas injective}, donc \textbf{pas bijective}. Le codage n'est pas inversible : deux signaux distincts différant d'un multiple de $(1,-2,-3)$ sont codés par le même message, l'information est perdue.
\end{exoresolu}

\begin{attention}[Les 5 erreurs qui coûtent des points au Bac]
\begin{enumerate}
\item \textbf{Conclure « base » sans compter les vecteurs.} Dire « la famille est libre donc c'est une base » n'est valable que parce qu'elle a \emph{exactement 3 vecteurs} dans $\mathbb{R}^3$ (dimension 3). Rédige toujours : « famille libre de 3 vecteurs dans un espace de dimension 3, donc base ».
\item \textbf{Croire que $f(\vec{0})=\vec{0}$ prouve la linéarité.} C'est une condition \emph{nécessaire} mais pas suffisante : l'application $(x,y,z)\mapsto(x^2,y,z)$ envoie $\vec{0}$ sur $\vec{0}$ sans être linéaire. En revanche, $f(\vec{0})\neq\vec{0}$ suffit à \emph{rejeter} la linéarité.
\item \textbf{Écrire les images en lignes dans la matrice.} Les coordonnées de $f(\vec{e_j})$ se placent en \emph{colonne} $j$. Une matrice transposée par erreur fausse tous les calculs de noyau et d'image qui suivent.
\item \textbf{Confondre noyau et image.} $\ker f$ est un sous-espace de l'espace de \emph{départ} (on résout $f(\vec{x})=\vec{0}$) ; $\mathrm{Im}\,f$ est un sous-espace de l'espace d'\emph{arrivée} (engendré par les colonnes). Ne jamais écrire « $\ker f=\mathrm{Vect}(f(\vec{i}),f(\vec{j}),f(\vec{k}))$ » : c'est l'image !
\item \textbf{Oublier le théorème du rang ou l'écrire faux.} La formule est $\dim E=\dim\ker f+\dim\mathrm{Im}\,f$, avec $\dim E=3$ ici — pas $\dim\ker f\times\dim\mathrm{Im}\,f$, pas $\dim\ker f-\mathrm{rg}(f)$. Et vérifie toujours la cohérence : $\dim\ker f+\mathrm{rg}(f)=3$.
\end{enumerate}
\end{attention}

\newpage

\coursec{Exercices}

\courssub{Je vérifie mes connaissances}

\begin{exercice}[QCM : une seule bonne réponse]
Pour chaque question, choisir la bonne réponse.
\begin{enumerate}
\item L'ensemble $F = \{(x, y, z) \in \mathbb{R}^3 \mid x + 2y - z = 0\}$ est :
\begin{enumerate}
\item un sous-espace vectoriel de $\mathbb{R}^3$ de dimension $1$ ;
\item un sous-espace vectoriel de $\mathbb{R}^3$ de dimension $2$ ;
\item pas un sous-espace vectoriel de $\mathbb{R}^3$.
\end{enumerate}
\item Une famille de $4$ vecteurs de $\mathbb{R}^3$ est :
\begin{enumerate}
\item toujours libre ; \item toujours liée ; \item parfois libre, parfois liée.
\end{enumerate}
\item Soit $f$ un endomorphisme de $\mathbb{R}^3$ tel que $\dim \ker f = 1$. Alors :
\begin{enumerate}
\item $\dim \operatorname{Im} f = 1$ ; \item $\dim \operatorname{Im} f = 2$ ; \item $f$ est bijective.
\end{enumerate}
\item La matrice de l'application linéaire $f(x, y, z) = (x + y, \, z, \, x - z)$ dans la base canonique est :
\begin{enumerate}
\item $\begin{pmatrix} 1 & 1 & 0 \\ 0 & 0 & 1 \\ 1 & 0 & -1 \end{pmatrix}$ ;
\item $\begin{pmatrix} 1 & 0 & 1 \\ 1 & 0 & 0 \\ 0 & 1 & -1 \end{pmatrix}$ ;
\item $\begin{pmatrix} 1 & 1 & 0 \\ 0 & 1 & 0 \\ 1 & 0 & -1 \end{pmatrix}$.
\end{enumerate}
\end{enumerate}
\end{exercice}

\begin{exercice}[Vrai ou faux, en justifiant]
Répondre par Vrai ou Faux en justifiant chaque réponse par un argument ou un contre-exemple.
\begin{enumerate}
\item Toute famille de $3$ vecteurs non nuls de $\mathbb{R}^3$ est une base de $\mathbb{R}^3$.
\item Si une famille de $3$ vecteurs de $\mathbb{R}^3$ est libre, alors c'est une base de $\mathbb{R}^3$.
\item L'ensemble $G = \{(x, y, z) \in \mathbb{R}^3 \mid x + y + z = 1\}$ est un sous-espace vectoriel de $\mathbb{R}^3$.
\item Si $f : \mathbb{R}^3 \to \mathbb{R}^3$ est linéaire et injective, alors $f$ est surjective.
\end{enumerate}
\end{exercice}

\begin{exercice}[Définitions à compléter]
Recopier et compléter les phrases suivantes.
\begin{enumerate}
\item Une famille $(\vec{u}, \vec{v}, \vec{w})$ de $\mathbb{R}^3$ est dite \cle{libre} lorsque l'égalité $\alpha \vec{u} + \beta \vec{v} + \gamma \vec{w} = \vec{0}$ entraîne \ldots
\item Une \cle{base} de $\mathbb{R}^3$ est une famille à la fois \ldots{} et \ldots{} de $\mathbb{R}^3$.
\item Soit $f : E \to F$ une application linéaire. Le \cle{noyau} de $f$ est l'ensemble $\ker f = \{\ldots\}$ ; c'est un sous-espace vectoriel de \ldots
\item Le théorème du rang affirme que pour tout endomorphisme $f$ de $\mathbb{R}^3$ : $\dim \mathbb{R}^3 = \ldots + \ldots$
\end{enumerate}
\end{exercice}

\begin{exercice}[Calculs directs]
\begin{enumerate}
\item Les vecteurs $\vec{u} = (1, 0, 2)$ et $\vec{v} = (2, 0, 4)$ sont-ils colinéaires ? La famille $(\vec{u}, \vec{v})$ est-elle libre ?
\item Soit $f(x, y, z) = (x + z, \, y - x, \, 2z)$. Calculer $f(1, 0, 1)$, $f(0, 1, 0)$ et $f(2, -1, 3)$.
\item Déterminer les coordonnées du vecteur $\vec{a} = (5, 3, 0)$ dans la base $(\vec{i}, \vec{j}, \vec{k})$ de $\mathbb{R}^3$ avec $\vec{i} = (1, 1, 0)$, $\vec{j} = (1, -1, 0)$ et $\vec{k} = (0, 0, 1)$.
\item Soit $f$ l'endomorphisme de $\mathbb{R}^3$ de matrice $A = \begin{pmatrix} 1 & 2 & 0 \\ 0 & 1 & 1 \\ 0 & 0 & 1 \end{pmatrix}$ dans la base canonique. Calculer l'image du vecteur $(1, 1, 1)$ par $f$.
\end{enumerate}
\end{exercice}

\courssub{Je m'entraîne}

\begin{methode}[Montrer qu'un ensemble est un sous-espace vectoriel]
Pour montrer qu'une partie $F$ d'un espace vectoriel $E$ est un sous-espace vectoriel, on vérifie trois conditions :
\begin{enumerate}
\item $\vec{0}_E \in F$ (non-vacuité) ;
\item $\forall \vec{u}, \vec{v} \in F,\; \vec{u} + \vec{v} \in F$ (stabilité par l'addition) ;
\item $\forall \vec{u} \in F,\; \forall \lambda \in \mathbb{R},\; \lambda \vec{u} \in F$ (stabilité par la multiplication par un scalaire).
\end{enumerate}
\emph{Astuce :} Il suffit souvent de vérifier que $F$ est l'ensemble des solutions d'un système homogène d'équations linéaires (noyau d'une application linéaire).
\end{methode}

\begin{exercice}[Un sous-espace vectoriel de $\mathbb{R}^3$]
On considère l'ensemble $F = \{(x, y, z) \in \mathbb{R}^3 \mid 2x - y + z = 0\}$.
\begin{enumerate}
\item Montrer que $F$ est un sous-espace vectoriel de $\mathbb{R}^3$.
\item Montrer que tout vecteur de $F$ s'écrit $(x, y, z) = x(1, 2, 0) + z(0, 1, 1)$.
\item En déduire une base de $F$ et préciser $\dim F$.
\item Le vecteur $(1, 3, 1)$ appartient-il à $F$ ?
\end{enumerate}
\end{exercice}

\begin{exercice}[Une base de $\mathbb{R}^3$ et des coordonnées]
On considère les vecteurs $\vec{u} = (1, 2, 1)$, $\vec{v} = (2, 1, 0)$ et $\vec{w} = (1, -1, 1)$.
\begin{enumerate}
\item Montrer que la famille $(\vec{u}, \vec{v}, \vec{w})$ est libre.
\item Justifier que $(\vec{u}, \vec{v}, \vec{w})$ est une base de $\mathbb{R}^3$.
\item Déterminer les coordonnées du vecteur $\vec{a} = (4, 5, 3)$ dans la base $(\vec{u}, \vec{v}, \vec{w})$.
\item Déterminer les coordonnées du vecteur $\vec{b} = (0, 0, 1)$ dans cette même base.
\end{enumerate}
\end{exercice}

\begin{methode}[Déterminer le noyau et l'image d'un endomorphisme]
Soit $f : \mathbb{R}^3 \to \mathbb{R}^3$ un endomorphisme de matrice $A$ dans la base canonique.
\begin{enumerate}
\item \textbf{Noyau :} Résoudre le système linéaire $A X = 0$ (ou $f(x,y,z) = (0,0,0)$). Les solutions forment $\ker f$. Une base de l'ensemble des solutions est une base de $\ker f$.
\item \textbf{Image :} $\operatorname{Im} f$ est engendré par les colonnes de $A$. En extraire une base (colonnes pivots après réduction de Gauss, ou colonnes correspondant aux variables principales).
\item \textbf{Théorème du rang :} Vérifier que $\dim \ker f + \dim \operatorname{Im} f = 3$.
\end{enumerate}
\end{methode}

\begin{exercice}[Étude complète d'un endomorphisme]
Soit $f$ l'application de $\mathbb{R}^3$ dans $\mathbb{R}^3$ définie par :
\[ f(x, y, z) = (x - y + z, \; 2x + y, \; x + 2y - z). \]
\begin{enumerate}
\item Montrer que $f$ est une application linéaire.
\item Écrire la matrice $A$ de $f$ dans la base canonique de $\mathbb{R}^3$.
\item Déterminer le noyau $\ker f$ : résoudre le système $f(x, y, z) = (0, 0, 0)$ et en donner une base.
\item Déterminer une base de $\operatorname{Im} f$.
\item Vérifier que le théorème du rang est satisfait.
\item L'application $f$ est-elle bijective ? Justifier.
\end{enumerate}
\end{exercice}

\begin{exercice}[Application définie par les images des vecteurs de base]
Soit $(\vec{e}_1, \vec{e}_2, \vec{e}_3)$ la base canonique de $\mathbb{R}^3$ et $f$ l'unique application linéaire de $\mathbb{R}^3$ dans $\mathbb{R}^3$ telle que :
\[ f(\vec{e}_1) = (1, 0, 1), \qquad f(\vec{e}_2) = (0, 1, 1), \qquad f(\vec{e}_3) = (1, 1, 2). \]
\begin{enumerate}
\item Justifier l'existence et l'unicité d'une telle application linéaire.
\item Écrire la matrice de $f$ dans la base canonique.
\item Donner l'expression de $f(x, y, z)$ pour tout $(x, y, z) \in \mathbb{R}^3$.
\item Montrer que $f(\vec{e}_3) = f(\vec{e}_1) + f(\vec{e}_2)$. En déduire un vecteur non nul de $\ker f$, puis que $f$ n'est pas injective.
\item Déterminer $\ker f$ et retrouver $\dim \operatorname{Im} f$ à l'aide du théorème du rang.
\end{enumerate}
\end{exercice}

\courssub{Je me prépare au Bac}

\begin{savaistu}[Le théorème du rang en télécommunications]
Dans les réseaux de téléphonie mobile (MTN, Orange Cameroun), le trafic entre antennes relais est modélisé par des matrices. Le \cle{rang} de la matrice de transfert indique le nombre de canaux indépendants. Si le rang est maximal (3 pour 3 antennes), toute configuration de trafic de sortie est atteignable (application bijective). Si le rang est inférieur, certaines pannes ou saturations rendent des configurations inaccessibles : le noyau représente les configurations d'entrée "invisibles" (perdues ou annulées).
\end{savaistu}

\begin{exercice}[Type Bac — Réseau de téléphonie et endomorphisme]
Une entreprise de télécommunications modélise la redistribution du trafic entre trois antennes relais de Yaoundé par un endomorphisme $f$ de $\mathbb{R}^3$, dont la matrice dans la base canonique est :
\[ A = \begin{pmatrix} 1 & 1 & 2 \\ 0 & 1 & 1 \\ 1 & 2 & 3 \end{pmatrix}. \]
\begin{enumerate}
\item Donner l'expression explicite de $f(x, y, z)$ pour tout $(x, y, z) \in \mathbb{R}^3$.
\item Montrer que les vecteurs $f(\vec{e}_1)$, $f(\vec{e}_2)$ et $f(\vec{e}_3)$ sont liés (on pourra chercher une relation entre les colonnes de $A$).
\item Déterminer le rang de $f$, c'est-à-dire $\dim \operatorname{Im} f$.
\item L'application $f$ est-elle bijective ? Justifier. Que peut-on en conclure sur la redistribution du trafic (toutes les configurations de sortie sont-elles atteignables) ?
\item Déterminer une base de $\ker f$ et interpréter : quelles configurations de trafic en entrée produisent un trafic nul en sortie ?
\item Déterminer une base de $\operatorname{Im} f$ et vérifier le théorème du rang.
\end{enumerate}
\end{exercice}

\begin{popfigure}
\begin{tikzpicture}[scale=1.2, >=Stealth]
% Plan P: x+y+z=0
\fill[popblue!10, opacity=0.5] (2,0,-2) -- (0,2,-2) -- (-2,0,2) -- (0,-2,2) -- cycle;
\draw[popblue, thick] (2,0,-2) -- (0,2,-2) -- (-2,0,2) -- (0,-2,2) -- cycle;
% Droite D: (1,1,1)
\draw[poporange, thick, ->] (0,0,0) -- (1.5,1.5,1.5) node[above right] {$\vec{d}=(1,1,1)$};
% Vecteur x
\draw[popdark, thick, ->] (0,0,0) -- (1,0.5,-0.2) node[midway, above] {$\vec{x}$};
% Projection p
\draw[popgreen, thick, ->] (0,0,0) -- (0.5,0.83,-1.33) node[midway, left] {$\vec{p}=f(\vec{x})$};
% q = t*d
\draw[poppurple, thick, ->] (0.5,0.83,-1.33) -- (1,0.5,-0.2) node[midway, right] {$\vec{q} \in D$};
% Axes
\draw[->] (-2,0,0) -- (2,0,0) node[right] {$x$};
\draw[->] (0,-2,0) -- (0,2,0) node[above] {$y$};
\draw[->] (0,0,-2) -- (0,0,2) node[above] {$z$};
\node[popblue] at (1.5,0,-1.5) {$P: x+y+z=0$};
\end{tikzpicture}
\legende{Projection sur le plan $P$ parallèlement à la droite $D$.}
\end{popfigure}

\begin{exercice}[Type Bac — Projection sur un plan]
Dans l'espace vectoriel $\mathbb{R}^3$, on considère le plan vectoriel $P$ d'équation $x + y + z = 0$ et la droite vectorielle $D$ dirigée par le vecteur $\vec{d} = (1, 1, 1)$. On admet que $\mathbb{R}^3 = P \oplus D$ : tout vecteur $\vec{x}$ de $\mathbb{R}^3$ s'écrit de manière unique $\vec{x} = \vec{p} + \vec{q}$ avec $\vec{p} \in P$ et $\vec{q} \in D$. On définit alors la projection $f$ sur $P$ parallèlement à $D$ par $f(\vec{x}) = \vec{p}$.
\begin{enumerate}
\item Sans aucun calcul, en utilisant l'interprétation géométrique de $f$, déterminer $\ker f$ et $\operatorname{Im} f$, puis leurs dimensions.
\item Vérifier que le théorème du rang est bien satisfait.
\item Montrer que $f$ est une application linéaire.
\item En écrivant $(x, y, z) = \vec{p} + t\vec{d}$ avec $\vec{p} \in P$ et $t \in \mathbb{R}$, montrer que $t = \dfrac{x + y + z}{3}$.
\item En déduire l'expression explicite de $f(x, y, z)$.
\item Écrire la matrice de $f$ dans la base canonique de $\mathbb{R}^3$.
\item Calculer $f \circ f$. Que constate-t-on ? Justifier géométriquement ce résultat.
\end{enumerate}
\end{exercice}

\begin{attention}[Pièges fréquents sur le changement de base]
\begin{itemize}
\item Ne pas confondre la matrice de passage $P_{\mathcal{B} \to \mathcal{C}}$ (colonnes = coordonnées des vecteurs de $\mathcal{B}$ dans $\mathcal{C}$) et la matrice de l'application identité.
\item Pour passer de la matrice dans $\mathcal{B}$ à la matrice dans $\mathcal{C}$ : $M_{\mathcal{C}} = P_{\mathcal{B} \to \mathcal{C}} \; M_{\mathcal{B}} \; P_{\mathcal{C} \to \mathcal{B}}$.
\item Vérifier toujours que la famille proposée est bien une base (libre + génératrice, ou libre de cardinal = dimension).
\end{itemize}
\end{attention}

\begin{exercice}[Type Bac — Changement de base et dimension]
Une coopérative agricole de l'Ouest conditionne ses produits (maïs, haricots, manioc) selon trois recettes de mélanges, modélisées par les vecteurs de $\mathbb{R}^3$ :
\[ \vec{u} = (1, 1, 0), \qquad \vec{v} = (1, 0, 1), \qquad \vec{w} = (0, 1, 1). \]
\textbf{Partie A : une nouvelle base de $\mathbb{R}^3$}
\begin{enumerate}
\item Montrer que la famille $\mathcal{B} = (\vec{u}, \vec{v}, \vec{w})$ est libre, puis que c'est une base de $\mathbb{R}^3$.
\item Déterminer les coordonnées du vecteur $\vec{a} = (3, 2, 5)$ dans la base $\mathcal{B}$.
\item Soit $(\vec{e}_1, \vec{e}_2, \vec{e}_3)$ la base canonique. Exprimer $\vec{e}_1$ dans la base $\mathcal{B}$.
\end{enumerate}
\textbf{Partie B : un endomorphisme de production}
On considère l'endomorphisme $g$ de $\mathbb{R}^3$ défini par :
\[ g(\vec{u}) = \vec{u} + \vec{v}, \qquad g(\vec{v}) = \vec{v} + \vec{w}, \qquad g(\vec{w}) = \vec{w} + \vec{u}. \]
\begin{enumerate}
\item Écrire la matrice $M$ de $g$ dans la base $\mathcal{B}$.
\item Déterminer $\ker g$ (on pourra travailler avec les coordonnées dans $\mathcal{B}$). L'application $g$ est-elle injective ?
\item À l'aide du théorème du rang, donner $\dim \operatorname{Im} g$ et conclure que $g$ est bijective.
\item Déterminer la matrice de $g$ dans la base canonique (on pourra calculer $g(\vec{e}_1)$, $g(\vec{e}_2)$, $g(\vec{e}_3)$ à l'aide de la question A.3).
\end{enumerate}
\end{exercice}

\newpage

\coursec{Corrigés des exercices}

\begin{corrige}[Exercice 1 — QCM : une seule bonne réponse]
\begin{enumerate}
\item \textbf{Réponse b.} $F$ est l'ensemble des solutions de l'équation linéaire homogène $x + 2y - z = 0$ : c'est le noyau de la forme linéaire $(x,y,z) \mapsto x + 2y - z$, donc un sous-espace vectoriel de $\mathbb{R}^3$. On a $z = x + 2y$, donc tout vecteur de $F$ s'écrit $(x, y, x+2y) = x(1,0,1) + y(0,1,2)$. La famille $((1,0,1), (0,1,2))$ est libre (vecteurs non colinéaires) et engendre $F$ : c'est une base de $F$, donc $\dim F = 2$.
\item \textbf{Réponse b.} Dans un espace vectoriel de dimension $3$, toute famille de plus de $3$ vecteurs est liée. Une famille de $4$ vecteurs de $\mathbb{R}^3$ est donc \textbf{toujours liée}.
\item \textbf{Réponse b.} Par le théorème du rang : $\dim \ker f + \dim \operatorname{Im} f = \dim \mathbb{R}^3 = 3$, donc $\dim \operatorname{Im} f = 3 - 1 = 2$.
\item \textbf{Réponse a.} Les colonnes de la matrice sont les images des vecteurs de la base canonique :
\[ f(\vec{e}_1) = f(1,0,0) = (1, 0, 1), \quad f(\vec{e}_2) = (1, 0, 0), \quad f(\vec{e}_3) = (0, 1, -1). \]
D'où $A = \begin{pmatrix} 1 & 1 & 0 \\ 0 & 0 & 1 \\ 1 & 0 & -1 \end{pmatrix}$.
\end{enumerate}
\end{corrige}

\begin{corrige}[Exercice 2 — Vrai ou faux, en justifiant]
\begin{enumerate}
\item \textbf{Faux.} Contre-exemple : les trois vecteurs $\vec{u} = (1,0,0)$, $\vec{v} = (0,1,0)$, $\vec{w} = (1,1,0)$ sont non nuls, mais $\vec{w} = \vec{u} + \vec{v}$, donc la famille est liée : ce n'est pas une base. (Être non nuls ne suffit pas ; il faut que la famille soit libre.)
\item \textbf{Vrai.} $\mathbb{R}^3$ est de dimension $3$. Une famille libre de $3$ vecteurs dans un espace de dimension $3$ est automatiquement une base (une famille libre de cardinal égal à la dimension est aussi génératrice).
\item \textbf{Faux.} Le vecteur nul $(0,0,0)$ ne vérifie pas $x + y + z = 1$ (car $0 + 0 + 0 = 0 \neq 1$), donc $\vec{0} \notin G$. Or tout sous-espace vectoriel contient le vecteur nul : $G$ n'est pas un sous-espace vectoriel de $\mathbb{R}^3$. (L'équation n'est pas \emph{homogène}.)
\item \textbf{Vrai.} $f$ injective signifie $\ker f = \{\vec{0}\}$, donc $\dim \ker f = 0$. Par le théorème du rang : $\dim \operatorname{Im} f = 3 - 0 = 3$. Comme $\operatorname{Im} f$ est un sous-espace de $\mathbb{R}^3$ de dimension $3$, on a $\operatorname{Im} f = \mathbb{R}^3$ : $f$ est surjective. (En dimension finie, pour un endomorphisme : injectif $\iff$ surjectif $\iff$ bijectif.)
\end{enumerate}
\end{corrige}

\begin{corrige}[Exercice 3 — Définitions à compléter]
\begin{enumerate}
\item Une famille $(\vec{u}, \vec{v}, \vec{w})$ de $\mathbb{R}^3$ est dite \cle{libre} lorsque l'égalité $\alpha \vec{u} + \beta \vec{v} + \gamma \vec{w} = \vec{0}$ entraîne \textbf{$\alpha = \beta = \gamma = 0$}.
\item Une \cle{base} de $\mathbb{R}^3$ est une famille à la fois \textbf{libre} et \textbf{génératrice} de $\mathbb{R}^3$.
\item Le \cle{noyau} de $f$ est l'ensemble $\ker f = \{\vec{x} \in E \mid f(\vec{x}) = \vec{0}_F\}$ ; c'est un sous-espace vectoriel de \textbf{$E$ (l'espace de départ)}.
\item Le théorème du rang affirme que pour tout endomorphisme $f$ de $\mathbb{R}^3$ : $\dim \mathbb{R}^3 = \dim \ker f + \dim \operatorname{Im} f$, c'est-à-dire $3 = \dim \ker f + \operatorname{rg} f$.
\end{enumerate}
\end{corrige}

\begin{corrige}[Exercice 4 — Calculs directs]
\begin{enumerate}
\item On remarque que $\vec{v} = (2, 0, 4) = 2(1, 0, 2) = 2\vec{u}$. Les vecteurs $\vec{u}$ et $\vec{v}$ sont donc \textbf{colinéaires}. La famille $(\vec{u}, \vec{v})$ est \textbf{liée} (donc pas libre), puisque $2\vec{u} - \vec{v} = \vec{0}$ est une combinaison linéaire nulle à coefficients non tous nuls.
\item $f(x, y, z) = (x + z, \, y - x, \, 2z)$ :
\begin{align*}
f(1, 0, 1) &= (1 + 1, \, 0 - 1, \, 2 \times 1) = (2, -1, 2) ; \\
f(0, 1, 0) &= (0 + 0, \, 1 - 0, \, 0) = (0, 1, 0) ; \\
f(2, -1, 3) &= (2 + 3, \, -1 - 2, \, 6) = (5, -3, 6).
\end{align*}
\item On cherche $\alpha, \beta, \gamma$ tels que $\vec{a} = \alpha \vec{i} + \beta \vec{j} + \gamma \vec{k}$ :
\[ (5, 3, 0) = \alpha(1,1,0) + \beta(1,-1,0) + \gamma(0,0,1) = (\alpha + \beta, \, \alpha - \beta, \, \gamma). \]
D'où le système : $\alpha + \beta = 5$, $\alpha - \beta = 3$, $\gamma = 0$. En additionnant : $2\alpha = 8$ donc $\alpha = 4$, puis $\beta = 1$, et $\gamma = 0$. Les coordonnées de $\vec{a}$ dans $(\vec{i}, \vec{j}, \vec{k})$ sont $\boxed{(4, 1, 0)}$. Vérification : $4(1,1,0) + 1(1,-1,0) = (5, 3, 0)$. \checkmark
\item L'image de $(1,1,1)$ s'obtient par le produit matriciel :
\[ A \begin{pmatrix} 1 \\ 1 \\ 1 \end{pmatrix} = \begin{pmatrix} 1 & 2 & 0 \\ 0 & 1 & 1 \\ 0 & 0 & 1 \end{pmatrix} \begin{pmatrix} 1 \\ 1 \\ 1 \end{pmatrix} = \begin{pmatrix} 1 + 2 + 0 \\ 0 + 1 + 1 \\ 0 + 0 + 1 \end{pmatrix} = \begin{pmatrix} 3 \\ 2 \\ 1 \end{pmatrix}. \]
Donc $f(1, 1, 1) = (3, 2, 1)$.
\end{enumerate}
\end{corrige}

\begin{corrige}[Exercice 5 — Un sous-espace vectoriel de $\mathbb{R}^3$]
\begin{enumerate}
\item On vérifie les trois conditions :
\begin{itemize}
\item $\vec{0} = (0,0,0)$ vérifie $2 \times 0 - 0 + 0 = 0$, donc $\vec{0} \in F$ ;
\item Soient $\vec{u} = (x, y, z) \in F$ et $\vec{u}' = (x', y', z') \in F$. Alors $\vec{u} + \vec{u}' = (x+x', y+y', z+z')$ et
\[ 2(x+x') - (y+y') + (z+z') = \underbrace{(2x - y + z)}_{=0} + \underbrace{(2x' - y' + z')}_{=0} = 0, \]
donc $\vec{u} + \vec{u}' \in F$ ;
\item Soient $\vec{u} = (x,y,z) \in F$ et $\lambda \in \mathbb{R}$. Alors $\lambda \vec{u} = (\lambda x, \lambda y, \lambda z)$ et $2\lambda x - \lambda y + \lambda z = \lambda(2x - y + z) = 0$, donc $\lambda \vec{u} \in F$.
\end{itemize}
$F$ est donc un sous-espace vectoriel de $\mathbb{R}^3$. \emph{(Autre rédaction : $F$ est le noyau de la forme linéaire $(x,y,z) \mapsto 2x - y + z$.)}
\item Si $(x, y, z) \in F$, alors $y = 2x + z$. Donc :
\[ (x, y, z) = (x, \, 2x + z, \, z) = x(1, 2, 0) + z(0, 1, 1). \]
\item Posons $\vec{a} = (1, 2, 0)$ et $\vec{b} = (0, 1, 1)$. D'après la question précédente, $(\vec{a}, \vec{b})$ engendre $F$. De plus, $\vec{a}$ et $\vec{b}$ ne sont pas colinéaires (par exemple $\vec{a}$ a une troisième coordonnée nulle et pas $\vec{b}$, et $\vec{b} \neq \vec{0}$), donc la famille $(\vec{a}, \vec{b})$ est libre. C'est une \textbf{base de $F$}, et $\boxed{\dim F = 2}$ ($F$ est un plan vectoriel).
\item Le vecteur $(1, 3, 1)$ vérifie-t-il $2x - y + z = 0$ ? On calcule : $2 \times 1 - 3 + 1 = 0$. \textbf{Oui}, $(1, 3, 1) \in F$. (D'ailleurs $(1,3,1) = \vec{a} + \vec{b}$.)
\end{enumerate}
\end{corrige}

\begin{corrige}[Exercice 6 — Une base de $\mathbb{R}^3$ et des coordonnées]
\begin{enumerate}
\item Soient $\alpha, \beta, \gamma \in \mathbb{R}$ tels que $\alpha \vec{u} + \beta \vec{v} + \gamma \vec{w} = \vec{0}$ :
\[ \alpha(1,2,1) + \beta(2,1,0) + \gamma(1,-1,1) = (0,0,0) \iff \begin{cases} \alpha + 2\beta + \gamma = 0 & (L_1) \\ 2\alpha + \beta - \gamma = 0 & (L_2) \\ \alpha + \gamma = 0 & (L_3) \end{cases} \]
De $(L_3)$ : $\gamma = -\alpha$. En reportant dans $(L_1)$ : $2\beta = 0$, donc $\beta = 0$. En reportant dans $(L_2)$ : $2\alpha + \alpha = 3\alpha = 0$, donc $\alpha = 0$, puis $\gamma = 0$.
Ainsi $\alpha = \beta = \gamma = 0$ : la famille $(\vec{u}, \vec{v}, \vec{w})$ est \textbf{libre}.
\item $(\vec{u}, \vec{v}, \vec{w})$ est une famille libre de $3$ vecteurs dans $\mathbb{R}^3$ qui est de dimension $3$ : c'est donc une \textbf{base de $\mathbb{R}^3$}.
\item On cherche $\alpha, \beta, \gamma$ tels que $\vec{a} = \alpha \vec{u} + \beta \vec{v} + \gamma \vec{w}$ :
\[ \begin{cases} \alpha + 2\beta + \gamma = 4 & (L_1) \\ 2\alpha + \beta - \gamma = 5 & (L_2) \\ \alpha + \gamma = 3 & (L_3) \end{cases} \]
$(L_1) + (L_2)$ : $3\alpha + 3\beta = 9$, soit $\alpha + \beta = 3$. Avec $(L_3)$ : $\gamma = 3 - \alpha$. Reportons dans $(L_2)$ : $2\alpha + \beta - (3 - \alpha) = 5$, soit $3\alpha + \beta = 8$. En soustrayant $\alpha + \beta = 3$ : $2\alpha = 5$, $\alpha = \dfrac{5}{2}$, puis $\beta = \dfrac{1}{2}$ et $\gamma = \dfrac{1}{2}$.
Vérification : $\dfrac{5}{2}(1,2,1) + \dfrac{1}{2}(2,1,0) + \dfrac{1}{2}(1,-1,1) = \left(\dfrac{5+2+1}{2}, \dfrac{10+1-1}{2}, \dfrac{5+0+1}{2}\right) = (4, 5, 3)$. \checkmark
Les coordonnées de $\vec{a}$ dans $(\vec{u}, \vec{v}, \vec{w})$ sont $\left(\dfrac{5}{2}, \dfrac{1}{2}, \dfrac{1}{2}\right)$.
\item Pour $\vec{b} = (0, 0, 1)$ :
\[ \begin{cases} \alpha + 2\beta + \gamma = 0 \\ 2\alpha + \beta - \gamma = 0 \\ \alpha + \gamma = 1 \end{cases} \]
De la troisième ligne : $\gamma = 1 - \alpha$. En reportant dans la première : $\alpha + 2\beta + 1 - \alpha = 0$, soit $\beta = -\dfrac{1}{2}$. Dans la deuxième : $2\alpha - \dfrac{1}{2} - 1 + \alpha = 0$, soit $3\alpha = \dfrac{3}{2}$, $\alpha = \dfrac{1}{2}$, puis $\gamma = \dfrac{1}{2}$.
Vérification : $\dfrac{1}{2}(1,2,1) - \dfrac{1}{2}(2,1,0) + \dfrac{1}{2}(1,-1,1) = \left(0, 0, 1\right)$. \checkmark
Les coordonnées de $\vec{b}$ sont $\left(\dfrac{1}{2}, -\dfrac{1}{2}, \dfrac{1}{2}\right)$.
\end{enumerate}
\end{corrige}

\begin{corrige}[Exercice 7 — Étude complète d'un endomorphisme]
\begin{enumerate}
\item Soient $\vec{x} = (x, y, z)$, $\vec{x}' = (x', y', z')$ dans $\mathbb{R}^3$ et $\lambda \in \mathbb{R}$.
\begin{align*}
f(\vec{x} + \vec{x}') &= f(x + x', \, y + y', \, z + z') \\
&= \big((x+x') - (y+y') + (z+z'), \; 2(x+x') + (y+y'), \; (x+x') + 2(y+y') - (z+z')\big) \\
&= (x - y + z, \, 2x + y, \, x + 2y - z) + (x' - y' + z', \, 2x' + y', \, x' + 2y' - z') \\
&= f(\vec{x}) + f(\vec{x}').
\end{align*}
\begin{align*}
f(\lambda \vec{x}) &= (\lambda x - \lambda y + \lambda z, \; 2\lambda x + \lambda y, \; \lambda x + 2\lambda y - \lambda z) = \lambda f(\vec{x}).
\end{align*}
Donc $f$ est une \textbf{application linéaire} (endomorphisme de $\mathbb{R}^3$).
\item Les colonnes de $A$ sont les images des vecteurs de la base canonique :
\[ f(\vec{e}_1) = (1, 2, 1), \quad f(\vec{e}_2) = (-1, 1, 2), \quad f(\vec{e}_3) = (1, 0, -1). \]
\[ A = \begin{pmatrix} 1 & -1 & 1 \\ 2 & 1 & 0 \\ 1 & 2 & -1 \end{pmatrix}. \]
\item $f(x, y, z) = (0, 0, 0) \iff \begin{cases} x - y + z = 0 & (L_1) \\ 2x + y = 0 & (L_2) \\ x + 2y - z = 0 & (L_3) \end{cases}$

De $(L_2)$ : $y = -2x$. De $(L_1)$ : $z = y - x = -3x$. Vérifions $(L_3)$ : $x + 2(-2x) - (-3x) = x - 4x + 3x = 0$. \checkmark{} Le système est compatible pour tout $x$.
Donc $(x, y, z) = (x, -2x, -3x) = x(1, -2, -3)$, et
\[ \ker f = \operatorname{Vect}\big((1, -2, -3)\big). \]
Le vecteur $(1, -2, -3)$ étant non nul, il forme une \textbf{base de $\ker f$}, et $\dim \ker f = 1$.
\item $\operatorname{Im} f$ est engendré par les colonnes de $A$ : $\vec{c}_1 = (1, 2, 1)$, $\vec{c}_2 = (-1, 1, 2)$, $\vec{c}_3 = (1, 0, -1)$. On remarque que $\vec{c}_3 = -\vec{c}_1 - \vec{c}_2$ (vérification : $-(1,2,1) - (-1,1,2) = (0, -3, -3)$... reprenons proprement : cherchons la relation). En fait, d'après le noyau : $f(\vec{e}_1) - 2 f(\vec{e}_2) - 3 f(\vec{e}_3) = \vec{0}$, c'est-à-dire $\vec{c}_1 = 2\vec{c}_2 + 3\vec{c}_3$. Vérification : $2(-1,1,2) + 3(1,0,-1) = (-2+3, \, 2, \, 4-3) = (1, 2, 1) = \vec{c}_1$. \checkmark

Donc $\operatorname{Im} f = \operatorname{Vect}(\vec{c}_2, \vec{c}_3)$. Les vecteurs $\vec{c}_2 = (-1, 1, 2)$ et $\vec{c}_3 = (1, 0, -1)$ ne sont pas colinéaires, donc forment une famille libre : c'est une \textbf{base de $\operatorname{Im} f$}, et $\dim \operatorname{Im} f = 2$.
\item Théorème du rang : $\dim \ker f + \dim \operatorname{Im} f = 1 + 2 = 3 = \dim \mathbb{R}^3$. \checkmark
\item $f$ n'est \textbf{pas bijective}, car $\ker f \neq \{\vec{0}\}$ ($f$ n'est pas injective ; par exemple $f(1, -2, -3) = \vec{0}$ avec $(1,-2,-3) \neq \vec{0}$). De même $\operatorname{Im} f \neq \mathbb{R}^3$, donc $f$ n'est pas surjective.
\end{enumerate}
\end{corrige}

\begin{corrige}[Exercice 8 — Application définie par les images des vecteurs de base]
\begin{enumerate}
\item $(\vec{e}_1, \vec{e}_2, \vec{e}_3)$ est une base de $\mathbb{R}^3$. Une application linéaire est entièrement déterminée par les images des vecteurs d'une base : tout $\vec{x} = x\vec{e}_1 + y\vec{e}_2 + z\vec{e}_3$ a pour image $f(\vec{x}) = x f(\vec{e}_1) + y f(\vec{e}_2) + z f(\vec{e}_3)$. D'où l'\textbf{existence et l'unicité} de $f$.
\item Les colonnes de la matrice sont $f(\vec{e}_1)$, $f(\vec{e}_2)$, $f(\vec{e}_3)$ :
\[ A = \begin{pmatrix} 1 & 0 & 1 \\ 0 & 1 & 1 \\ 1 & 1 & 2 \end{pmatrix}. \]
\item Pour tout $(x, y, z) \in \mathbb{R}^3$ :
\[ f(x, y, z) = x(1,0,1) + y(0,1,1) + z(1,1,2) = (x + z, \; y + z, \; x + y + 2z). \]
\item $f(\vec{e}_1) + f(\vec{e}_2) = (1, 0, 1) + (0, 1, 1) = (1, 1, 2) = f(\vec{e}_3)$. \checkmark

Par linéarité : $f(\vec{e}_1 + \vec{e}_2 - \vec{e}_3) = f(\vec{e}_1) + f(\vec{e}_2) - f(\vec{e}_3) = \vec{0}$. Le vecteur $\vec{e}_1 + \vec{e}_2 - \vec{e}_3 = (1, 1, -1)$ est non nul et appartient à $\ker f$. Donc $\ker f \neq \{\vec{0}\}$ : $f$ \textbf{n'est pas injective}.
\item Résolvons $f(x, y, z) = (0,0,0)$ :
\[ \begin{cases} x + z = 0 \\ y + z = 0 \\ x + y + 2z = 0 \end{cases} \iff \begin{cases} x = -z \\ y = -z \\ -z - z + 2z = 0 \end{cases} \]
La troisième équation est automatiquement vérifiée. Donc $(x, y, z) = (-z, -z, z) = z(-1, -1, 1)$, et
\[ \ker f = \operatorname{Vect}\big((-1, -1, 1)\big) = \operatorname{Vect}\big((1, 1, -1)\big), \qquad \dim \ker f = 1. \]
Par le théorème du rang : $\dim \operatorname{Im} f = 3 - 1 = 2$. (On vérifie : $f(\vec{e}_1)$ et $f(\vec{e}_2)$ ne sont pas colinéaires et engendrent $\operatorname{Im} f$ puisque $f(\vec{e}_3) = f(\vec{e}_1) + f(\vec{e}_2)$.)
\end{enumerate}
\end{corrige}

\begin{corrige}[Exercice 9 — Type Bac — Réseau de téléphonie et endomorphisme]
\emph{Barème indicatif : 1. : 1 pt ; 2. : 1,5 pt ; 3. : 1 pt ; 4. : 1,5 pt ; 5. : 2 pts ; 6. : 2 pts.}

\begin{enumerate}
\item La première colonne de $A$ donne $f(\vec{e}_1) = (1, 0, 1)$, la deuxième $f(\vec{e}_2) = (1, 1, 2)$, la troisième $f(\vec{e}_3) = (2, 1, 3)$. D'où :
\[ f(x, y, z) = x f(\vec{e}_1) + y f(\vec{e}_2) + z f(\vec{e}_3) = (x + y + 2z, \; y + z, \; x + 2y + 3z). \]
\item Notons $\vec{c}_1 = (1, 0, 1)$, $\vec{c}_2 = (1, 1, 2)$, $\vec{c}_3 = (2, 1, 3)$ les colonnes de $A$. On observe que :
\[ \vec{c}_1 + \vec{c}_2 = (1+1, \, 0+1, \, 1+2) = (2, 1, 3) = \vec{c}_3. \]
La relation $\vec{c}_1 + \vec{c}_2 - \vec{c}_3 = \vec{0}$ est une combinaison linéaire nulle à coefficients non tous nuls : la famille $(f(\vec{e}_1), f(\vec{e}_2), f(\vec{e}_3))$ est \textbf{liée}.
\item $\operatorname{Im} f = \operatorname{Vect}(\vec{c}_1, \vec{c}_2, \vec{c}_3) = \operatorname{Vect}(\vec{c}_1, \vec{c}_2)$ puisque $\vec{c}_3 = \vec{c}_1 + \vec{c}_2$. Les vecteurs $\vec{c}_1 = (1, 0, 1)$ et $\vec{c}_2 = (1, 1, 2)$ ne sont pas colinéaires (leurs coordonnées ne sont pas proportionnelles), donc forment une famille libre. Ainsi :
\[ \operatorname{rg} f = \dim \operatorname{Im} f = 2. \]
\item $\operatorname{Im} f \neq \mathbb{R}^3$ (dimension $2 < 3$), donc $f$ n'est pas surjective, et par suite \textbf{pas bijective}. \emph{Interprétation :} toutes les configurations de trafic en sortie ne sont \textbf{pas atteignables} ; seules celles appartenant au plan $\operatorname{Im} f$ le sont. Le réseau présente une redondance (la troisième colonne est combinaison des deux premières) qui limite les sorties possibles.
\item $f(x, y, z) = (0,0,0) \iff \begin{cases} x + y + 2z = 0 & (L_1) \\ y + z = 0 & (L_2) \\ x + 2y + 3z = 0 & (L_3) \end{cases}$

De $(L_2)$ : $y = -z$. De $(L_1)$ : $x = -y - 2z = z - 2z = -z$. Vérification de $(L_3)$ : $-z + 2(-z) + 3z = 0$. \checkmark
Donc $(x, y, z) = (-z, -z, z) = z(-1, -1, 1)$ et
\[ \ker f = \operatorname{Vect}\big((-1, -1, 1)\big), \quad \text{base : } \big((-1, -1, 1)\big), \quad \dim \ker f = 1. \]
\emph{Interprétation :} les configurations d'entrée proportionnelles à $(-1, -1, 1)$ (par exemple retirer une unité de trafic des antennes 1 et 2 et en ajouter une à l'antenne 3) produisent un trafic nul en sortie : elles sont « invisibles » pour le réseau.
\item D'après la question 3, une base de $\operatorname{Im} f$ est $\big((1, 0, 1), \, (1, 1, 2)\big)$. Vérification du théorème du rang :
\[ \dim \ker f + \dim \operatorname{Im} f = 1 + 2 = 3 = \dim \mathbb{R}^3. \quad \checkmark \]
\end{enumerate}

\textbf{Points de vigilance :} citer explicitement le théorème du rang par son nom ; justifier la liberté de la famille retenue pour l'image (non-colinéarité) ; ne pas conclure à la bijectivité sans avoir comparé $\dim \operatorname{Im} f$ à $3$.
\end{corrige}

\begin{corrige}[Exercice 10 — Type Bac — Projection sur un plan]
\emph{Barème indicatif : 1. : 1,5 pt ; 2. : 0,5 pt ; 3. : 1,5 pt ; 4. : 1,5 pt ; 5. : 1 pt ; 6. : 1 pt ; 7. : 1,5 pt.}

\begin{enumerate}
\item Géométriquement, $f(\vec{x})$ est le projeté de $\vec{x}$ sur $P$ parallèlement à $D$.
\begin{itemize}
\item $f(\vec{x}) = \vec{0}$ signifie que $\vec{x}$ se projette en $\vec{0}$, c'est-à-dire que $\vec{x} \in D$. Donc $\ker f = D = \operatorname{Vect}(\vec{d})$ et $\dim \ker f = 1$.
\item L'image de $f$ est l'ensemble des projetés, c'est-à-dire le plan $P$ tout entier (tout point de $P$ est son propre projeté). Donc $\operatorname{Im} f = P$ et $\dim \operatorname{Im} f = 2$.
\end{itemize}
\item $\dim \ker f + \dim \operatorname{Im} f = 1 + 2 = 3 = \dim \mathbb{R}^3$ : le théorème du rang est satisfait. \checkmark
\item Soient $\vec{x}, \vec{x}' \in \mathbb{R}^3$ et $\lambda \in \mathbb{R}$. Écrivons les décompositions uniques : $\vec{x} = \vec{p} + \vec{q}$ et $\vec{x}' = \vec{p}' + \vec{q}'$ avec $\vec{p}, \vec{p}' \in P$ et $\vec{q}, \vec{q}' \in D$.
Alors $\vec{x} + \vec{x}' = (\vec{p} + \vec{p}') + (\vec{q} + \vec{q}')$. Comme $P$ et $D$ sont des sous-espaces vectoriels, $\vec{p} + \vec{p}' \in P$ et $\vec{q} + \vec{q}' \in D$. Par unicité de la décomposition :
\[ f(\vec{x} + \vec{x}') = \vec{p} + \vec{p}' = f(\vec{x}) + f(\vec{x}'). \]
De même, $\lambda \vec{x} = \lambda \vec{p} + \lambda \vec{q}$ avec $\lambda \vec{p} \in P$ et $\lambda \vec{q} \in D$, donc $f(\lambda \vec{x}) = \lambda \vec{p} = \lambda f(\vec{x})$. $f$ est bien \textbf{linéaire}.
\item Écrivons $(x, y, z) = \vec{p} + t\vec{d}$ avec $\vec{p} = (p_1, p_2, p_3) \in P$, c'est-à-dire $p_1 + p_2 + p_3 = 0$. Alors :
\[ (x, y, z) = (p_1 + t, \, p_2 + t, \, p_3 + t). \]
En additionnant les trois coordonnées :
\[ x + y + z = (p_1 + p_2 + p_3) + 3t = 0 + 3t = 3t, \quad \text{donc } t = \frac{x + y + z}{3}. \]
\item On en déduit $\vec{p} = (x, y, z) - t(1, 1, 1)$ :
\[ f(x, y, z) = \left( x - \frac{x+y+z}{3}, \; y - \frac{x+y+z}{3}, \; z - \frac{x+y+z}{3} \right) = \frac{1}{3}\left( 2x - y - z, \; -x + 2y - z, \; -x - y + 2z \right). \]
\item Les images des vecteurs de la base canonique :
\[ f(\vec{e}_1) = \frac{1}{3}(2, -1, -1), \quad f(\vec{e}_2) = \frac{1}{3}(-1, 2, -1), \quad f(\vec{e}_3) = \frac{1}{3}(-1, -1, 2). \]
\[ M = \frac{1}{3}\begin{pmatrix} 2 & -1 & -1 \\ -1 & 2 & -1 \\ -1 & -1 & 2 \end{pmatrix}. \]
\item Pour $\vec{x} = \vec{p} + \vec{q}$ avec $\vec{p} \in P$, $\vec{q} \in D$ : $f(\vec{x}) = \vec{p}$. Or $\vec{p} \in P$, donc sa décomposition est $\vec{p} = \vec{p} + \vec{0}$, et $f(f(\vec{x})) = f(\vec{p}) = \vec{p} = f(\vec{x})$. Ainsi :
\[ f \circ f = f. \]
On dit que $f$ est un \cle{projecteur} (application \emph{idempotente}). \emph{Justification géométrique :} une fois le vecteur projeté sur $P$, le résultat est déjà dans $P$ ; le projeter à nouveau ne change rien.
\end{enumerate}

\textbf{Points de vigilance :} pour la question 3, l'unicité de la décomposition $\vec{x} = \vec{p} + \vec{q}$ (justifiée par $\mathbb{R}^3 = P \oplus D$) doit être invoquée ; pour la question 7, on attend la propriété $f \circ f = f$ et son interprétation géométrique, pas seulement un calcul matriciel.
\end{corrige}

\begin{corrige}[Exercice 11 — Type Bac — Changement de base et dimension]
\emph{Barème indicatif : Partie A : 1. : 1,5 pt ; 2. : 1,5 pt ; 3. : 1 pt. Partie B : 1. : 1 pt ; 2. : 2 pts ; 3. : 1 pt ; 4. : 2 pts.}

\textbf{Partie A}
\begin{enumerate}
\item Soient $\alpha, \beta, \gamma \in \mathbb{R}$ tels que $\alpha \vec{u} + \beta \vec{v} + \gamma \vec{w} = \vec{0}$ :
\[ \alpha(1,1,0) + \beta(1,0,1) + \gamma(0,1,1) = (0,0,0) \iff \begin{cases} \alpha + \beta = 0 \\ \alpha + \gamma = 0 \\ \beta + \gamma = 0 \end{cases} \]
Des deux premières équations : $\beta = -\alpha$ et $\gamma = -\alpha$. La troisième donne $-2\alpha = 0$, donc $\alpha = 0$, puis $\beta = \gamma = 0$. La famille $\mathcal{B}$ est \textbf{libre}. Étant libre et de cardinal $3 = \dim \mathbb{R}^3$, c'est une \textbf{base de $\mathbb{R}^3$}.
\item On cherche $\alpha, \beta, \gamma$ tels que $\vec{a} = \alpha \vec{u} + \beta \vec{v} + \gamma \vec{w}$ :
\[ \begin{cases} \alpha + \beta = 3 & (L_1) \\ \alpha + \gamma = 2 & (L_2) \\ \beta + \gamma = 5 & (L_3) \end{cases} \]
$(L_1) + (L_2) - (L_3)$ : $2\alpha = 3 + 2 - 5 = 0$, donc $\alpha = 0$ ; puis $\beta = 3$ et $\gamma = 2$.
Vérification : $0 \cdot \vec{u} + 3(1,0,1) + 2(0,1,1) = (3, 2, 5)$. \checkmark
Les coordonnées de $\vec{a}$ dans $\mathcal{B}$ sont $(0, 3, 2)$.
\item On résout $\vec{e}_1 = (1, 0, 0) = \alpha \vec{u} + \beta \vec{v} + \gamma \vec{w}$ :
\[ \begin{cases} \alpha + \beta = 1 \\ \alpha + \gamma = 0 \\ \beta + \gamma = 0 \end{cases} \]
Des deux dernières : $\gamma = -\alpha$ et $\gamma = -\beta$, donc $\alpha = \beta$. La première donne $2\alpha = 1$, $\alpha = \beta = \dfrac{1}{2}$, $\gamma = -\dfrac{1}{2}$.
\[ \vec{e}_1 = \frac{1}{2}\vec{u} + \frac{1}{2}\vec{v} - \frac{1}{2}\vec{w}. \]
Vérification : $\dfrac{1}{2}(1,1,0) + \dfrac{1}{2}(1,0,1) - \dfrac{1}{2}(0,1,1) = (1, 0, 0)$. \checkmark
\end{enumerate}

\textbf{Partie B}
\begin{enumerate}
\item Les colonnes de $M$ sont les coordonnées de $g(\vec{u})$, $g(\vec{v})$, $g(\vec{w})$ dans $\mathcal{B}$ :
\[ g(\vec{u}) = \vec{u} + \vec{v} = 1\vec{u} + 1\vec{v} + 0\vec{w}, \quad g(\vec{v}) = 0\vec{u} + 1\vec{v} + 1\vec{w}, \quad g(\vec{w}) = 1\vec{u} + 0\vec{v} + 1\vec{w}. \]
\[ M = \begin{pmatrix} 1 & 0 & 1 \\ 1 & 1 & 0 \\ 0 & 1 & 1 \end{pmatrix}. \]
\item Soit $\vec{x}$ de coordonnées $(\alpha, \beta, \gamma)$ dans $\mathcal{B}$. Alors $g(\vec{x}) = \vec{0}$ si et seulement si $M \begin{pmatrix} \alpha \\ \beta \\ \gamma \end{pmatrix} = \vec{0}$ :
\[ \begin{cases} \alpha + \gamma = 0 \\ \alpha + \beta = 0 \\ \beta + \gamma = 0 \end{cases} \iff \begin{cases} \gamma = -\alpha \\ \beta = -\alpha \\ -\alpha - \alpha = 0 \end{cases} \iff \alpha = \beta = \gamma = 0. \]
Donc $\ker g = \{\vec{0}\}$ : $g$ est \textbf{injective}.
\item Par le théorème du rang : $\dim \operatorname{Im} g = 3 - \dim \ker g = 3 - 0 = 3$. Comme $\operatorname{Im} g$ est un sous-espace de $\mathbb{R}^3$ de dimension $3$, $\operatorname{Im} g = \mathbb{R}^3$ : $g$ est surjective. Étant injective et surjective, $g$ est \textbf{bijective}.
\item Calculons les images des vecteurs canoniques. D'après A.3 : $\vec{e}_1 = \dfrac{1}{2}\vec{u} + \dfrac{1}{2}\vec{v} - \dfrac{1}{2}\vec{w}$. Par linéarité :
\begin{align*}
g(\vec{e}_1) &= \frac{1}{2}g(\vec{u}) + \frac{1}{2}g(\vec{v}) - \frac{1}{2}g(\vec{w}) = \frac{1}{2}(\vec{u} + \vec{v}) + \frac{1}{2}(\vec{v} + \vec{w}) - \frac{1}{2}(\vec{w} + \vec{u}) = \vec{v} = (1, 0, 1).
\end{align*}
Par un raisonnement identique, exprimons $\vec{e}_2$ et $\vec{e}_3$ dans $\mathcal{B}$. Pour $\vec{e}_2 = (0,1,0)$ : $\alpha + \beta = 0$, $\alpha + \gamma = 1$, $\beta + \gamma = 0$ donnent $\alpha = \dfrac{1}{2}$, $\beta = -\dfrac{1}{2}$, $\gamma = \dfrac{1}{2}$, soit $\vec{e}_2 = \dfrac{1}{2}\vec{u} - \dfrac{1}{2}\vec{v} + \dfrac{1}{2}\vec{w}$. D'où :
\[ g(\vec{e}_2) = \frac{1}{2}(\vec{u} + \vec{v}) - \frac{1}{2}(\vec{v} + \vec{w}) + \frac{1}{2}(\vec{w} + \vec{u}) = \vec{u} = (1, 1, 0). \]
Pour $\vec{e}_3 = (0,0,1)$ : $\alpha + \beta = 0$, $\alpha + \gamma = 0$, $\beta + \gamma = 1$ donnent $\alpha = -\dfrac{1}{2}$, $\beta = \dfrac{1}{2}$, $\gamma = \dfrac{1}{2}$, soit $\vec{e}_3 = -\dfrac{1}{2}\vec{u} + \dfrac{1}{2}\vec{v} + \dfrac{1}{2}\vec{w}$. D'où :
\[ g(\vec{e}_3) = -\frac{1}{2}(\vec{u} + \vec{v}) + \frac{1}{2}(\vec{v} + \vec{w}) + \frac{1}{2}(\vec{w} + \vec{u}) = \vec{w} = (0, 1, 1). \]
La matrice de $g$ dans la base canonique a pour colonnes $g(\vec{e}_1)$, $g(\vec{e}_2)$, $g(\vec{e}_3)$ :
\[ M' = \begin{pmatrix} 1 & 1 & 0 \\ 0 & 1 & 1 \\ 1 & 0 & 1 \end{pmatrix}. \]
\emph{Vérification de cohérence :} $g(\vec{u}) = g(\vec{e}_1 + \vec{e}_2) = (1,0,1) + (1,1,0) = (2, 1, 1)$, et $\vec{u} + \vec{v} = (1,1,0) + (1,0,1) = (2,1,1)$. \checkmark
\end{enumerate}

\textbf{Points de vigilance :} dans la matrice $M$, ce sont les \emph{coordonnées dans $\mathcal{B}$} des images qui forment les colonnes, pas les coordonnées canoniques ; pour la question B.4, il faut justifier chaque passage par la linéarité de $g$ et réutiliser la méthode de A.3 ; conclure à la bijectivité exige de citer le théorème du rang \emph{et} l'inclusion $\operatorname{Im} g \subset \mathbb{R}^3$.
\end{corrige}

\newpage

\coursec{Fiche bilan}

\begin{popfigure}
\begin{tikzpicture}[mindmap, grow cyclic, every node/.style={concept, font=\small},
  root concept/.append style={concept color=popblue, font=\bfseries},
  level 1/.append style={level distance=3.4cm, sibling angle=60},
  level 2/.append style={level distance=2.2cm, font=\scriptsize}]
\node[root concept]{Espaces vectoriels et applications linéaires}
  child[concept color=poporange]{ node {Espace vectoriel $\mathbb{R}^3$}
    child { node {8 axiomes} }
    child { node {Sous-espace : $0_E\in F$, stable par comb. lin.} }
  }
  child[concept color=popgreen]{ node {Bases de $\mathbb{R}^3$}
    child { node {Famille libre} }
    child { node {Famille génératrice} }
    child { node {Base = libre + génératrice (3 vecteurs)} }
  }
  child[concept color=poppurple]{ node {Coordonnées}
    child { node {$\vec{u}=x\vec{e}_1+y\vec{e}_2+z\vec{e}_3$} }
    child { node {Unicité du triplet $(x;y;z)$} }
  }
  child[concept color=poppink]{ node {Applications linéaires}
    child { node {$f(u+v)=f(u)+f(v)$} }
    child { node {$f(\lambda u)=\lambda f(u)$} }
    child { node {$f(\lambda u+v)=\lambda f(u)+f(v)$} }
  }
  child[concept color=popgold]{ node {Matrice de $f$}
    child { node {Colonnes = images des vecteurs de base} }
    child { node {$Y=AX$} }
  }
  child[concept color=popdark]{ node {Noyau et image}
    child { node {$\ker f=\{u\,;\,f(u)=0\}$} }
    child { node {$\operatorname{Im}f=\{f(u)\}$} }
    child { node {Théorème du rang} }
  };
\end{tikzpicture}
\legende{Carte mentale de la leçon : les six blocs à maîtriser pour le Bac.}
\end{popfigure}

\begin{aretenir}[Synthèse de la leçon]
\textbf{Définitions clés.}
\begin{itemize}
  \item \cle{Espace vectoriel} sur $\mathbb{R}$ : ensemble $E$ muni d'une addition interne et d'une multiplication par un scalaire vérifiant les 8 axiomes (associativité, commutativité, élément neutre, symétrique, distributivités, associativité mixte, unité) ; $\mathbb{R}^3$ en est le modèle canonique de \cle{dimension 3}.
  \item \cle{Sous-espace vectoriel} $F$ de $E$ : $F \neq \varnothing$ (équivalent à $0_E\in F$) et pour tous $u,v\in F$, $\lambda\in\mathbb{R}$ : $\lambda u+v\in F$ (critère de stabilité par combinaison linéaire).
  \item \cle{Famille libre} $(\vec{u}_1,\dots,\vec{u}_k)$ : $\lambda_1\vec{u}_1+\dots+\lambda_k\vec{u}_k=\vec{0}\ \Rightarrow\ \lambda_1=\dots=\lambda_k=0$.
  \item \cle{Famille génératrice} de $E$ : tout vecteur de $E$ s'écrit comme combinaison linéaire des vecteurs de la famille.
  \item \cle{Base} de $E$ : famille à la fois libre et génératrice ; dans $\mathbb{R}^3$, toute base a exactement \textbf{3 vecteurs}.
  \item \cle{Application linéaire} $f:E\to F$ : $f(u+v)=f(u)+f(v)$ et $f(\lambda u)=\lambda f(u)$ pour tous $u,v\in E$, $\lambda\in\mathbb{R}$ ; condition équivalente unique : $f(\lambda u+v)=\lambda f(u)+f(v)$.
  \item \cle{Noyau} : $\ker f=\{u\in E\,;\,f(u)=0_F\}$, sous-espace vectoriel de $E$. \cle{Image} : $\operatorname{Im}f=\{f(u)\,;\,u\in E\}$, sous-espace vectoriel de $F$.
\end{itemize}

\textbf{Théorèmes et formules à connaître par cœur.}
\begin{center}
\begin{tabularx}{\linewidth}{|l|X|}
\hline
\rowcolor{popblueL} \textbf{Résultat} & \textbf{Énoncé} \\
\hline
Coordonnées & Dans une base $(\vec{e}_1,\vec{e}_2,\vec{e}_3)$, tout $\vec{u}$ s'écrit de manière \textbf{unique} $\vec{u}=x\vec{e}_1+y\vec{e}_2+z\vec{e}_3$. Le triplet $(x;y;z)$ est les coordonnées de $\vec{u}$. \\
\hline
Base de $\mathbb{R}^3$ & Une famille de 3 vecteurs de $\mathbb{R}^3$ est une base \emph{si et seulement si} elle est libre (il suffit de vérifier la liberté, ou que le déterminant est non nul). \\
\hline
Image d'une base & Une application linéaire est entièrement déterminée par les images des vecteurs d'une base de départ. \\
\hline
Matrice de $f$ & La $j$-ième colonne de la matrice $A$ contient les coordonnées de $f(\vec{e}_j)$ dans la base d'arrivée ; si $X$ sont les coordonnées de $\vec{u}$ et $Y$ celles de $f(\vec{u})$, alors $Y=AX$. \\
\hline
Injectivité & $f$ injective $\iff \ker f=\{0_E\}$. \\
\hline
Théorème du rang & $\dim E=\dim\ker f+\dim\operatorname{Im}f$ (en dimension 3 : $\dim\ker f+\operatorname{rg} f=3$). \\
\hline
\end{tabularx}
\end{center}

\textbf{Méthodes express.}
\begin{enumerate}
  \item \emph{Montrer qu'une famille de 3 vecteurs est une base de $\mathbb{R}^3$} :
  \begin{itemize}
    \item Méthode système : poser $\lambda_1\vec{e}_1+\lambda_2\vec{e}_2+\lambda_3\vec{e}_3=\vec{0}$, résoudre le système $3\times 3$ ; si l'unique solution est $(0;0;0)$, la famille est libre, donc c'est une base.
    \item Méthode déterminant : former la matrice $P$ dont les colonnes sont les vecteurs ; si $\det P \neq 0$, c'est une base.
  \end{itemize}
  \item \emph{Trouver les coordonnées de $\vec{u}$ dans une base $(\vec{e}_1,\vec{e}_2,\vec{e}_3)$} :
  \begin{itemize}
    \item Résoudre $\vec{u}=x\vec{e}_1+y\vec{e}_2+z\vec{e}_3$ (système de 3 équations).
    \item Si $P$ est la matrice de passage (colonnes = $\vec{e}_i$), alors $\begin{pmatrix}x\\y\\z\end{pmatrix}=P^{-1}\vec{u}$.
  \end{itemize}
  \item \emph{Prouver qu'une application est linéaire} : vérifier $f(\lambda u+v)=\lambda f(u)+f(v)$ pour $u,v$ quelconques ; contrôler d'abord $f(0)=0$ (si $f(0)\neq 0$, ce n'est pas linéaire).
  \item \emph{Écrire la matrice de $f$ dans les bases $(\vec{e}_i)$ et $(\vec{f}_j)$} : calculer $f(\vec{e}_1)$, $f(\vec{e}_2)$, $f(\vec{e}_3)$ et placer leurs coordonnées \textbf{en colonnes} dans la base d'arrivée $(\vec{f}_j)$.
  \item \emph{Déterminer $\ker f$} : résoudre $f(x;y;z)=(0;0;0)$ ; exprimer les solutions en fonction de paramètres libres ; une base de $\ker f$ est formée par les vecteurs directeurs.
  \item \emph{Déterminer $\operatorname{Im}f$} : $\operatorname{Im}f=\operatorname{Vect}\big(f(\vec{e}_1),f(\vec{e}_2),f(\vec{e}_3)\big)$ ; extraire une famille libre (par élimination de Gauss sur la matrice de $f$).
  \item \emph{Utiliser le théorème du rang} : dès qu'une des deux dimensions ($\ker f$ ou $\operatorname{Im}f$) est connue, l'autre s'obtient par $\dim\ker f+\dim\operatorname{Im}f=3$.
\end{enumerate}
\end{aretenir}

\begin{aretenir}[Checklist avant le Bac]
Avant l'épreuve, vérifie que tu sais :
\begin{itemize}
  \item[$\square$] Citer les deux opérations et l'idée des 8 axiomes d'un espace vectoriel.
  \item[$\square$] Prouver qu'un ensemble est un sous-espace vectoriel avec le test $\lambda u+v\in F$ (sans oublier $F\neq\varnothing$ ou $0_E\in F$).
  \item[$\square$] Démontrer qu'une famille de 3 vecteurs de $\mathbb{R}^3$ est libre en résolvant un système homogène (ou en calculant un déterminant).
  \item[$\square$] Conclure qu'une famille libre de 3 vecteurs est automatiquement une base de $\mathbb{R}^3$.
  \item[$\square$] Calculer les coordonnées d'un vecteur dans une base donnée par résolution d'un système (ou inversion de matrice).
  \item[$\square$] Vérifier la linéarité d'une application avec la condition unique $f(\lambda u+v)=\lambda f(u)+f(v)$.
  \item[$\square$] Construire la matrice de $f$ en plaçant les images des vecteurs de base \textbf{en colonnes} dans la base d'arrivée.
  \item[$\square$] Déterminer le noyau de $f$ en résolvant $f(u)=0$ et en donnant une base de $\ker f$.
  \item[$\square$] Déterminer l'image de $f$ comme sous-espace engendré par les images des vecteurs de base, puis en extraire une base.
  \item[$\square$] Énoncer et appliquer le théorème du rang : $\dim\ker f+\dim\operatorname{Im}f=3$.
  \item[$\square$] Utiliser $\ker f=\{0\}\iff f$ injective, et en déduire la bijectivité quand $\operatorname{rg}f=3$.
  \item[$\square$] Rédiger proprement : annoncer la méthode, justifier chaque étape, conclure par une phrase.
\end{itemize}
\end{aretenir}

\findecours

\end{document}
